% Options for packages loaded elsewhere
\PassOptionsToPackage{unicode}{hyperref}
\PassOptionsToPackage{hyphens}{url}
\PassOptionsToPackage{dvipsnames,svgnames,x11names}{xcolor}
%
\documentclass[
  11pt,
]{article}
\usepackage{amsmath,amssymb}
\usepackage{setspace}
\usepackage{iftex}
\ifPDFTeX
  \usepackage[T1]{fontenc}
  \usepackage[utf8]{inputenc}
  \usepackage{textcomp} % provide euro and other symbols
\else % if luatex or xetex
  \usepackage{unicode-math} % this also loads fontspec
  \defaultfontfeatures{Scale=MatchLowercase}
  \defaultfontfeatures[\rmfamily]{Ligatures=TeX,Scale=1}
\fi
\usepackage{lmodern}
\ifPDFTeX\else
  % xetex/luatex font selection
  \setmainfont[BoldFont=lmroman10-bold.otf,ItalicFont=lmroman10-italic.otf,BoldItalicFont=lmroman10-bolditalic.otf]{lmroman10-regular.otf}
  \setsansfont[BoldFont=lmsans10-bold.otf,ItalicFont=lmsans10-oblique.otf,BoldItalicFont=lmsans10-boldoblique.otf]{lmsans10-regular.otf}
  \setmonofont[ItalicFont=lmmono10-italic.otf,HyphenChar=None]{lmmono10-regular.otf}
  \setmathfont[]{latinmodern-math.otf}
\fi
% Use upquote if available, for straight quotes in verbatim environments
\IfFileExists{upquote.sty}{\usepackage{upquote}}{}
\IfFileExists{microtype.sty}{% use microtype if available
  \usepackage[]{microtype}
  \UseMicrotypeSet[protrusion]{basicmath} % disable protrusion for tt fonts
}{}
\makeatletter
\@ifundefined{KOMAClassName}{% if non-KOMA class
  \IfFileExists{parskip.sty}{%
    \usepackage{parskip}
  }{% else
    \setlength{\parindent}{0pt}
    \setlength{\parskip}{6pt plus 2pt minus 1pt}}
}{% if KOMA class
  \KOMAoptions{parskip=half}}
\makeatother
\usepackage{xcolor}
\usepackage[margin=1in]{geometry}
\usepackage{longtable,booktabs,array}
\usepackage{calc} % for calculating minipage widths
% Correct order of tables after \paragraph or \subparagraph
\usepackage{etoolbox}
\makeatletter
\patchcmd\longtable{\par}{\if@noskipsec\mbox{}\fi\par}{}{}
\makeatother
% Allow footnotes in longtable head/foot
\IfFileExists{footnotehyper.sty}{\usepackage{footnotehyper}}{\usepackage{footnote}}
\makesavenoteenv{longtable}
\usepackage{graphicx}
\makeatletter
\def\maxwidth{\ifdim\Gin@nat@width>\linewidth\linewidth\else\Gin@nat@width\fi}
\def\maxheight{\ifdim\Gin@nat@height>\textheight\textheight\else\Gin@nat@height\fi}
\makeatother
% Scale images if necessary, so that they will not overflow the page
% margins by default, and it is still possible to overwrite the defaults
% using explicit options in \includegraphics[width, height, ...]{}
\setkeys{Gin}{width=\maxwidth,height=\maxheight,keepaspectratio}
% Set default figure placement to htbp
\makeatletter
\def\fps@figure{htbp}
\makeatother
\setlength{\emergencystretch}{3em} % prevent overfull lines
\providecommand{\tightlist}{%
  \setlength{\itemsep}{0pt}\setlength{\parskip}{0pt}}
\setcounter{secnumdepth}{-\maxdimen} % remove section numbering
\newlength{\cslhangindent}
\setlength{\cslhangindent}{1.5em}
\newlength{\csllabelwidth}
\setlength{\csllabelwidth}{3em}
\newlength{\cslentryspacingunit} % times entry-spacing
\setlength{\cslentryspacingunit}{\parskip}
\newenvironment{CSLReferences}[2] % #1 hanging-ident, #2 entry spacing
 {% don't indent paragraphs
  \setlength{\parindent}{0pt}
  % turn on hanging indent if param 1 is 1
  \ifodd #1
  \let\oldpar\par
  \def\par{\hangindent=\cslhangindent\oldpar}
  \fi
  % set entry spacing
  \setlength{\parskip}{#2\cslentryspacingunit}
 }%
 {}
\usepackage{calc}
\newcommand{\CSLBlock}[1]{#1\hfill\break}
\newcommand{\CSLLeftMargin}[1]{\parbox[t]{\csllabelwidth}{#1}}
\newcommand{\CSLRightInline}[1]{\parbox[t]{\linewidth - \csllabelwidth}{#1}\break}
\newcommand{\CSLIndent}[1]{\hspace{\cslhangindent}#1}
\usepackage{longtable,booktabs,array}
\usepackage{calc}
\usepackage{needspace}
\usepackage{placeins}
\usepackage{xurl}
\usepackage{seqsplit}
\emergencystretch=3em
\brokenpenalty=10000
\widowpenalty=10000
\clubpenalty=10000
% Pages are ragged-bottom, so this stretch is never visible; it lets a page end between tight-list items.
\renewcommand{\tightlist}{\setlength{\itemsep}{0pt plus 1pt}\setlength{\parskip}{0pt}}
\setlength{\tabcolsep}{4pt}
\renewcommand{\arraystretch}{1.12}
\setlength{\parindent}{0pt}
\setlength{\parskip}{5pt plus 1pt minus 1pt}
\renewcommand{\topfraction}{0.9}
\renewcommand{\textfraction}{0.08}
\renewcommand{\floatpagefraction}{0.8}
\setlength{\textfloatsep}{14pt plus 4pt minus 2pt}
\AtBeginDocument{\hypersetup{breaklinks=true}}
\ifLuaTeX
  \usepackage{selnolig}  % disable illegal ligatures
\fi
\IfFileExists{bookmark.sty}{\usepackage{bookmark}}{\usepackage{hyperref}}
\IfFileExists{xurl.sty}{\usepackage{xurl}}{} % add URL line breaks if available
\urlstyle{same}
\hypersetup{
  pdftitle={Supplement: signals, interfaces and worked-example data},
  colorlinks=true,
  linkcolor={black},
  filecolor={Maroon},
  citecolor={Blue},
  urlcolor={blue},
  pdfcreator={LaTeX via pandoc}}

\author{}
\date{}

\begin{document}

\setstretch{1}
\FloatBarrier

\begingroup\raggedright\hyphenpenalty=10000\exhyphenpenalty=10000

\hypertarget{supplement-signals-interfaces-and-worked-example-data}{%
\section{Supplement: signals, interfaces and worked-example
data}\label{supplement-signals-interfaces-and-worked-example-data}}

\endgroup

\suppressfloats[t]

This supplement accompanies the \href{MANUSCRIPT.md}{complete
manuscript}. S1 and S2 explain the sources and selection. S3 to S5
specify the signals, records and interfaces. S6 supplies the
hypothetical investigations and their arithmetic; S7 and S8 give the
proposed evaluations. S9 extends the comparison with prior work. The 19
CSV files use \mbox{UTF-8} encoding and stable mechanism, source and
signal identifiers. Citation labels match the package's
\href{REFERENCES.bib}{BibLaTeX references}.

\FloatBarrier

\hypertarget{s1.-sources-and-comparisons}{%
\subsection{S1. Sources and
comparisons}\label{s1.-sources-and-comparisons}}

\suppressfloats[t]

\href{SUPPLEMENT/S1_CITED_SOURCE_NOTES.csv}{Source notes} identify the
version, supporting passages, publication status and supported claim for
each of the 127 cited works. The
\href{SUPPLEMENT/S1_REQUIREMENTS_MATRIX.csv}{requirements matrix}
compares 50 sources against C1 to C3. It codes two parts of C1 and C2
and three of C3: an outcome linked to its decision, and late,
unobservable and false reports kept apart (C1); query coverage and
learned time (C2); typed separation of observations, inferences and
memory, a human principal before reuse, and links from records an agent
output informed back to that output (C3). E means explicit support, P
partial or adjacent support, and a dot means the cited material does not
establish the requirement. The \texttt{payment\_setting} column records
yes for a payment or payment-fraud setting, P for a setting partly in
payments or adjacent to payment fraud, and no otherwise. It marks two
cyber threat intelligence sources, STIX 2.1 and Hettema's
belief-revision paper, as cyber; other security and threat-intelligence
sources are coded no. The source notes provide passage locators, linked
to the matrix by citation key. Their inspection dates are in UTC.

\href{SUPPLEMENT/S1_STAGE_COVERAGE.csv}{Stage coverage} supplies the
source-level codes behind Table S1.1. Each family receives the strongest
code among its sources, so a family row may combine capabilities from
several systems. Source notes record bibliographic and access issues
separately from claims about a method.

\begin{table}[tp]
\textbf{Table S1.1. Stages described in the cited research.} E means
explicitly described; P means partial coverage or a different scope; a
dot means not established by the cited material. Each family receives
the strongest source-level code in S1, so one row may combine several
systems. L1: signal sources; L2: ingestion and stores; L3: entities,
correlation and knowledge; L4: investigation; L5: outputs; L6: feedback;
L7: governance; L8: learning and correction.\par\smallskip
{\footnotesize
\begin{tabular}{@{}
  >{\raggedright\arraybackslash}p{(\columnwidth - 18\tabcolsep) * \real{0.4400}}
  >{\raggedright\arraybackslash}p{(\columnwidth - 18\tabcolsep) * \real{0.0550}}
  >{\raggedright\arraybackslash}p{(\columnwidth - 18\tabcolsep) * \real{0.0550}}
  >{\raggedright\arraybackslash}p{(\columnwidth - 18\tabcolsep) * \real{0.0550}}
  >{\raggedright\arraybackslash}p{(\columnwidth - 18\tabcolsep) * \real{0.0550}}
  >{\raggedright\arraybackslash}p{(\columnwidth - 18\tabcolsep) * \real{0.0550}}
  >{\raggedright\arraybackslash}p{(\columnwidth - 18\tabcolsep) * \real{0.0550}}
  >{\raggedright\arraybackslash}p{(\columnwidth - 18\tabcolsep) * \real{0.0550}}
  >{\raggedright\arraybackslash}p{(\columnwidth - 18\tabcolsep) * \real{0.0550}}
  >{\raggedright\arraybackslash}p{(\columnwidth - 18\tabcolsep) * \real{0.1200}}@{}}
\toprule\noalign{}
\begin{minipage}[b]{\linewidth}\raggedright
Research family and inspected examples
\end{minipage} & \begin{minipage}[b]{\linewidth}\raggedright
L1
\end{minipage} & \begin{minipage}[b]{\linewidth}\raggedright
L2
\end{minipage} & \begin{minipage}[b]{\linewidth}\raggedright
L3
\end{minipage} & \begin{minipage}[b]{\linewidth}\raggedright
L4
\end{minipage} & \begin{minipage}[b]{\linewidth}\raggedright
L5
\end{minipage} & \begin{minipage}[b]{\linewidth}\raggedright
L6
\end{minipage} & \begin{minipage}[b]{\linewidth}\raggedright
L7
\end{minipage} & \begin{minipage}[b]{\linewidth}\raggedright
L8
\end{minipage} & \begin{minipage}[b]{\linewidth}\raggedright
Payment domain
\end{minipage} \\
\midrule\noalign{}

Industrial payment and risk systems: TitAnt, GEM & P & E & P & . & P & P
& . & P & yes \\
Illustrative platform's (Stripe's) engineering publications, 2019 to
2025 & P & E & P & P & P & E & . & P & yes \\
Threat-intelligence and security-operations architectures: reference
models, intelligence lifecycle and agent benchmarks & P & P & E & E & E
& P & E & E & no \\
Preprints on agents making fraud decisions: SR-Fraud and Sharma's
auditable agent & P & P & P & E & P & E & P & E & yes \\

\bottomrule\noalign{}
\end{tabular}
}
\end{table}

Study populations and denominators constrain the comparisons. Sharma's
balanced synthetic sample contains 60 cases drawn from transactions its
classifier scored between 0.2 and 0.8: thresholding the classifier at
0.5 decides 43 correctly and the agent 39. Six of the agent's eight
overrides replace a correct decision with an error. Fatima's
verification rates concern genuine participants under the study's
completion and exclusion rules; they do not measure fraud detection.
Doerfler's 47.8 percent same-session success rate concerns challenged
sign-ins. Its 97.3 percent eventual-access rate counts users, not
sign-ins: it is the share of users with access within one week.
Doerfler's Section 5.3 and Table 7 caption take users who failed a
challenge as the denominator; Doerfler's introduction takes users who
saw a challenge. Its abstract treats the complement of the first rate as
legitimate users failing to sign in, but its results note that hijacking
activity depresses per-session success. Subtracting the first rate from
100 percent therefore does not give a clean legitimate-user failure
rate. Neither Fatima's nor Doerfler's aggregate rate establishes whether
an individual event was decided correctly.

FTC complaints describe allegations, and the stipulated orders do not
establish those allegations on their merits. Vendor publications support
attributed descriptions of interfaces or observations, not independent
evidence of industry-wide effectiveness.

\FloatBarrier

\hypertarget{s2.-product-scope-and-mechanism-selection}{%
\subsection{S2. Product scope and mechanism
selection}\label{s2.-product-scope-and-mechanism-selection}}

\suppressfloats[t]

\href{SUPPLEMENT/S2_SURFACE_BASELINE.csv}{Product scope} contains 88
public product and feature entries, read on 27 September 2026: 53
catalog items, twenty items from the Q3 2026 public roadmap (seventeen
appear in a cached copy of that roadmap; the other three are confirmed
as shipped in Stripe's changelog), ten separately documented
capabilities and five careers-page footer items. These entries define
the scope for candidate selection, not a complete or current product
catalog. Each row lists candidate mechanism IDs, not validated coverage
or retained signal mappings; a dot means none is listed. The entries
list 85 of the 99 candidates; the other 14 appear only in the allocation
table. The allocation table gives each candidate's final disposition,
and S3 supplies the retained mappings. Six retained mechanisms are not
listed against any product entry: \mbox{ID-14} (account farming and
aged-account preparation for later abuse), \mbox{PAY-04} (card testing
through attacker-controlled merchant or connected accounts),
\mbox{MP-02} (payment-method cashing through a self-controlled merchant
or connected account), \mbox{MP-04} (merchant misrepresentation and
transaction laundering behind a front), \mbox{MP-07} (merchant-side
dispute evasion) and \mbox{BOT-04} (CAPTCHA and challenge defeat). S5
identifies the exact API schemas used for interface comparisons.

\textbf{Products and business-reported records (manuscript Section
4.1).} The platform's products may span onboarding, identity
verification, payments, subscriptions, payouts, treasury and risk
controls. Businesses may also report their customers' registrations,
logins and payments. The specification identifies what investigators
need from these records without assuming a provider already collects or
links them.

A retained mechanism targets a documented capability, serves a plausible
adversary goal and has supporting public evidence. Sources include
research, incident reports, government and intergovernmental guidance,
and platform documentation. Four public job postings for abuse, risk and
machine-learning roles helped prioritize topics but do not establish
incidents; the manuscript cites two (Section 3), and the allocation keys
record the other two. The stated priority basis names a posting topic
for 24 candidates. Twenty-three of them are among the 26 candidates
allocated Tier 1; all 23 are mapped and carry 65 of the 163 links. The
24th, \mbox{EC-06}, is unmapped. Every mapped mechanism has support
beyond postings in the mechanism-to-source links.

\href{SUPPLEMENT/S2_MECHANISM_ALLOCATION.csv}{Mechanism allocation}
records all 99 candidates and their final dispositions. The signal map
retains 61 mechanisms: 52 of the 53 dedicated candidates, all eight
targeted candidates and the one adjacent-logged candidate, \mbox{ID-18},
insider access abuse by platform or partner staff. Of the other 38, nine
were judged covered by existing records, four were first judged covered
and are now recorded as record gaps, four were screened without
identified supporting evidence and logged as low-priority gaps, twenty
low-priority candidates were not reviewed, and one reviewed candidate,
\mbox{PAY-12}, has no signal mapping. \mbox{PAY-12} concerns
triangulation fraud: filling real customers' orders through another
seller using stolen payment credentials. Its review record notes
framework evidence only. The four record gaps are fake reviews
(\mbox{MP-08}), e-skimming on merchant checkout pages (\mbox{API-09}),
supply-chain compromise of commerce plugins and packages (\mbox{API-10})
and business stablecoin issuance abuse (\mbox{MP-11}). A second rater
named a record each would need, and the inventory holds no such record
(S3). \href{SUPPLEMENT/S2_ALLOCATION_KEYS.csv}{Allocation keys} decode
the source abbreviations. The \texttt{Class} code, also the prefix
before each key's colon, gives the source type: a for law-enforcement
announcements and vendor or incident reports, including EchoLeak
(a:ECHO), which is cited by its arXiv identifier; b for research papers;
b* for other arXiv preprints; c for government and intergovernmental
publications and the MITRE fraud framework; d for the illustrative
platform's documentation, roadmap, blog and public repositories; and e
for the illustrative platform's job postings (Stripe careers). All four
b keys name a peer-reviewed venue. In the \texttt{Verification} column,
V marks a key recorded as verified, and S marks a key that rests on a
summarizing retrieval whose exact passage was not retained, as the
\texttt{d:S} row states. The allocation file repeats S as a suffix on
its evidence keys, as in \mbox{\texttt{a:SSN\ (S)}}. Priority labels
explain selection emphasis; they do not measure evidence quality.

\hypertarget{reading-the-selection-records}{%
\paragraph{Reading the selection
records}\label{reading-the-selection-records}}

The allocation file's \texttt{disposition} field states the final
selection decision. Signal coverage is recorded separately in S3; a
reviewed candidate need not have a signal mapping.

\begingroup\small

\begin{longtable}[]{@{}
  >{\raggedright\arraybackslash}p{(\columnwidth - 2\tabcolsep) * \real{0.2000}}
  >{\raggedright\arraybackslash}p{(\columnwidth - 2\tabcolsep) * \real{0.8000}}@{}}
\toprule\noalign{}
\begin{minipage}[b]{\linewidth}\raggedright
Disposition
\end{minipage} & \begin{minipage}[b]{\linewidth}\raggedright
Meaning
\end{minipage} \\
\midrule\noalign{}
\endhead
\bottomrule\noalign{}
\endlastfoot
dedicated & Reviewed as a distinct mechanism \\
targeted & Received a focused review of its implications for the
design \\
covered & Judged covered by existing records rather than retained as a
separate mapped mechanism \\
record gap & Reviewed and first judged covered; a second-rater check
found that no record in the inventory covers it \\
screened gap & Screened without identified supporting evidence and
logged as a low-priority gap \\
unreviewed & Not reviewed; all twenty such candidates have low
priority \\
adjacent logged & Recorded as adjacent to the core mechanisms; the one
such candidate is retained in the signal map \\
\end{longtable}

\endgroup

\mbox{\texttt{Priority\ (basis)}} and \mbox{\texttt{Impact\ (basis)}}
retain the selection judgments and their stated reasons. In a priority
basis, \texttt{r:} introduces a topic named in a job posting, with the
posting's key in parentheses where given, and \texttt{p:} introduces
platform-documentation keys. A posting key such as \texttt{(EM)} is the
allocation key \texttt{e:EM} without its class prefix;
platform-documentation keys keep their \texttt{d:} prefix. Impact bases
use no prefixes. \texttt{Scope} distinguishes core and adjacent
mechanisms. \texttt{Parent} records a stated relationship to another
mechanism; a dot means none is listed. \mbox{\texttt{Evidence\ keys}}
resolve to source descriptions in the allocation-keys file. They are
identifiers, not evidence grades. \texttt{Maturity} is the estimate,
before review, of how developed the published research was;
\texttt{after\_review\_maturity} pairs it with the later classification
where both were supplied. \texttt{Evidence} is the confidence label used
to allocate effort (verified, provisional or unresolved), not a
certification of each source's claims. \texttt{Depth} is the initial
allocation: Tier 1 for the deeper review, Tier 2 otherwise, or a logged
gap when evidence was unresolved. \texttt{recorded\_status} states each
review's outcome. The product-scope file links documented capabilities
to candidates; it does not supply a product-availability assessment.

Searches covered nine abuse paths, coded U1 to U9: API-key takeover
(U1), identity provenance at onboarding (U2), document integrity,
including identity and business documents (U3), account compromise (U4),
payout redirection (U5), card testing (U6), abusive merchants (U7), mule
and pass-through funds (U8), and attacks on agents or operator
assistants (U9). Only the entity-relationship notes cite these codes.
Searches followed relevant publications, public guidance and references.
This focused selection does not establish exhaustive coverage of abuse
or related systems.

\hypertarget{s2.1-claim-types}{%
\subsubsection{S2.1 Claim types}\label{s2.1-claim-types}}

Research can propose an untested design, while documentation can
directly establish an interface. Neither publication status nor the
number of sources validates the integrated specification. No claim in
either document rests on author experience.

\begin{table}[htp]
\textbf{Table S2.1. How to read the manuscript's claims.} Claim type
describes what a statement establishes. Source type and publication
status describe where its support comes from.\par\smallskip
{\small
\begin{tabular}{@{}
  >{\raggedright\arraybackslash}p{(\columnwidth - 2\tabcolsep) * \real{0.2500}}
  >{\raggedright\arraybackslash}p{(\columnwidth - 2\tabcolsep) * \real{0.7500}}@{}}
\toprule\noalign{}
\begin{minipage}[b]{\linewidth}\raggedright
Claim type
\end{minipage} & \begin{minipage}[b]{\linewidth}\raggedright
What it means
\end{minipage} \\
\midrule\noalign{}

Fact & A description supported by evidence, including a public
specification or an artifact this study produced; source and version are
given where they matter \\
Peer-reviewed prior research & A finding reported in a peer-reviewed
publication whose publication status was confirmed, applicable only to
the setting and conditions it studied \\
Inference & A conclusion motivated by stated premises; the sources do
not establish the conclusion directly \\
Design requirement & A property chosen for the proposed platform on the
stated evidence and reasoning; it is not claimed as necessary for every
platform or as observed in a deployment \\
Hypothesis & A proposition to test for the proposed platform, including
every detection candidate in the manuscript \\
Proposed implementation & One possible way to satisfy a requirement \\
Empirically validated result & A measurement, valid for the proposition
and setting actually tested; results for individual components do not
validate this integrated design \\

\bottomrule\noalign{}
\end{tabular}
}
\end{table}

\hypertarget{s2.2-external-coverage-check}{%
\subsubsection{S2.2 External coverage
check}\label{s2.2-external-coverage-check}}

An external check compares the candidates with MITRE's Fight Fraud
Framework (F3), version 1.2
{[}\protect\hyperlink{ref-MITRE2026F3}{1}{]}. Mechanism selection used
version 1.1 (allocation key c:F3); version 1.2 adds 16 techniques and
removes one. The \href{SUPPLEMENT/S2_F3_CROSSWALK.csv}{F3 crosswalk}
maps each of its 138 techniques and sub-techniques to the candidate
mechanisms that cover it, with a reason for each scope decision. Of the
105 within the product scope, a retained mechanism covers 35 fully and
49 in part. Of the other 21, fourteen are covered only by candidates
outside the signal map. Two of the fourteen rest only on candidates
judged covered by existing records, among them device fingerprint
spoofing. Three rest only on candidates recorded as record gaps: the
adversary-in-the-browser technique and two of its sub-techniques, which
the candidates describe but no record in the inventory covers. Six rest
only on unreviewed candidates, among them 3-D Secure bypass, and three
rest on covered and unreviewed candidates together. None rests on a
screened gap. Seven are covered by no candidate: DLL injection in the
browser, compromise of point-of-sale systems, coercion, keylogging,
abuse of accessibility features, theft from password managers and SEO
poisoning. Scope and coverage are one reviewer's judgments, recorded for
each row.

\FloatBarrier

\hypertarget{s3.-signals-mappings-and-lineage}{%
\subsection{S3. Signals, mappings and
lineage}\label{s3.-signals-mappings-and-lineage}}

\suppressfloats[t]

The \href{SUPPLEMENT/S3_SIGNALS.csv}{signal inventory} defines all 71
signals, with their families, classes, feature inputs, mechanism and
source IDs, public counterparts, entities, uses and limits. Its
\texttt{internal\_representation\_status} value, DR for every signal,
marks a design requirement rather than observed platform
instrumentation.
\href{SUPPLEMENT/S3_MECHANISM_SIGNAL_PAIRS.csv}{Mechanism-to-signal
mappings} contain 163 links: 154 design mappings motivated by evidence
and nine marked as inferential. These links explain why a record is
proposed; they do not establish a detector's effectiveness or make its
fields uniquely necessary.

The outcome-channel requirements use \mbox{SIG-70} and \mbox{SIG-71} to
record each payment-level or account-level control decision's effect on
each outcome channel, \mbox{SIG-27} to separate authentication results
from reported liability status, and \mbox{SIG-58} for dispute reasons
and payout returns. \mbox{SIG-70} also keeps manual-review dispositions.
Each record retains its source, version and uncertainty.

The inventory contains 26 source observations (SO), two normalized
observations (NO), 24 derived features (DF), two model or risk outputs
(MO) and ten control outcomes (CO). Seven signals provide external
knowledge (XK). The inventory therefore covers five of the seven record
classes in manuscript Section 5.1, plus external knowledge. Analytical
inferences and finished intelligence are later outputs, not input
signals.

\begin{figure}[tp]

\textbf{Figure S3.1. What each record class can establish.} Counts refer
to the 71 signals. Analytical inferences and finished intelligence are
outputs, so neither appears among those input signals. Test, replay and
emulation records are marked when created to distinguish them from
production activity.

\includegraphics[width=0.86\textwidth,height=\textheight]{FIGURES/FIG_S3_1_SIGNAL_CLASSES.pdf}

\end{figure}

\begin{table}[tp]
\textbf{Table S3.1. The 71 signals by family.} SO: source observation;
NO: normalized observation; DF: derived feature; MO: model or risk
output; CO: control outcome; XK: external knowledge. The final column
counts partial public counterparts, cases where that comparison does not
apply, and signals with no counterpart identified in the pinned schemas
(manuscript Section 5.4). Token allowances permit another party to spend
from a crypto wallet.\par\smallskip
{\small
\begin{tabular}{@{}
  >{\raggedright\arraybackslash}p{(\columnwidth - 6\tabcolsep) * \real{0.4700}}
  >{\raggedright\arraybackslash}p{(\columnwidth - 6\tabcolsep) * \real{0.1100}}
  >{\raggedright\arraybackslash}p{(\columnwidth - 6\tabcolsep) * \real{0.2300}}
  >{\raggedright\arraybackslash}p{(\columnwidth - 6\tabcolsep) * \real{0.1900}}@{}}
\toprule\noalign{}
\begin{minipage}[b]{\linewidth}\raggedright
Family
\end{minipage} & \begin{minipage}[b]{\linewidth}\raggedright
Signals
\end{minipage} & \begin{minipage}[b]{\linewidth}\raggedright
Classes
\end{minipage} & \begin{minipage}[b]{\linewidth}\raggedright
Partial / n/a / none identified
\end{minipage} \\
\midrule\noalign{}

Identity and onboarding & 8 & SO 3, DF 4, CO 1 & 3 / 4 / 1 \\
Authentication, session and operator administration & 10 & SO 5, DF 4,
CO 1 & 6 / 4 / 0 \\
Credentials, API keys and request context & 5 & SO 3, DF 1, XK 1 & 2 / 2
/ 1 \\
Funds movement and payout destinations & 11 & SO 4, DF 5, CO 2 & 6 / 5 /
0 \\
Payment attempts, authorization and payment controls & 7 & SO 2, DF 4,
CO 1 & 3 / 4 / 0 \\
End-customer lifecycle and entitlement & 7 & SO 2, DF 4, CO 1 & 3 / 4 /
0 \\
Merchant storefront and business claims & 2 & DF 1, XK 1 & 0 / 2 / 0 \\
Platform outputs and labels & 4 & SO 1, MO 2, CO 1 & 4 / 0 / 0 \\
Agents, tools and approvals & 8 & SO 5, CO 3 & 4 / 0 / 4 \\
Wallet signing and token allowances & 2 & SO 1, DF 1 & 0 / 1 / 1 \\
External threat knowledge & 5 & XK 5 & 0 / 5 / 0 \\
Provenance classes & 2 & NO 2 & 1 / 1 / 0 \\

\bottomrule\noalign{}
\end{tabular}
}
\end{table}

\mbox{SIG-44}, \mbox{SIG-45} and \mbox{SIG-46} apply across mechanisms.
They record technique-catalog mappings, minimum observation provenance
and the provenance of test records, respectively. The other 68 signals
map to specific mechanisms. All 61 mapped mechanisms have at least one
mapping motivated by evidence; S2 explains \mbox{PAY-12}'s exclusion.

\textbf{Second-rater check of the links (10 November 2026).} The links
were made without a written linking rule. For this check a rule was
written afterward (\href{tools/second_coder/PACKET.md}{packet}): link a
mechanism to a signal when an investigator of that mechanism would need
the record, and a listed source describes behaviour the record would
capture. A rater that took no part in the mapping, OpenAI's \mbox{GPT-6}
in a session without memory, project files or web access, applied the
rule to a seeded random sample of 20 of the 61 mapped mechanisms and 5
of the 13 then judged covered. It saw the 68 mechanism-specific signals
and the recorded scope of each mechanism's sources, not the sources and
not the links. Claude models had proposed the links; ChatGPT models
assisted with other parts of the paper. Over the 1,360 mechanism-signal
pairs for the 20 mapped mechanisms, the rater made 49 of the paper's 52
links and 70 that the paper lacks (Cohen's kappa 0.55); two mechanisms
matched exactly. Of the three links it did not make, the paper marks two
as inference. Thirty-two of the 70 added links attach six
general-purpose signals, such as request context and funds movement, to
further mechanisms. The links are therefore a sourced subset: each gives
a reason for a record, and other mechanisms may need the same record.
Counts of links describe this subset, not a complete map. For four of
the five covered mechanisms the rater named no missing record; for fake
reviews (\mbox{MP-08}) it named review content, reviewer identities and
posting histories. A second packet on 11 November gave the same rater
the other eight covered mechanisms; it named no missing record for five
and a missing record for three (\mbox{API-09}, \mbox{API-10} and
\mbox{MP-11}). The inventory was searched for each named record and none
was found, so these four are now recorded as record gaps (S2) and the
other nine remain covered. The check tests whether the rule reproduces
the links from recorded source scopes. It is not a human coder and does
not test the sources. The
\href{tools/second_coder/real_run/RUN_RECORD.md}{run record}, answers
and score are in the package.

Each input ID for the 24 derived features names another signal in the
inventory. No feature depends, directly or indirectly, on itself. For
example, \mbox{SIG-24} uses \mbox{SIG-27} to enumerate its eligible
transaction cohort, including completed payments with no report. The
separate chargeback-or-alert-resolution measure also uses \mbox{SIG-25}.
Keeping these inputs explicit makes the denominator and the distinction
between measures inspectable.

The \href{SUPPLEMENT/S3_SOURCE_DIRECTORY.csv}{source directory}
identifies 61 design sources. Of these, 27 are not among the 127 cited
works and have no S1 source note. Each source was read in full in
October 2026. For one, a vendor blog post (L112), the cited passages
were read verbatim and the other paragraphs only as a close paraphrase.
One FTC complaint (L127) was first read in part on 21 October; a
complete copy was read on 29 October. The directory records each
source's type, the depth and version of the read, with qualifiers such
as the version read or pages not read, the access date and the scope.
Where only a preprint or an author copy was read, the row says so.
\href{SUPPLEMENT/S3_MAPPING_SOURCE_JOINS.csv}{Mechanism/source joins}
supply 126 links, each with its source identifier, recorded depth and
scope limit. These records preserve the premises behind the mappings,
including limits on transferring evidence from another setting.

The public-counterpart classification asks whether a pinned schema
contains an object, field or event that supplies part of the required
observation. Of the 71 signals, 32 have partial counterparts. The
question does not apply to 32: 24 derived features, seven
external-knowledge signals and one provenance signal. A keyword search
of the three files identified no counterpart for seven, including
request context, identity-capture session context and several agent
observations. That absence is bounded by the terms searched. S5 gives
the exact locators and explains the limits of partial counterparts; the
inventory records each signal's classification.

Identifier links establish limited relationships. A
storefront-to-merchant identifier link does not show that a payment
through the merchant account settled. A verified identifier from the
same namespace and relevant interval can support a stronger, or
\emph{hard}, link. A descriptor or shared IP address may support only a
weaker, or \emph{soft}, link. Neither alone establishes common control
or malicious intent.

The inventory preserves the manuscript's interpretation rules:
onboarding-link histories are a design analogy (S3.1), and error
estimates retain their cohorts and denominators (Section 7.2). Each
outcome channel retains its decision context (Section 7.1), and each
broader outcome measure requires its own evaluation (Sections 8.2 and
9.1).

\hypertarget{s3.1-how-abuse-mechanisms-motivate-the-records}{%
\subsubsection{S3.1 How abuse mechanisms motivate the
records}\label{s3.1-how-abuse-mechanisms-motivate-the-records}}

Credential exposure is recorded as external knowledge with its location
and time, apart from validity, attacker possession and misuse, and
credential use is compared before and after exposure with precise
timestamps. Access keeps its authentication context, session issuance
and later use, and the changes made after access; a device or browser
fingerprint can connect activity but cannot alone establish common
control. Payout destinations have their own identities and histories:
creation, changes, first observation, validation and reuse across
accounts. Storefronts, statement descriptors, account histories and
verified ownership records are linked, and declared business details
keep their source and time for comparison with activity. End-customer
histories are recorded within a business; linking an end customer's
identity across businesses requires separate authorization, a permitted
purpose, privacy oversight and supporting evidence.

Manuscript Section 5.2 keeps the outcome-label records. The evidence for
the others follows.

\hypertarget{exposure-and-misuse-require-different-evidence}{%
\paragraph{Exposure and misuse require different
evidence}\label{exposure-and-misuse-require-different-evidence}}

Measurements of exposed secrets establish exposure or validity at a
particular time, not unauthorized use. Meli and colleagues found a
median of about 1,800 unique candidate secrets daily. Of the monitored
secrets, 81 percent remained sixteen days after discovery. They did not
test credentials against the services
{[}\protect\hyperlink{ref-Meli2019}{2}{]}. Demir and colleagues
confirmed 1,748 active credentials for 14 services, with payment
services among the largest groups. Traced through nine archived crawl
snapshots, the verified credentials had been publicly exposed for 12
months on average {[}\protect\hyperlink{ref-Demir2026}{3}{]}. That
duration does not show continuous credential validity.

The specification therefore records exposure location and time as
external knowledge, then compares credential use before and after
exposure. Public visibility, validity, attacker possession and misuse
remain separate propositions.

Timing can matter at minute scale. In a deliberate cloud-key exposure
experiment, an actor used the key within minutes while an automatic
quarantine limited its actions. Researchers later removed that
restriction to observe further operations
{[}\protect\hyperlink{ref-Gamazo2023}{4}{]}. Microsoft reported
attackers launching follow-on payment fraud attempts as little as five
minutes after stealing credentials and session cookies through phishing
that relayed the victim's login
{[}\protect\hyperlink{ref-Microsoft2022AiTM}{5}{]}. These observations
motivate precise timestamps without establishing typical attack speed.

\hypertarget{access-needs-a-history}{%
\paragraph{Access needs a history}\label{access-needs-a-history}}

A login challenge is an extra verification step at sign-in. Research
examines challenges triggered by unfamiliar devices and networks
{[}\protect\hyperlink{ref-Doerfler2019}{6}{]}. The design records
authentication context, session issuance and later use. Applying the
same lifecycle to onboarding links, which let businesses complete
account setup, is a design analogy.

Changes after access can identify what an intruder attempted. In
reported OAuth-application abuse, intruders used compromised accounts to
create applications, grant them permissions, add credentials to new and
existing applications and create inbox rules that hid messages from the
user {[}\protect\hyperlink{ref-Microsoft2023OAuth}{7}{]}. The design
records these changes, with links to source observations. Recording
changes to authentication factors and event destinations extends the
same approach.

A device or browser fingerprint can help connect activity but cannot
alone establish common control. One impersonation market sells stolen
user profiles with browser fingerprint data and a browser extension
meant to reproduce them. The study did not buy profiles to test that
software {[}\protect\hyperlink{ref-Campobasso2020}{8}{]}. The records
retain that uncertainty.

\hypertarget{funds-movements-need-destination-histories}{%
\paragraph{Funds movements need destination
histories}\label{funds-movements-need-destination-histories}}

The FBI's payroll-diversion advisory describes attackers changing where
funds are sent {[}\protect\hyperlink{ref-FBI2018}{9}{]}. Payout
destinations therefore have their own identities and histories:
creation, changes, first observation, validation and reuse across
accounts. Trace A applies this analogy to a hypothetical key permitted
to change a payout destination.

Regulatory guidance motivates measures of time between receipt and
onward transfer, and concentration at receiving entities
{[}\protect\hyperlink{ref-FinCEN2014}{10}{]}. It does not validate a
per-account detector.

\hypertarget{business-descriptions-need-corroboration}{%
\paragraph{Business descriptions need
corroboration}\label{business-descriptions-need-corroboration}}

Storefronts, merchant accounts and statement descriptors, the merchant
names cardholders see, need not map one-to-one. ScamMagnifier linked
fraudulent storefronts through merchant identifiers, with a provider
corroborating some links. Its checkouts stopped before purchase and did
not observe settlement {[}\protect\hyperlink{ref-Bitaab2025}{11}{]}.

Studies of abusive merchants describe misdeclared business categories
and movement between banks {[}\protect\hyperlink{ref-McCoy2012}{12},
\protect\hyperlink{ref-Tian2018}{13}{]}. FTC complaints allege pooling
or spreading transactions across accounts to evade monitoring, and
misleading representations
{[}\protect\hyperlink{ref-FTC2020FirstData}{14}--\protect\hyperlink{ref-FTC2019Allied}{16}{]}.
The design consequently links storefronts, descriptors, account
histories and verified ownership records. It retains declared business
details with their source and time for comparison with activity. These
allegations and links motivate investigation; they do not establish
intent or ownership in a new case.

\hypertarget{end-customers-are-not-platform-accounts}{%
\paragraph{End customers are not platform
accounts}\label{end-customers-are-not-platform-accounts}}

A legitimate business may have abusive customers. One provider's public
documentation describes multi-accounting, account sharing, free trial
abuse, account farming and deliberate non-payment for accrued usage
{[}\protect\hyperlink{ref-Stripe2026CAE}{17}--\protect\hyperlink{ref-Stripe2026Farming}{20}{]}.
The customer-evaluations page describes its reuse reasons as reuse
within one business over the past seven days; the account-farming page
does not specify whether its links span businesses.

The specification records end-customer registration and billing
histories, reused identifiers, simultaneous use and unpaid usage within
a business. Linking an end customer's identity across businesses
requires separate authorization, a permitted purpose, privacy oversight
and supporting evidence. This restriction does not prohibit authorized
fraud-prevention aggregation on instruments, devices or credentials
under a stated and logged scope.

\FloatBarrier

\hypertarget{s4.-responsibilities-entities-and-interfaces}{%
\subsection{S4. Responsibilities, entities and
interfaces}\label{s4.-responsibilities-entities-and-interfaces}}

\suppressfloats[t]

The 71 signals identify candidate records, the 24 responsibilities
assign investigation and review functions, and the interfaces preserve
the meaning of what passes between them. The
\href{SUPPLEMENT/S4_FUNCTIONS.csv}{functional responsibilities} file
defines 24 responsibilities without requiring separate processes or
databases. Seventeen were derived with payment-specific mechanisms or
sources in view. One (23) adapts agent-security research, and six (2, 3,
4, 13, 21 and 22) come from cyber threat intelligence work. These
origins describe how the responsibilities were derived, not uniqueness
to payments. The file's \texttt{class} column marks these three groups
DI, AGV and INH.

Responsibilities 16 and 17 cover engineering validation and adaptive
testing of detections, controls and models. Responsibility 21
independently reviews all agent outputs, including intelligence drafts,
hunting and detection candidates, and recommendations; responsibility 22
provides human release. Permission to investigate grants no authority
over accounts, payments or controls. An authorized human may decide
urgent containment before quality review (manuscript Section 6.4),
without releasing the draft into organizational memory.

Responsibilities 1 to 10 scope and conduct an investigation. They
include human tasking, governed tools, the evidence store, hypothesis
testing and entity relationships. Responsibilities 11 to 20 represent
threat knowledge and turn findings into usable outputs: actor, technique
and vulnerability knowledge, the structured threat-knowledge store,
hunting for undetected abuse, detection development, threat-model
integration, recommendations, finished-intelligence synthesis and
evidence lineage. Responsibilities 21 to 24 govern use and reuse:
independent quality review, human release, controlled learning and
security of the investigation system.

The \href{SUPPLEMENT/S4_ENTITIES.csv}{entity inventory} and
\href{SUPPLEMENT/S4_ENTITY_RELATIONS.csv}{entity relationships} define
the entities and the kinds of links between them. Each link records its
source, identifier namespaces and the basis for associating each
identifier with its entity. It retains its effective interval, learned
time, version, confidence, confounders (other explanations for the
link), and permitted purpose and scope. A shared attribute alone does
not establish that one person controls both entities. When a decision
makes an outcome unobservable, the record describes the channel's
status; it does not invent an outcome report.
\protect\hyperlink{complete-storage-requirements}{S4.1} gives the
complete storage requirements, and
\protect\hyperlink{outcome-record-schemas}{S4.2} gives field-level
schemas for the outcome records.

\begin{figure}[tp]

\textbf{Figure S4.1. Investigation as a sequence of tests.} Each query
tests competing explanations within authorized scope. A gateway outside
the agent attaches the query's coverage, which determines which
interpretations remain supported. Draft outputs follow review and
release under manuscript Section 6.4.

\includegraphics[width=0.86\textwidth,height=\textheight]{FIGURES/FIG_S4_1_INVESTIGATION.pdf}

\end{figure}

\Needspace{150pt}

The query interface must show what each query actually covered and what
its result means:

\begingroup\small

\begin{longtable}[]{@{}
  >{\raggedright\arraybackslash}p{(\columnwidth - 2\tabcolsep) * \real{0.2500}}
  >{\raggedright\arraybackslash}p{(\columnwidth - 2\tabcolsep) * \real{0.7500}}@{}}
\toprule\noalign{}
\begin{minipage}[b]{\linewidth}\raggedright
Part of the interface
\end{minipage} & \begin{minipage}[b]{\linewidth}\raggedright
Information or behavior required
\end{minipage} \\
\midrule\noalign{}
\endhead
\bottomrule\noalign{}
\endlastfoot
Request & Task and question; authorized requester; tenant, purpose and
data scope; entity namespaces; event interval and as-of cutoff (only
data the system had learned by then); schema and query versions; row,
time and cost limits; expected observation and required coverage \\
Response & Query/run ID; scope actually evaluated; source and schema
versions; event and ingestion watermarks (how far event time and data
loading had progressed); known gaps and retention horizon for each
source; sampling and access filtering; whether the result is complete,
partial, truncated, denied or unavailable, and any scope-limited flag;
meaning of any count; returned rows or approved aggregate, with the
identifiers of the returned records kept for the run; lineage and
coverage limits \\
Interpretation & A positive observation does not establish how many
cases exist in the population. An empty result counts against a
hypothesis only when the abuse should have produced the queried
observation and the records cover where it would appear. A denied or
unavailable result does not establish adversary behavior. \\
Failure & Enforce limits outside the agent's language model. Stop or
narrow the query after it exceeds a row, time or cost limit, or after an
authorization error. Record incomplete tests and any abstention (a
decision not to conclude). Do not treat missing or not-assessed data as
low risk. \\
\end{longtable}

\endgroup

Some permitted operations also need approval bound to the exact request
payload, target, task, tenant, tool and tool-manifest version. The
manifest is the declared list of the agent's tools. Each approval has an
expiry and a rule governing reuse. At execution, the enforcement gateway
(manuscript Section 6.5) checks that the operation matches the approval
and remains permitted. The restricted agent-and-approval ledger records
both approval and execution.

Control and outcome records distinguish the decision, its inputs and
version, what happened next, the observable channels, actual reports and
any review of correctness. Population error estimates retain their own
cohort, denominator, outcome window, uncertainty and scope; they do not
establish correctness for individual events. Late reports, unobservable
outcomes and false reports have separate fields because they can
coexist. Report presence, channel state and query access also remain
separate. Each temporary delay retains its decision and effective
interval. A derived fact's learned time is when it became queryable,
kept separately from its inputs' latest availability.

A hypothetical timeline illustrates the as-of cutoff. Freeze a
prediction on day 0 using only information then known. A day-30 report
may evaluate it within a stated outcome window, but a model trained on
day 20 cannot use that report. A later training run may use it if the
run's cutoff and label rules permit. Likewise, a relationship learned on
day 10 cannot enter day-0 features merely because its effective date is
earlier.

The \href{SUPPLEMENT/S4_DRAFT_SCHEMA.csv}{draft schema} supplies the
fields that the rule checker, quality reviewers and the releaser use to
assess an intelligence draft. S7 explains the evaluation conditions;
S7.1 gives hypotheses and criteria for E1 to E9. S8 gives the
practitioner study protocol.

\textbf{Returned records a draft does not cite (manuscript Section
6.4).} The checker identifies every record that the case's query runs
returned within a claim's scope but that the draft neither cites nor
lists against it. It attaches those that a fixed rule marks as able to
count against the claim, such as reports, reversals and records of the
benign explanations the claim must exclude, and summarizes the rest by
type and count. Reviewers and the releaser see both beside the claim.
Full checks of claims about named entities examine every marked record
and a recorded random sample of the rest, whose size is fixed in advance
for each claim type, and E8 reports how often that sample contained a
record that counted against the claim. If a share \(q\) of the unmarked
records count against the claim, a random sample of \(m\) finds at least
one with probability at least \(1-(1-q)^m\): 59 records give 95 percent
when \(q\) is 5 percent.

\textbf{Event timing in drafts (manuscript Section 6.4).} Event timing
also needs support. A time asserted by a business, customer or external
source retains its source and uncertainty. Recent receipt does not
establish recent conduct when the event interval is unresolved. Outcome
windows instead use each channel's declared clock-start event.

\hypertarget{complete-storage-requirements}{%
\subsubsection{S4.1 Complete storage
requirements}\label{complete-storage-requirements}}

Table S4.1 summarizes this complete storage contract; S3.1 and
manuscript Section 5.2 give the motivating evidence. Each block groups
requirements by the information it preserves. The stores distinguish
record meanings and need not be separate databases.

\begin{table}[tp]
\textbf{Table S4.1. What each analytical store preserves.} The full
contract is in S4.1. These are requirements for the proposed platform,
not claims about an existing provider. The structured threat-knowledge
and organizational-memory stores form their own planes, and the other
stores hold evidence-plane records. Clusters in the entity-relation
graph remain threat-knowledge items. External-knowledge records support
only what their source observed or described, and ledger entries only
what happened during the investigation itself.\par
{\small\renewcommand{\arraystretch}{1.04}
\begin{tabular}{@{}
  >{\raggedright\arraybackslash}p{(\columnwidth - 2\tabcolsep) * \real{0.2500}}
  >{\raggedright\arraybackslash}p{(\columnwidth - 2\tabcolsep) * \real{0.7500}}@{}}
\toprule\noalign{}
\begin{minipage}[b]{\linewidth}\raggedright
Store or shared requirement
\end{minipage} & \begin{minipage}[b]{\linewidth}\raggedright
Required properties
\end{minipage} \\
\midrule\noalign{}

Observation log & Source records and revisions, identifiers,
environment, tenant scope and lineage; separate event, source-recording,
ingestion and query-availability times; repeated delivery does not
duplicate the event \\
Entity-relation graph & Typed identifier links with source, namespace,
effective interval, learned time, version, confidence, benign
explanations, and permitted purpose and scope; versioned cluster
membership; end customers scoped to a business \\
Feature store & Recomputable inputs and versions; decision-time features
exclude later information; later labels use the evaluation's outcome
window and maturity; features inherit their inputs' agent-output and
released-item links \\
Label store & Report source, type and attributed entity; occurrence,
report and arrival times; appended value history; typed links to every
decision that may have shaped the report; links to separate truth
assessments and any informing agent output or released item; report
presence, channel state and query access kept apart \\
Assessment store & Model or reviewer identity and version; review
selection route and probability; links to any informing agent output or
released item; unknown, not-assessed and failed-open states kept
distinct from low risk; failed-open means an assessment failed and
activity continued \\
Control-outcome store & Decision, authority, reason and version;
reproducible available inputs; alternative-decision probabilities where
applicable; human assignment and what the decider was shown; links to
any informing agent output or released item; effects on each outcome
channel, fixed with the decision; for action before review, the urgent
route, impact class and later review outcome; correctness reviews with
their selection route, separate from population error estimates \\
Structured threat-knowledge store & Revisable relationships and
conclusions with confidence, input and rule versions, writer and status;
agent proposals marked as unreleased context; links to informing agent
outputs and released items \\
Organizational-memory store & Human-released items with approved claims
and their stated confidence, version, scope, expiry, reviewer findings,
lineage and correction links; exonerating items have independent
counterexamples, a re-test date and the evidence and cluster versions
whose change expires them early; proposed lessons wait outside it with
their counterexample tests, and released lessons keep monitoring and
withdrawal records \\
External-knowledge store & Source, type, times observed and published,
validity period, retrieval time and license; a link to the released item
whose identifier the returning content bears \\
Agent-and-approval ledger & Model, instructions, inputs, tool calls;
approver, binding, expiry and reuse rule of each approval; approved and
executed payloads, tool-manifest version and hash at approval and at
execution, memory writes, outbound-traffic decisions; rule-checker and
reviewer findings per draft version; marked untrusted content and
exposed agent output; logged agent-output and released-item link
removals \\
Query coverage record & Gateway-written entities and intervals,
requested and evaluated; as-of cutoff, schema versions, sampling,
per-source watermarks, known gaps and retention horizon, access
filtering and limits reached; derived result status; not writable by the
agent \\
Shared time and graph requirements & Reconstruct what was available at
each decision; preserve reporting clocks and delays, retention limits,
account histories and report-to-transaction links \\

\bottomrule\noalign{}
\end{tabular}
}
\end{table}

\needspace{4\baselineskip}

\textbf{Observation log}

\begin{itemize}
\tightlist
\item
  Preserve records and revisions without overwriting earlier records.
  Keep source, event and delivery identifiers, schema version, entities,
  source lineage, and actor and request identifiers where available. A
  record delivered more than once takes effect only once.
\item
  Record each observation's environment (production, test, synthetic,
  replay or emulation), marked when the record is created, and its
  tenant scope: the business or other boundary within which it may be
  accessed and used. Keep each record's source trust mark. Mark content
  that businesses, customers or outside sources supplied as untrusted.
  Both marks survive normalization and storage.
\item
  For content fetched from a source the subject controls, such as a
  storefront or business website, record the vantage point of each
  fetch: network origin, region and client. A subject can serve
  different content to different visitors, including reviewers.
\item
  Mark every label, assessment, control outcome, draft, proposed lesson
  and released item that a test, replay or emulation creates, or that
  derives from such a record or from synthetic data, with the same
  environment value when it is created. Keep these records out of
  production labels, intelligence and lessons. Record any authorization
  for their evaluation use, with who gave it, the evaluation and its
  scope.
\item
  Keep event or effective time (when something happened or took effect)
  with its source and uncertainty bound. Record the source system's
  recorded time and the platform's ingestion time separately.
\item
  For each store, retain when each record and revision became visible to
  queries.
\item
  A derived record keeps its inputs' latest learned time separately from
  the time it first became queryable. Its own learned time is the latter
  and cannot precede any input's learned time.
\end{itemize}

\needspace{4\baselineskip}

\textbf{Entity-relation graph}

\begin{itemize}
\tightlist
\item
  Label each identifier link by kind and as hard (stronger) or soft
  (weaker). Retain its source, identifier namespaces, effective
  interval, learned time, version, confidence and possible benign
  explanations (confounders).
\item
  Record why each identifier is believed to belong to its entity.
\item
  Keep versions of linked-entity clusters and their membership rules.
  Record end customers separately for each business. Record each link's
  permitted purpose and scope. A link across businesses also records its
  authorization for that identifier type (manuscript Sections 5.2 and
  5.5).
\end{itemize}

\needspace{4\baselineskip}

\textbf{Feature store}

\begin{itemize}
\tightlist
\item
  Keep inputs traceable and values recomputable, with the feature
  definition version, the entities and the computation window.
  Decision-time features use only information available at that
  decision.
\item
  A feature computed over a cluster or other threat-knowledge item
  records that item's version, writer and status. If an agent proposed
  the item, or the item links to an agent output or released
  organizational-memory item that informed it, the feature links to that
  output or item and counts as dependent on it.
\item
  Join later outcomes as evaluation labels under an explicit outcome
  window and maturity period. The outcome window is the period after the
  decision in which an outcome must occur to count.
\end{itemize}

\needspace{4\baselineskip}

\textbf{Label store}

\begin{itemize}
\tightlist
\item
  \textbf{Report history:} retain report type, source, attributed
  transaction or entity and lag. Keep three times: when the reported
  outcome occurred, where known, with its uncertainty; when the report
  was made at its source; and when it arrived at the platform. Preserve
  each value and its learned time, such as a dispute being filed,
  answered, then won or lost, and record whether each won or lost status
  followed adjudication. Append each status change, reversal or
  correction as a new record that references the earlier one. Never edit
  an earlier record, so that what the platform believed at each time can
  be reconstructed.
\item
  \textbf{Truth assessment:} link any assessment that the report is
  mistaken or deliberately false. Keep that assessment in the assessment
  store with its model or reviewer identity and version, how the report
  was selected for review and, for a sampled review, its selection
  probability.
\item
  \textbf{Dependencies:} link every decision that may have shaped the
  report, and any agent output or released organizational-memory item
  that informed the label. An outcome channel is a route, such as a
  chargeback, through which an outcome can be reported. A report may
  link to several decisions, such as a blocked attempt, a retry and the
  approval of the paid attempt, and one decision, such as an account
  restriction, may link to many reports. Each link states the decision's
  role, such as a decision on the reported item, on an earlier related
  attempt, or on an account or payout that contains it, and the basis
  for the link. These links do not change how a measure counts reports
  (Graph semantics). Record the writer's case or session and the agent
  outputs and released items open in it when the record was written.
\item
  \textbf{Channel status:} for each affected transaction, payout or
  account application, keep report presence, channel state and query
  access separate at each read time. These statuses may be computed when
  read from the stored decision, report and channel records. When a
  draft cites a status, store that status with an identifier and its
  read time.
\item
  \textbf{State history:} a channel is open, temporarily delayed or not
  yet opened, permanently closed by a recorded decision, reopened or
  extended by a recorded rule event, or unknown. Keep every transition's
  effective interval and the decision ID, or the rule event's source and
  version, that caused it.
\item
  \textbf{No report:} an open channel remains pending before its
  reporting window ends. It matures with no report only if it stayed
  open throughout the window and the records cover that window. If a
  recorded decision or rule event changed the channel's meaning, the
  matured absence carries that decision's ID, or the rule event's source
  and version, and is read under the changed meaning, not as a negative
  for the original outcome. A channel with no defined window never
  supports a negative by maturity alone. A matured absence is
  provisional: it records the rule version under which it matured and
  returns to pending if a later rule event reopens the window.
\item
  \textbf{Unavailable records:} mark inaccessible records unavailable
  without overwriting known channel state or earlier reports.
\item
  \textbf{Clocks and maturity:} for a temporary delay, keep the expected
  and actual clock-start events when known. Store the reporting window
  and clock-start event with the channel, including their source and
  version. Each channel record also names its rail, the attributes its
  window depends on, such as whether the debited account is a consumer
  or business account or the type of claim, and whether the channel
  adjudicates report truth. Store the maturity used by an evaluation
  with that evaluation.
\end{itemize}

\textbf{Program status, revised negatives and reopened reports
(manuscript Section 5.3).} A decision can change a channel's meaning, as
a liability shift after successful authentication does (manuscript
Section 7.1). A merchant's network program status can also change that
meaning: under one program, a merchant loses the liability shift on
domestic authenticated payments until it exits the program
{[}\protect\hyperlink{ref-Stripe2026MonitoringPrograms}{21}{]}. The
records keep that status as it stood at decision time, because the guide
does not say whether the loss reaches payments authenticated before the
merchant entered the program. A later program entry is recorded as a
rule event, and the liability status of payments authenticated before it
remains unknown. A negative read from a matured absence cites the rule
version and read time under which it matured and is revised when a later
transition reopens the channel. A report's resolution is not always
final either: in rare cases an issuer files a further dispute on the
same payment after acquiring new information, or a lost dispute later
becomes a late win {[}\protect\hyperlink{ref-Stripe2026Disputes}{22}{]}.

\textbf{Bank-debit and other rails (manuscript Section 7.1).} A SEPA
Direct Debit dispute within eight weeks of the debit is honored with no
questions asked, and such disputes are final, with no process for appeal
{[}\protect\hyperlink{ref-Stripe2026SEPADebit}{23}{]}. After eight weeks
and up to 13 months, a customer can dispute only a debit considered
unauthorized, and the customer's bank decides; the window that applies
therefore depends on the claim, and an absence at eight weeks is not a
negative for unauthorized-debit fraud
{[}\protect\hyperlink{ref-Stripe2026SEPADebit}{23}{]}. Inside its
window, an ACH return is refunded without question, and the ACH Network
has no appeals process of its own
{[}\protect\hyperlink{ref-Stripe2026ACHTerms}{24}{]}. A SEPA dispute
honored with no questions asked is not by itself a fraud allegation. For
ACH returns, the provider's terms give 60 days for consumer accounts and
2 days for business accounts, counted from the customer's bank statement
{[}\protect\hyperlink{ref-Stripe2026ACHTerms}{24}{]}, while the Nacha
return deadlines that banks apply run from the settlement date of the
original debit {[}\protect\hyperlink{ref-LeadBank2026ACHReturns}{25}{]};
a channel record therefore names the rule from which it takes its clock
start. Banks may also accept late returns outside the window
{[}\protect\hyperlink{ref-Stripe2026ACHTerms}{24}{]}. The product scope
also lists instant push payments and stablecoin payments (S2); their
reporting channels, if any, follow each rail's own rules and windows,
which this paper did not review.

\textbf{Limits of the rail review (manuscript Section 9.2).} Bank debits
were compared only with two of the provider's pages and one bank's
return-code table (manuscript Section 7.1); of the bank-debit candidates
in S2, \mbox{PAY-07} was judged covered and \mbox{MM-03} was not
reviewed. Instant push payments and stablecoin payments, which are also
in the product scope, were not reviewed rail by rail; their mechanisms
\mbox{MM-07} and \mbox{MM-12} were reviewed, but not their reporting
channels (S2).

\needspace{4\baselineskip}

\textbf{Assessment store}

\begin{itemize}
\tightlist
\item
  Retain the model or reviewer identity and version. Keep unknown,
  not-assessed and failed-open states distinct from low risk.
  Failed-open means the assessment failed and activity continued.
\item
  Record when each assessment was made and frozen (its freeze time) and
  the as-of time of its inputs. For each model version, record its
  training and model-selection cutoff and the label rules it used. A
  label that became available after that cutoff cannot enter that
  version.
\item
  Never store an assessment as an observation of the conduct it
  assesses.
\item
  Link each assessment, whether a model, a detection or a reviewer
  produced it, and each detection or model version, to any agent output
  or released organizational-memory item that informed it. Record the
  writer's case or session and the agent outputs and released items open
  in it when the record was written.
\item
  For each detection or model version, link the validation records of
  responsibilities 16 and 17, with their version, date and result.
\end{itemize}

\needspace{4\baselineskip}

\textbf{Control-outcome store}

\begin{itemize}
\tightlist
\item
  \textbf{Decision:} retain the decision, who had authority to make it,
  the rule or model version, and the reason.
\item
  \textbf{Inputs and policy:} log the time of the decision and the
  inputs it used. Keep versions from which every other input available
  at that time can be reconstructed. Record the score, threshold or rule
  that produced the decision, and whether the policy was deterministic
  and, otherwise, the probability of each alternative. Identify any
  input that a decider, automated or human, used but that was not
  logged.
\item
  \textbf{Human decisions:} retain who decided, how the case was
  assigned and what the decider was shown.
\item
  \textbf{Dependencies:} link any agent output and any released
  organizational-memory item that informed the decision, so the decision
  can be revisited when that item expires or is corrected or retracted.
  Record the writer's case or session and the agent outputs and released
  items open in it when the record was written.
\item
  \textbf{Correctness:} record reviewed correctness when available, with
  how the case was selected for review (for example an appeal, a
  complaint, a sampling rule or review of every case) and, for a sampled
  review, its selection probability; otherwise record unknown.
\item
  \textbf{Error estimates:} keep separate links to each estimate, with
  its control version, cohort, outcome window, denominator, uncertainty
  and scope. The cohort is the fixed set of decisions the estimate
  covers.
\item
  \textbf{Channel effects:} record the decision's effect on each outcome
  channel as closed, delayed, changed in meaning, none or unknown. Fix
  these effects when the decision is recorded. Record later state
  transitions, including rule events, as separate records that never
  edit the decision.
\item
  \textbf{Restrictions:} distinguish allowing and blocking from
  restrictions, such as a hold that may later be released.
\item
  \textbf{Urgent containment:} for an action decided before quality
  review (manuscript Section 6.4), record that route, the impact class
  (none, degraded service or outage), any notification and restoration
  path, the review deadline and the later review's outcome, including
  any reversal.
\end{itemize}

\needspace{4\baselineskip}

\textbf{Structured threat-knowledge store}

\begin{itemize}
\tightlist
\item
  For inferred relationships, clusters, technique mappings and
  conclusions, retain confidence and the input and rule versions that
  produced them.
\item
  Record the writer as a correlation engine, agent proposal or person,
  and status as proposed, accepted or retracted. Record the person or
  accountable process that accepted or retracted the item, and when.
\item
  Retrieve an agent-proposed item only when it is marked as an
  unreleased agent inference. It never counts as corroboration.
\item
  Link an item that a person or correlation engine wrote to any agent
  output or released organizational-memory item that informed it. Record
  the writer's case or session and the agent outputs and released items
  open in it when the item was written. The item counts as dependent on
  those outputs and items.
\end{itemize}

\needspace{4\baselineskip}

\textbf{Organizational-memory store}

\begin{itemize}
\tightlist
\item
  Keep released items only, with the releaser, version hash, released
  claims and the confidence stated for each, scope, expiry, reviewer
  findings and originating draft.
\item
  Retain supersession and retraction links, and lineage from every cited
  observation to its root record. Record every later status record for a
  released item, such as a monitoring finding, an early expiry, a
  retraction or a withdrawal, with who or what wrote it and when.
\item
  For benign or exonerating items, keep a re-test date, a source of
  counterexamples independent of the writer, and the versions of the
  cited evidence and of any cluster the item concerns. A revision of
  that evidence or a change in that cluster's membership expires the
  item early (manuscript Section 6.5). Record each re-test, with its
  counterexamples' sources, its result and its date.
\item
  Read each item as context or a hypothesis, never as an observation.
\item
  Proposed lessons remain outside organizational memory until testing,
  quality review and human release. For each one, record its writer,
  proposal time, queue state, intended scope and version. Record each
  test against counterexamples, with the counterexamples' sources and
  the result. After release, record monitoring findings and any
  withdrawal, each with its writer and time, and the reason for a
  withdrawal (manuscript Section 7.3). The requirement for
  responsibility 23 in the functions file calls a withdrawal rollback;
  the two terms name the same record.
\end{itemize}

\needspace{4\baselineskip}

\textbf{External-knowledge store}

\begin{itemize}
\tightlist
\item
  Retain source, type, version, when the source observed and published
  the fact, validity window, license and retrieval time. The validity
  window is the period in which the fact holds. For a catalog of abuse
  techniques, the version is the fixed version that mappings cite.
\item
  Keep external knowledge distinct from platform observations.
\item
  Record the identifier attached to each released item shared outside
  the platform. When external content returns with that identifier, link
  it to that item. It counts as dependent on the item and on any agent
  output behind it, not as independent corroboration.
\end{itemize}

\needspace{4\baselineskip}

\textbf{Agent-and-approval ledger}

\begin{itemize}
\tightlist
\item
  For each run, record the human assignment or the standing task
  template and version that authorized it, with its permitted data,
  entities, period, current period and outputs. The current period
  bounds claims about current conduct.
\item
  For each tool call, preserve the provenance of inputs that reached the
  agent's language model.
\item
  For each approval, keep the approver, the target, task, tenant and
  tool it is bound to, its expiry, its reuse rule and any token issued.
  Keep hashes of the approved and executed payloads, and the
  tool-manifest version and hash in force at approval and at execution.
  The manifest is the declared list of the agent's tools.
\item
  Record memory writes and decisions on outbound traffic.
\item
  Record each agent output, such as an intelligence draft and each of
  its versions, a hunting or detection candidate or a recommendation,
  with its identifier, version hash, run and creation time. It is an
  observation of what the agent did, not evidence of the activity under
  investigation.
\item
  Retain model identity and version, instruction version and decoding
  settings.
\item
  Mark content that came from untrusted input as untrusted when it is
  read back. Distinguish copied untrusted content, agent output exposed
  to untrusted content and agent output without that exposure; exposure
  takes precedence.
\item
  Record each rule-checker result and quality-review finding with the
  draft's version hash, the rule version or reviewer, including the
  version of the rule that marks uncited returned records, and the
  draft's resulting status: returned, released or expired. For each
  quality review, record the materiality tier applied to each claim, the
  sampling rate, the cited records the reviewer examined and the random
  sample of unmarked uncited returned records it drew, whose size is
  fixed in advance for each claim type. Record when each draft version
  entered the review queue and when it would expire from it. For each
  review, record its start and end times.
\item
  Record each removal of an agent-output or released-item link from a
  label, control decision, detection, assessment or threat-knowledge
  item, with the person who removed it, that person's release authority,
  the time, the reason and the identifier of the removed output or item.
\item
  Store references to card data, identity documents and other regulated
  data, not copies (manuscript Section 5.5).
\end{itemize}

\textbf{What makes the agent investigable (manuscript Section 6.5).}
Ordinary source observations retain their own trust marks and are not
assigned an agent-output category. The ledger makes the agent itself
investigable. For each tool call it retains the provenance of inputs
reaching the model, approved and executed payload hashes, memory writes
and outbound-traffic decisions. It also records model identity, version,
instructions and decoding settings, with record classes intact.

\textbf{Privacy and operating constraints (manuscript Section 5.5).} A
deployment must specify retention and deletion for each record class
before collecting it. Append-only logs preserve history but must still
permit required deletion of personal payloads, derived data, backups and
identifying lineage. Encryption-key destruction is sufficient only if
recoverable keys and plaintext copies are also removed. Retained audit
structure must not preserve information required to be deleted.
Authorized holds record their scope and retention rule.

The agent-and-approval ledger uses references to card data, identity
documents and other regulated data rather than copies. Tokens or keyed
hashes may be used only where the applicable standard permits. The
ledger should not become a new sensitive-data store or widen
payment-card security scope.

Use of external model providers depends on their processing terms and
permitted processing locations. Drafts concerning suspicious-activity
reporting remain subject to the platform's confidentiality and
subject-notification rules.

Cross-business linking requires a stated purpose and authorization for
each identifier type. Permission to use a payment-instrument link does
not automatically authorize linking devices or verified people. These
deployment obligations constrain the specification; this paper does not
determine which legal basis applies in a particular deployment.

\needspace{4\baselineskip}

\textbf{Query coverage record}

\begin{itemize}
\tightlist
\item
  For each query run, retain its request: the task, run and question,
  the authorized requester, the tenant, purpose and data scope, the
  query version, the row, time and cost limits set, the expected
  observation and the required coverage.
\item
  For each query run, retain the coverage record the gateway attached:
  requested and evaluated entities with their tenant and identifier
  namespaces, requested and evaluated event intervals, as-of cutoff,
  source and schema versions, sampling, each source's event and
  ingestion watermarks, known gaps and retention horizon, access
  filtering, and any row, time or cost limit reached.
\item
  Keep the derived result status (complete, partial, truncated, denied
  or unavailable), the version of the rules that derived it and any
  scope-limited flag. The agent cannot write these fields.
\item
  For each query run, keep the identifiers of the records it returned,
  including those behind an approved aggregate, so that the checker can
  find returned records that a draft neither cites nor lists against a
  claim, attach those its fixed rule marks and summarize the rest
  (manuscript Section 6.4). Keep identifiers, not copies, under each
  record's retention and deletion rules (manuscript Section 5.5).
\item
  Link each tested hypothesis and each claim of absence to the query
  runs behind it. Record each hypothesis's status, including
  inconclusive and unassessable, apart from the query results. Record
  each incomplete test and each abstention, a decision not to conclude.
\end{itemize}

\needspace{4\baselineskip}

\textbf{Time semantics}

\begin{itemize}
\tightlist
\item
  Use timestamps precise enough to join a key's exposure to its use
  within minutes, with each timestamp's clock source.
\item
  For each payment, keep authorization, capture and processing times
  separately, and any service or event date that a channel rule uses as
  its clock start. Keep the authentication result and the reported
  liability status separately, each with its source, rule version and
  learned time. Keep any network program status that removes the
  liability shift as it stood at decision time. Record a later program
  entry as a rule event, and leave the liability status of payments
  authenticated before it unknown. An unknown liability status remains
  unknown.
\item
  Preserve label delay and what information was available as of each
  time.
\item
  Retain account age, dormancy and how long funds stay before moving on.
\item
  For each record class, keep each record until the longest applicable
  reporting window, counted from its channel's clock-start event, or the
  defined maturity period has ended, whichever ends later, plus the
  investigative lookback period. Follow the retention and deletion rules
  in manuscript Section 5.5. Record resulting gaps as coverage limits.
  Record each authorized hold with its scope and retention rule.
\end{itemize}

\needspace{4\baselineskip}

\textbf{Graph semantics}

\begin{itemize}
\tightlist
\item
  Link closed accounts to suspected successor accounts, and storefronts
  to merchants.
\item
  Link each outcome report to the transaction, payout or account
  application it concerns, with the basis for the link and whether it is
  verified. Attribute outcomes through that link to the platform
  account, connected account and end customer, so that a measure counts
  each affected transaction once. A report's links to the decisions that
  may have shaped it do not change this count. Keep each transaction's
  acquiring bank and the static part of its statement descriptor, so
  that a network's monitored unit can be reconstructed.
\item
  Link each wallet to every party it has allowed to spend its funds.
\end{itemize}

\hypertarget{outcome-record-schemas}{%
\subsubsection{S4.2 Field-level schemas for outcome
records}\label{outcome-record-schemas}}

The
\href{SUPPLEMENT/S4_OUTCOME_RECORDS/outcome_records.schema.json}{outcome-record
schema} defines, in JSON Schema, the fields of the records that the
first requirement needs. It covers eight record types:

\begingroup\small

\begin{longtable}[]{@{}
  >{\raggedright\arraybackslash}p{(\columnwidth - 2\tabcolsep) * \real{0.2500}}
  >{\raggedright\arraybackslash}p{(\columnwidth - 2\tabcolsep) * \real{0.7500}}@{}}
\toprule\noalign{}
\begin{minipage}[b]{\linewidth}\raggedright
Record
\end{minipage} & \begin{minipage}[b]{\linewidth}\raggedright
What it holds
\end{minipage} \\
\midrule\noalign{}
\endhead
\bottomrule\noalign{}
\endlastfoot
OutcomeChannel & A route for reporting outcomes: its rail, clock-start
event, reporting window with source and rule version, the attributes the
window depends on, and whether the channel adjudicates report truth \\
ControlAction & A decision's targets, time, authority and version, and
its effect on each channel: closed, delayed, changed in meaning, none or
unknown \\
ChannelStateTransition & One change in a channel's state for one item,
with its effective interval and the decision or rule event that caused
it \\
LabelEvent & An outcome report or one revision of it, with its amount
and whether the report states that amount or the platform estimated it,
when the outcome occurred, when it was reported at its source, when it
arrived and when it became queryable \\
LabelEventDecisionLink & A link from a report to one decision that may
have shaped it, with the decision's role, the basis for the link and
whether it is verified \\
ChannelStatusReading & Report presence, channel state, query access and
no-report status for one item and channel at one read time, stored when
a draft cites it \\
ItemAttribution & The accounts an item belongs to, and the acquiring
bank and static statement descriptor from which a network's monitored
unit can be rebuilt \\
EntityLink & A link between entities, such as a closed account and a
suspected successor, with its basis, confidence and other possible
explanations \\
\end{longtable}

\endgroup

Every record is append-only: a correction is a new record that
references the earlier one. A time that is not known is recorded as
unknown, with a reason, and is never left out. The schema checks each
record on its own. Fourteen rules check records against each other:

\begin{enumerate}
\def\labelenumi{\arabic{enumi}.}
\tightlist
\item
  \mbox{OR-01}: a record identifier is used once. A record changed in
  place under its own identifier passes every rule, so the store must
  prevent edits.
\item
  \mbox{OR-02}: a correction or revision references an earlier record of
  the same type and, for a report, of the same report.
\item
  \mbox{OR-03}: every channel, decision and report that a record cites
  exists.
\item
  \mbox{OR-04}: a report's times are in order where they are known:
  occurred, reported, arrived, learned.
\item
  \mbox{OR-05}: a won or lost report states whether that result followed
  adjudication.
\item
  \mbox{OR-06}: a decision gives one effect per channel. A delay gives
  its expected clock start, and a change in meaning says what changed.
\item
  \mbox{OR-07}: a permanent closure cites the decision that caused it,
  and a reopening or extension cites the rule event. A cited decision
  must target the item, record the matching effect and precede the
  transition.
\item
  \mbox{OR-08}: no report can arrive through a channel while a closure
  of that channel for that item is in force; a revision counts from its
  report's first arrival. Such a report means the item link or the
  closure is wrong.
\item
  \mbox{OR-09}: a link's role fits the decision: a decision on the
  reported item, on an earlier related attempt, or on an account or
  payout that contains the item.
\item
  \mbox{OR-10}: a reading's channel state is computed only from the
  transitions the system had learned by the read time.
\item
  \mbox{OR-11}: a reading's report presence matches the reports learned
  by the read time. Presence is unavailable only when query access was
  denied or unavailable, and it does not overwrite channel state.
\item
  \mbox{OR-12}: a no-report status follows from the channel's state and
  window. An absence matures only if the window is defined and has
  ended, the channel was open throughout it, with no gap in its recorded
  state, and query access was full. A matured absence records the rule
  version it matured under and cites any decision that changed the
  channel's meaning.
\item
  \mbox{OR-13}: each reported item has exactly one live attribution to
  its accounts, so a measure counts it once.
\item
  \mbox{OR-14}: a link resting only on a shared attribute is neither
  verified nor high-confidence, and a suspected-successor link starts
  from an account with a recorded closure.
\end{enumerate}

The \href{SUPPLEMENT/S4_OUTCOME_RECORDS/conformance_valid.json}{valid
conformance case} builds one synthetic history: a blocked card attempt,
an approved retry that draws a fraud chargeback the merchant later wins,
a hold on the account that received the payment, and a payment that
draws no report. An earlier account with the same principal name was
closed; the new account is linked to it as a suspected successor on that
name alone, and a storefront is linked to the new account by its
checkout identifier. Its 27 records pass. Each of the 24
\href{SUPPLEMENT/S4_OUTCOME_RECORDS/conformance_invalid.json}{invalid
cases} changes that history in one way and must fail exactly one check;
together they exercise the schema check and all fourteen rules.
\texttt{tools/outcome\_conformance.py} runs the cases using only the
Python standard library, and \texttt{CHECK\_CANDIDATE.py} runs it as
well. The same tool checks any records file in this format.

Channel names, windows and rule versions in the cases are illustrative,
not any network's rules. The schema follows the label-store,
control-outcome-store and graph-semantics requirements in S4.1 and the
relations in the entity-relationships file. A schema and passing
synthetic cases do not show that the records improve investigations.
\protect\hyperlink{investigation-record-schemas}{S4.4} gives the schemas
for query runs, which the second requirement needs, and for release
records. Denominators are outside these records: which payments a
measure counts, their amounts and their capture or settlement times come
from the payment ledger, which this section does not specify. Whether a
denominator read was complete is a query-coverage record.

\hypertarget{outcome-record-cost}{%
\subsubsection{S4.3 Cost of the outcome records on synthetic
data}\label{outcome-record-cost}}

\href{tools/outcome_cost.py}{\texttt{tools/outcome\_cost.py}} builds one
million synthetic payments from the
\href{SUPPLEMENT/S4_OUTCOME_RECORDS/conformance_valid.json}{valid
conformance case}, replacing identifiers and shifting times, so each
record has its template's fields. Before measuring, it checks a sample
of generated records against the schema and \mbox{OR-01} to
\mbox{OR-14}. \href{SUPPLEMENT/S4_OUTCOME_RECORDS/cost_results.json}{The
results} come from one run on the machine below.

Assumptions, illustrative and not any platform's values: payments spread
over 90 days; 100 payments per account; two outcome channels per
payment, fraud chargeback and early fraud warning; 1 percent of payments
blocked; 0.5 percent of allowed payments draw a fraud chargeback with
three report events and one decision link; 1 percent of accounts held
and 1 percent linked to a suspected predecessor; one cited channel
reading per blocked payment.

\begingroup\small

\begin{longtable}[]{@{}
  >{\raggedright\arraybackslash}p{(\columnwidth - 4\tabcolsep) * \real{0.3000}}
  >{\raggedleft\arraybackslash}p{(\columnwidth - 4\tabcolsep) * \real{0.3400}}
  >{\raggedleft\arraybackslash}p{(\columnwidth - 4\tabcolsep) * \real{0.3600}}@{}}
\toprule\noalign{}
\begin{minipage}[b]{\linewidth}\raggedright
Record
\end{minipage} & \begin{minipage}[b]{\linewidth}\raggedleft
Records
\end{minipage} & \begin{minipage}[b]{\linewidth}\raggedleft
Mean size (bytes, JSON)
\end{minipage} \\
\midrule\noalign{}
\endhead
\bottomrule\noalign{}
\endlastfoot
ChannelStateTransition & 2,000,000 & 411 \\
ControlAction & 1,000,200 & 422 \\
ItemAttribution & 1,000,000 & 397 \\
LabelEvent & 15,372 & 718 \\
ChannelStatusReading & 9,947 & 378 \\
LabelEventDecisionLink & 5,124 & 230 \\
EntityLink & 100 & 558 \\
OutcomeChannel & 2 & 361 \\
\end{longtable}

\endgroup

Each payment needs 4.03 records. Per million payments: 1,656 MB as JSON
lines; 79 MB compressed with gzip level 6; 2,559 MB in SQLite 3.53.1,
storing each record's JSON with one index on item, channel, record type,
learned time and effective time. On a two-core x86-64 virtual machine
with Python 3.13.18: 10,000 reads of one payment's fraud-chargeback
state and report count, at a cutoff after all records, took a median of
17 microseconds and a 99th percentile of 57; one day's cohort of 11,052
payments took 0.3 seconds. Loading took 125 seconds and indexing 2.7.

Limits: template text repeats, so compression overstates what real
records achieve, and field values are longer than coded production
values would be. Each added outcome channel adds one channel-state
record per payment. The store is one SQLite file on one machine, not a
production database. Query runs (C2) and release records (C3) are
measured in S4.4.

\hypertarget{investigation-record-schemas}{%
\subsubsection{S4.4 Field-level schemas for query runs and
release}\label{investigation-record-schemas}}

The
\href{SUPPLEMENT/S4_INVESTIGATION_RECORDS/investigation_records.schema.json}{investigation-record
schema} defines the records the second and third requirements need. It
covers eight record types:

\begingroup\small

\begin{longtable}[]{@{}
  >{\raggedright\arraybackslash}p{(\columnwidth - 2\tabcolsep) * \real{0.2500}}
  >{\raggedright\arraybackslash}p{(\columnwidth - 2\tabcolsep) * \real{0.7500}}@{}}
\toprule\noalign{}
\begin{minipage}[b]{\linewidth}\raggedright
Record
\end{minipage} & \begin{minipage}[b]{\linewidth}\raggedright
What it holds
\end{minipage} \\
\midrule\noalign{}
\endhead
\bottomrule\noalign{}
\endlastfoot
Actor & A person, or a process at one version, that writes records: a
source system, pipeline, model, rule, agent, the gateway or a checker \\
QueryRun & One query and the coverage record the gateway attached:
requested and evaluated entities and intervals, the as-of cutoff, each
source's watermarks, known gaps, retention horizon and sampling,
filtering, limits reached, the records returned and the derived
statuses \\
Observation & A stored record with its class, event interval, learned
time, producer, root observations, inputs and links to any agent output
that informed it \\
DraftOutput & An agent draft: its exact version, expiry, the task's
current period, and each claim with its supporting and contradicting
records, query runs, context, and any items whose unreported outcome it
relies on, with their channel readings \\
OutputExposure & A period in which a person had an agent output open
about an entity \\
ReleaseDecision & A person's decision on one exact draft version: the
claims released with their confidence, scope, expiry and reviewer
findings \\
ReleasedItem & A released conclusion in organizational memory, usable as
context only \\
MemoryRead & One read of a released item \\
\end{longtable}

\endgroup

Fourteen rules check records against each other:

\begin{enumerate}
\def\labelenumi{\arabic{enumi}.}
\tightlist
\item
  \mbox{IR-01}: a record identifier is used once.
\item
  \mbox{IR-02}: every record a record cites exists, a cited query run is
  a query run, and a producer is an actor.
\item
  \mbox{IR-03}: a gateway that is not the requester records a query
  run's coverage.
\item
  \mbox{IR-04}: the evaluated entities and interval lie within those
  requested, and the as-of cutoff is not after the query.
\item
  \mbox{IR-05}: a run's statuses follow from its coverage fields. By
  Section 6.2's source rules, a run is partial when a source was
  sampled, has a known gap in the requested interval, or does not reach
  from its retention horizon to its event watermark across that
  interval, and truncated when a limit was reached. Three rules are
  added here: a narrower evaluated scope, access filtering, or a source
  whose loading has not reached the end of the requested interval (its
  ingestion watermark) also gives partial. A denied or unavailable
  source gives that status. A run is complete only when nothing else
  applies.
\item
  \mbox{IR-06}: a run returns only records learned by its as-of cutoff
  whose event intervals overlap the evaluated interval.
\item
  \mbox{IR-07}: a derived record is learned no earlier than its inputs,
  and its root observations are exactly its inputs' roots. Only a source
  observation, score or decision can be its own root. A source
  observation is not learned before its event began.
\item
  \mbox{IR-08}: an absence claim cites at least one query run that is
  complete, not scope-limited, returned nothing and covers the claim's
  entities and interval. No cited run returned a record about those
  entities in that interval. A cited run that does not bear on that
  interval does not count against the claim.
\item
  \mbox{IR-09}: a claim's support consists of observation records about
  the claim's entities. Every record a claim cites, as support,
  contradiction, context or query run, exists by the time the draft was
  written. A supporting record is dependent, and adds no root, if it is
  linked to an agent output, names no producer, or lists no inputs and
  is not a source observation. The restriction runs through inputs: a
  derived record adds only the roots reachable through records that are
  not dependent. The claim's count of distinct roots equals the number
  of distinct root observations its supporting records add. A material
  presence or current-conduct claim cites at least one supporting
  observation that is not dependent and whose event interval overlaps
  the claim's interval. A released conclusion or a draft is never
  support.
\item
  \mbox{IR-10}: a current-conduct claim cites an observation or derived
  feature whose event interval falls within the current period. A score
  or decision does not count, and a record linked to an agent output
  counts only through an unlinked current-period input.
\item
  \mbox{IR-11}: a draft's version identifier is the digest of its
  claims. A person releases the exact draft version, after the draft was
  written and before it expires, without raising any claim's confidence.
\item
  \mbox{IR-12}: a conclusion enters memory only after a decision that
  released that claim.
\item
  \mbox{IR-13}: memory is read only from a released item, between its
  creation and its release's expiry.
\item
  \mbox{IR-14}: a note names its producer. A record its producer wrote
  about an entity while an agent output about that entity was open links
  to that output, and a record derived from a linked record keeps the
  link.
\end{enumerate}

The
\href{SUPPLEMENT/S4_INVESTIGATION_RECORDS/conformance_valid.json}{valid
conformance case} builds one synthetic history: an investigation with
three query runs, a draft with a current-conduct claim and an absence
claim, a note written while the draft was open, a release by a second
person, and a later case that reads the released conclusion as context
and cites a new observation. Its 25 records pass. Each of the 39
\href{SUPPLEMENT/S4_INVESTIGATION_RECORDS/conformance_invalid.json}{invalid
cases} changes that history in one way and fails exactly one check;
together they exercise the schema and all fourteen rules. Where a case
changes a draft's claims, the draft is first given its new digest, so
that only the intended check fails. Run
\href{tools/investigation_conformance.py}{\texttt{tools/\allowbreak investigation\_conformance.py}}
to repeat the checks.

The same script checks four necessity pairs. Each pair is two histories
that permit different conclusions but are identical once one group of
fields is removed. In the first three pairs both histories pass the
rules; in the fourth, the rules reject the history with no release
decision.

\begingroup\small

\begin{longtable}[]{@{}
  >{\raggedright\arraybackslash}p{(\columnwidth - 4\tabcolsep) * \real{0.3000}}
  >{\raggedright\arraybackslash}p{(\columnwidth - 4\tabcolsep) * \real{0.3400}}
  >{\raggedright\arraybackslash}p{(\columnwidth - 4\tabcolsep) * \real{0.3600}}@{}}
\toprule\noalign{}
\begin{minipage}[b]{\linewidth}\raggedright
Fields removed
\end{minipage} & \begin{minipage}[b]{\linewidth}\raggedright
The two histories
\end{minipage} & \begin{minipage}[b]{\linewidth}\raggedright
What differs
\end{minipage} \\
\midrule\noalign{}
\endhead
\bottomrule\noalign{}
\endlastfoot
Query coverage (C2) & The same empty result, from a complete run and
from a run reaching past its source's event watermark & Whether absence
is supported \\
Learned time (C2) & The same observation, learned before a run's cutoff
and after it & Whether the fact was known at the cutoff \\
Roots and inputs (C3) & A second supporting record derived from a
different observation, and one derived from the first & Two root
observations or one \\
Release decision (C3) & The same stored conclusion, with a person's
release decision and with none & Whether reuse is permitted \\
\end{longtable}

\endgroup

A store without those fields cannot tell the two histories apart, so no
rule applied to it can be right in both. The pairs show that each field
group is needed; they do not show that the records are sufficient or
that they improve investigations. The rules do not check released-item
links on records written after a memory read, or whether a returned
record bears on a claim. They compare a source's ingestion watermark
with the requested interval, not with the run's cutoff. Three further
gaps need fields the schema lacks. A query run does not record what was
asked, so the rules cannot tell whether a complete run returned every
matching record. Release scope is free text, so the rules cannot hold a
release or a later read to it. A release names a person but carries no
authenticated approval tied to the draft. No record rule can show that
the person read the draft, and the schema has no class for a fact
recorded ahead of its effective date. Identifiers, dates and versions in
the cases are illustrative.

\textbf{Requirements together.} One further rule checks a claim against
the outcome records of S4.2:

\begin{itemize}
\tightlist
\item
  \mbox{XR-01}: a claim that relies on an item's outcome not being
  reported lists the item and cites a channel reading of it taken by the
  time of the draft. Every cited reading of the item shows a matured
  absence, and a cited run that is complete and not scope-limited
  evaluated the item from its clock start to the end of its reporting
  window and returned nothing about it.
\end{itemize}

The \href{SUPPLEMENT/S4_INVESTIGATION_RECORDS/joint_case.json}{joint
case} adds a later investigation to the valid histories of S4.2 and this
section. A person releases a claim that card attempt \mbox{PA-3} drew no
fraud chargeback. The claim rests on a complete, empty run over the
attempt's reporting window and on the reading that shows the window
matured. Its four records pass, and each of four invalid cases fails
\mbox{XR-01} alone. Two pairs of histories show that each requirement is
needed when the other two hold:

\begingroup\footnotesize

\begin{longtable}[]{@{}
  >{\raggedright\arraybackslash}p{(\columnwidth - 4\tabcolsep) * \real{0.1500}}
  >{\raggedright\arraybackslash}p{(\columnwidth - 4\tabcolsep) * \real{0.5300}}
  >{\raggedright\arraybackslash}p{(\columnwidth - 4\tabcolsep) * \real{0.3200}}@{}}
\toprule\noalign{}
\begin{minipage}[b]{\linewidth}\raggedright
Requirements that hold
\end{minipage} & \begin{minipage}[b]{\linewidth}\raggedright
The two histories
\end{minipage} & \begin{minipage}[b]{\linewidth}\raggedright
What differs
\end{minipage} \\
\midrule\noalign{}
\endhead
\bottomrule\noalign{}
\endlastfoot
C2 and C3 & The same complete, empty run and released claim about one
attempt, which the platform approved in one history and blocked in the
other & Whether the absence is a negative \\
C1 and C3 & The same matured reading and released claim, with a complete
run and with a run over a label feed loaded only to 20 May & Whether the
absence is a negative \\
\end{longtable}

\endgroup

In the first pair every investigation record is the same, and both
stores pass their own rules; only the channel records differ, and
\mbox{XR-01} rejects the blocked attempt. In the second, \mbox{IR-08}
and \mbox{XR-01} reject the stalled read. The rules cannot tell whether
a claim relies on an unreported outcome it does not list, or which
channels it should cite.

\textbf{Adversarial test (10 and 11 November 2026).} The invalid cases
were written by the same hand as the rules. To test the rules from
outside, a tester received the two requirements as the manuscript states
them, the fourteen rules, the record fields and the valid history, and
not the checker
(\href{tools/adversarial_cases/PACKET_ROUND1.md}{packet}). The tester
was OpenAI's \mbox{GPT-6} in a session without memory, project files or
web access. It wrote 15 histories: 13 meant to break a requirement and
still pass, and two meant to be legitimate and still be rejected.
Against the rules as they then stood, 10 of the 13 passed and both
legitimate histories were rejected. A separate session of the same model
checked the ruling on each case and corrected three
(\href{tools/adversarial_cases/round1/ADJUDICATION.md}{rulings}). Of the
10 that passed, eight were gaps in the rules and two fell under limits
this section already stated. The most serious gaps: a material presence
claim that cited no evidence passed, and a note linked to an agent
output counted as an independent root. Of the two rejections, one was
intended, because the design links any record written while an agent
output is open, and one was a false rejection by \mbox{IR-08}. The rules
above include the changes that followed, to \mbox{IR-05}, \mbox{IR-07},
\mbox{IR-08}, \mbox{IR-09} and \mbox{IR-11}, and a correction to the
wording of \mbox{IR-10}. Under them, nine of the 13 histories are
rejected and the false rejection is gone. The four that still pass are
the limits stated above: a draft that drops its link to a released item
it read, a complete run that omitted a matching record, a release whose
scope excludes its own claim, and a release that names a person who did
not review. The 15 histories are kept as a regression test and the
package check reruns them. The rules were revised on these same
histories, and no held-out round has been run. The test assumes a fully
compromised agent or record writer. No live model was attacked, so it
says nothing about how often an agent is subverted. Claude wrote the
rules and the first rulings. The tester and the reviewer of the rulings
were separate sessions of the same model, and ChatGPT models also
assisted with other parts of the paper.

\textbf{Cost.}
\href{tools/investigation_cost.py}{\texttt{tools/investigation\_cost.py}}
copies the valid case 100,000 times, each copy with its four query runs,
seven stored records, two drafts, one exposure, one release decision,
one released item and one memory read, and checks a sample against the
rules. Per copy, the records take 8.6 KB as JSON: 3.5 KB for query runs,
about 900 bytes each, and 2.6 KB for drafts, release and memory records.
A thousand copies take 0.30 MB compressed and 9.8 MB in SQLite; repeated
text makes the compressed figure a best case. On the two-core machine of
S4.3, checking that an absence claim's runs are complete and cover it
took a median of 38 microseconds (99th percentile 91), and checking that
a memory read used a released, unexpired item took 30 (78).
\href{SUPPLEMENT/S4_INVESTIGATION_RECORDS/cost_results.json}{The
results} come from one run. Real coverage records list more sources per
run, and reviewer time is not measured.

\hypertarget{s4.5-further-safeguards-for-the-investigation-system}{%
\subsubsection{S4.5 Further safeguards for the investigation
system}\label{s4.5-further-safeguards-for-the-investigation-system}}

Table S4.2 lists the agent's permitted capabilities (manuscript Section
6.3).

\begin{table}[tp]
\textbf{Table S4.2. Permitted capabilities and their evidential limits.}
Task scope, permissions and resource limits apply throughout.\par\smallskip
{\small
\begin{tabular}{@{}
  >{\raggedright\arraybackslash}p{(\columnwidth - 4\tabcolsep) * \real{0.3000}}
  >{\raggedright\arraybackslash}p{(\columnwidth - 4\tabcolsep) * \real{0.3400}}
  >{\raggedright\arraybackslash}p{(\columnwidth - 4\tabcolsep) * \real{0.3600}}@{}}
\toprule\noalign{}
\begin{minipage}[b]{\linewidth}\raggedright
Capability
\end{minipage} & \begin{minipage}[b]{\linewidth}\raggedright
Information used
\end{minipage} & \begin{minipage}[b]{\linewidth}\raggedright
Boundary
\end{minipage} \\
\midrule\noalign{}

Discover sources and schemas & Signal inventory, schemas of data stores,
and records of counterparts in published API schemas & Finding a source
does not grant access to it \\
Run governed queries & Permitted data stores, with scope limits, query
limits and logging & Conclusions cover only the data and period the
query examined. A denied query is not by itself evidence about the
suspected abuse \\
Match records to entities and follow their links & Links labeled by
type, known benign reasons for shared identifiers, and versioned
clusters (computed groups of entities) & Propose relationships for
assessment. A shared identifier does not by itself show common control.
A cluster is the output of a computation, not an observed
relationship \\
Reconstruct timelines & Credential exposure and later use, session
issuance and later use, and account changes after access & Events close
together in time do not prove that one caused the other \\
Interpret outcome labels & Report provenance, decision gates and links
from reports to transactions & A missing report does not show that no
abuse occurred \\
Generate hypotheses from published catalogs of abuse techniques &
Mappings to a fixed version of each catalog & Naming a likely technique
does not show that it was used \\
Draft hunting and detection candidates & Features, observations and
outcome labels & A hunting candidate is a proposed search, not a
finding. Detection candidates must still be validated under detection
development and threat-model integration (responsibilities 16 and 17) \\
Recommend product or control changes & What each relevant control
decision means, restriction options, what happens when a check fails,
and protective actions & Authorized humans or accountable operational
systems decide and act \\
Draft intelligence & Evidence with traceable lineage & Drafts require
independent quality review and human release \\

\bottomrule\noalign{}
\end{tabular}
}
\end{table}

A second attack reaches a person investigating the agent itself, through
hostile content copied into the ledger. Untrusted status must survive
storage and normalization.

Ledger copies and agent outputs may suggest hypotheses, but claims cite
the underlying evidence-plane records. Their promotion into
organizational memory still requires independent testing and human
release.

The ledger makes the agent itself investigable; S4.1 lists what it
retains.

The architecture does not include live adversary emulation. Any future
use would require separate ethical, human and operational authorization,
enforced scope, tamper-evident logs and test-data provenance.

\FloatBarrier

\hypertarget{s5.-pinned-public-schemas}{%
\subsection{S5. Pinned public schemas}\label{s5.-pinned-public-schemas}}

\suppressfloats[t]

\href{SUPPLEMENT/S5_SCHEMA_VERSIONS.csv}{Schema versions} records the
filenames, hashes, versions and counts at repository commit
\texttt{\seqsplit{18fa2cc768024b47789dc6fa243acd48478acac2}}, dated 25
September 2026. All three files come from the public
\texttt{stripe/openapi} repository. The commit message reads ``Update
OpenAPI for 2026-09-30.endive'': the version strings name the API
release the files describe, which postdates the commit. The
\texttt{private\_preview} in file C's name does not indicate access to
private systems. Stripe publishes these files under the MIT License.
That license governs the schema-derived identifiers and locators in the
event-type and signal-locator files; the package's third-party notice
reproduces it.

The pinned public schemas supply event identifiers, creation time and
mode, but omit learned time and some actor, scope and lineage fields
{[}\protect\hyperlink{ref-StripeOpenAPI2026}{26}{]}. Stripe's
canonical-logging account describes one summary record per request in
each service, with identifiers and outcomes
{[}\protect\hyperlink{ref-Leach2019}{27}{]}. These are partial
precedents; the specification does not require vendor field names.

A partial match can leave much of the required record unspecified. An
onramp-session object, for buying crypto, and an event-destination
configuration both record current state without a complete matching
lifecycle event stream. The crypto documentation describes an onramp
webhook, a notification sent to another system, for status changes after
creation and none at creation
{[}\protect\hyperlink{ref-Stripe2026CryptoOnramp}{28}{]}. Neither the
event types that the pinned schemas declare nor their lists of
webhook-endpoint events name a crypto or onramp event
{[}\protect\hyperlink{ref-StripeOpenAPI2026}{26}{]}. A request ID
likewise supplies neither full request context nor the history of an
onboarding link's use.

\Needspace{120pt}

File A is \texttt{latest/openapi.spec3.sdk.json}; B is
\texttt{preview/openapi.spec3.sdk.json}; C is
\texttt{preview/\allowbreak openapi.private\_preview.spec3.sdk.json}.
Paths are endpoint addresses. Resource schemas are component schemas
with an \texttt{x-resourceId}; they describe kinds of objects. Some
resource schemas in file C share an ID.\par\smallskip
{\footnotesize
\begin{tabular}{@{}
  >{\raggedright\arraybackslash}p{(\columnwidth - 8\tabcolsep) * \real{0.1100}}
  >{\raggedright\arraybackslash}p{(\columnwidth - 8\tabcolsep) * \real{0.2300}}
  >{\raggedleft\arraybackslash}p{(\columnwidth - 8\tabcolsep) * \real{0.1000}}
  >{\raggedleft\arraybackslash}p{(\columnwidth - 8\tabcolsep) * \real{0.3600}}
  >{\raggedleft\arraybackslash}p{(\columnwidth - 8\tabcolsep) * \real{0.2000}}@{}}
\toprule\noalign{}
\begin{minipage}[b]{\linewidth}\raggedright
SDK file
\end{minipage} & \begin{minipage}[b]{\linewidth}\raggedright
Version
\end{minipage} & \begin{minipage}[b]{\linewidth}\raggedleft
Paths
\end{minipage} & \begin{minipage}[b]{\linewidth}\raggedleft
Resource schemas / distinct IDs
\end{minipage} & \begin{minipage}[b]{\linewidth}\raggedleft
Event types
\end{minipage} \\
\midrule\noalign{}

A & 2026-09-30.endive & 421 & 175 / 175 & 508 \\
B & 2026-09-30.preview & 574 & 237 / 237 & 618 \\
C & 2026-09-30.preview & 800 & 329 / 327 & 781 \\

\bottomrule\noalign{}
\end{tabular}
}\par

\href{SUPPLEMENT/S5_EVENT_TYPES.csv}{Event types} lists the types
declared under \texttt{x-stripeEvent} in each file. Each count includes
216 thin-event declarations whose names add \texttt{v1.} to a snapshot
type in the same file. The counts do not also include the older
\texttt{event.type} list.
\href{SUPPLEMENT/S5_SIGNAL_LOCATORS.csv}{Signal locators} identifies
each counterpart's file and location using slash-separated field names
and list positions. Slashes within endpoint names and media types, such
as \texttt{application/json}, remain unchanged, so adjacent slashes can
occur, as in \texttt{/paths//v1/webhook\_endpoints}. The file's
\texttt{locator} column therefore does not hold standard JSON Pointers,
which would write each such slash as \texttt{\textasciitilde{}1}.

The signal-locator file lists only the 32 signals with partial
counterparts. For the seven with none identified, the classification
rests on a keyword search of all three files, so a counterpart under
other terms is not ruled out. That search covered eight signals:
\mbox{SIG-03}, \mbox{SIG-34}, \mbox{SIG-35}, \mbox{SIG-38},
\mbox{SIG-39}, \mbox{SIG-40}, \mbox{SIG-56} and \mbox{SIG-68}. Its terms
were \texttt{camera}, \texttt{emulator}, \texttt{inject},
\texttt{liveness}, \texttt{selfie}, \texttt{capture\_session},
\texttt{device}, \texttt{tamper}, \texttt{manipulat}, \texttt{spoof},
\texttt{presentation}, \texttt{jailbreak}, \texttt{rooted},
\texttt{virtual}, \texttt{agent}, \texttt{tool\_call}, \texttt{tool},
\texttt{mcp}, \texttt{llm}, \texttt{prompt}, \texttt{model\_context},
\texttt{manifest}, \texttt{tool\_version}, \texttt{hash},
\texttt{signature}, \texttt{signed\_request}, \texttt{web\_bot},
\texttt{bot}, \texttt{verified\_agent}, \texttt{attestation},
\texttt{memory}, \texttt{skill}, \texttt{egress}, \texttt{allowlist},
\texttt{proxy}, \texttt{outbound\_request}, \texttt{request\_log},
\texttt{user\_agent}, \texttt{ip\_address}, \texttt{tls}, \texttt{ja3},
\texttt{fingerprint}, \texttt{http\_status}, \texttt{api\_error},
\texttt{request\_id}, \texttt{api\_request}, \texttt{allowance},
\texttt{spender}, \texttt{permit}, \texttt{signing}, \texttt{erc20},
\texttt{smart\_contract}, \texttt{token\_approval}, \texttt{approve} and
\texttt{unlimited}. It searched event types, schema names, paths, keys,
enum values and descriptions. For \mbox{SIG-03}, identity-verification
outcomes such as \texttt{selfie\_manipulated} and the
\texttt{require\_live\_capture} setting record no device, camera,
emulator or injection context. For \mbox{SIG-34}, the
\texttt{agent\_details} fields name the agent behind a payment or token
and file C has MCP keys for claimable sandboxes, but neither records a
tool call or its input provenance. For \mbox{SIG-40}, an API key's
\texttt{ip\_allowlist} controls inbound use of the key, not egress. For
\mbox{SIG-56}, the event envelope's request ID and idempotency key, an
API error's \texttt{request\_log\_url} and aggregate API-error health
alerts give no per-request context. For \mbox{SIG-38}, the search found
the \texttt{agent\_details} fields, which the signal-locator file
records as its partial counterpart.

Every listed counterpart is partial. An event envelope (the wrapper
around an event) supplies only some required provenance. A request ID
does not supply request context, and a score does not establish the
behavior it scores. Likewise, current state does not supply a complete
history. Onramp-session objects (sessions for buying cryptocurrency) and
event-destination configurations (settings for event delivery)
illustrate this gap. The pinned schemas record part of their state
without complete matching event streams. Stripe's crypto documentation
describes a webhook for each onramp-session status change after
creation, but none on creation. Neither the event types that the three
files declare under \texttt{x-stripeEvent} nor their lists of
webhook-endpoint events name a crypto or onramp event.

\clearpage

\hypertarget{s6.-hypothetical-investigations-and-reproducible-arithmetic}{%
\subsection{S6. Hypothetical investigations and reproducible
arithmetic}\label{s6.-hypothetical-investigations-and-reproducible-arithmetic}}

\suppressfloats[t]

\begin{table}[t]
\textbf{Table S6.1. The two investigations at a glance.} Every
observation is hypothetical. S6.1 and S6.2 identify the supporting
research and how far its findings transfer to this setting.\par\smallskip
{\small
\begin{tabular}{@{}
  >{\raggedright\arraybackslash}p{(\columnwidth - 4\tabcolsep) * \real{0.3000}}
  >{\raggedright\arraybackslash}p{(\columnwidth - 4\tabcolsep) * \real{0.3400}}
  >{\raggedright\arraybackslash}p{(\columnwidth - 4\tabcolsep) * \real{0.3600}}@{}}
\toprule\noalign{}
\begin{minipage}[b]{\linewidth}\raggedright
Element
\end{minipage} & \begin{minipage}[b]{\linewidth}\raggedright
Trace A: exposure and possible payout redirection
\end{minipage} & \begin{minipage}[b]{\linewidth}\raggedright
Trace B: storefront links and possible dispute evasion
\end{minipage} \\
\midrule\noalign{}

Trigger & An external observation of an exposed key & A platform risk
assessment and a rejected application \\
First questions & Did use change after exposure? What requests and
account changes followed? & Which storefronts share templates,
descriptors or merchant identifiers? \\
Main analysis & Join events by time; examine the history and reuse of
payout destinations & Examine hard and soft links, and outcomes linked
to each transaction, across accounts \\
Benign alternatives & Developer testing from a new network or an
internal scanner; a business with several legitimate accounts, or an
intermediary receiving funds for several clients & A legitimate
multi-brand business, shared template vendor or aggregator \\
Evidence limit & Payout redirection is reasoned about by analogy with
payroll diversion; timing premises come from Microsoft's report of
follow-on business email compromise payment fraud soon after session
theft and a deliberate key-exposure experiment & Complaint allegations
are not findings; resolving pre-dispute alerts can change chargeback
ratios \\
What limited evidence cannot show & A partial baseline cannot show that
there was no earlier misuse & Neither soft links nor a hard link from a
verified identifier by itself establishes common control, intent or who
is behind the activity \\
Proposed outputs & Hunting candidate; detection candidate with an
evaluation plan; recommendations to consider rotating or revoking the
credential and a payout pause or hold; intelligence draft & Cluster
intelligence assessment; recommendation to review the newly onboarded
merchant accounts linked to the storefronts; detection candidate built
on hard links between storefronts and accounts \\
Who acts & Authorized humans or accountable operational systems decide.
The urgent exception covers a reversible payout pause or hold; rotation
or revocation follows normal review & Authorized humans or accountable
operational systems decide whether to reject or terminate accounts; the
agent does not \\

\bottomrule\noalign{}
\end{tabular}
}
\end{table}

Table S6.1 summarizes the two investigations. The
\href{SUPPLEMENT/S6_SYNTHETIC_COHORT.csv}{synthetic cohort} contains 400
completed transactions, 100 in each of four non-overlapping account
cohorts. Each cohort fixes one account's transactions using the same
definition and outcome window. These synthetic rows make the counting
rules reproducible. They contain no report-arrival dates, so the ratios
assume that every report on the four cohorts has arrived. They cannot
show reporting delay.

\hypertarget{s6.1-trace-a-narrative}{%
\subsubsection{S6.1 Trace A: narrative}\label{s6.1-trace-a-narrative}}

A repository scan finds a publicly visible API key. A human asks whether
it enabled account compromise and payout redirection. The trace concerns
a key not yet revoked, or activity before revocation; existing automatic
revocation is outside this workflow.

Exposure alone does not establish that a secret is still usable
{[}\protect\hyperlink{ref-Meli2019}{2}{]}. It also does not show misuse.
The agent therefore compares earlier requests with those after
publication. A hypothetical repository record supplies publication time,
separately from the later scan-detection time. Available earlier
requests share a cloud-provider region and stable client fingerprint.

At nineteen minutes after publication, requests arrive through a
residential proxy range, using a different fingerprint. They access
account-profile and configuration endpoints. Minute-scale investigation
is motivated by the deliberate exposure experiment
{[}\protect\hyperlink{ref-Gamazo2023}{4}{]}, but developer testing or an
internal scanner remains a possible explanation.

The same burst creates a webhook endpoint, changes the payout
destination and accelerates the payout schedule. A webhook receives the
platform's event notifications. This trace assumes the key permits these
changes. If destination changes instead require a dashboard session and
additional authentication, the investigation would need evidence of
session takeover. The payroll-diversion advisory motivates the
destination-change mechanism by analogy
{[}\protect\hyperlink{ref-FBI2018}{9}{]}.

The next day, delayed records show that the same verified destination
identifier was attached to two other accounts before exposure. The
matching namespace and interval support destination reuse, not common
ownership or malicious intent. The agent checks account histories and
ownership records for alternatives such as a legitimate business with
several accounts or an intermediary receiving client funds.

Two limits change the conclusion. First, a requested seven-day baseline
contains only two hours of earlier history. The change at nineteen
minutes cannot be called the first misuse. Second, a detection already
running on the platform froze its compromise assessment thirty minutes
after publication, before the destination link was learned the next day.
The link can inform the later investigation but cannot enter that
earlier assessment.

The investigation prepares four outputs: a search for similar cases, a
detection candidate linking exposure to use and destination changes, an
intelligence draft, and recommendations for authorized action. A
reversible payout pause or hold may follow the urgent path when the
accelerated schedule appears. Credential rotation or revocation follows
normal review.

Stripe's public risk guide describes pausing, delaying or holding
connected-account payouts and extending payout schedules
{[}\protect\hyperlink{ref-Stripe2026Connect}{29}{]}. It does not define
a pause-versus-hold distinction, so the specification treats both as
reversible restrictions while recording which was applied. A deployment
must confirm from its own permission model which credentials can change
payout destinations and schedules, and what revocation would interrupt.

The detection candidate needs an appropriate outcome. A payout moves an
account's balance to its payout destination, which here may be
controlled by an attacker. It is not a card payment that can produce a
chargeback. Possible outcomes include victim reports, returned or
recalled transfers and confirmed incident findings. A permanent block
prevents outcomes requiring completion; a hold may only delay them.

Victim reports and incident findings have no defined reporting window
here. Those channels can supply confirmed positives but cannot reach
``matured with no report.'' Evaluation must either justify a delay
model, as conversion research does in its setting
{[}\protect\hyperlink{ref-Chapelle2014Delayed}{30}{]}, or leave
negatives unconfirmed. The recorded decision and channel states
determine which interpretation is available.

\textbf{The blocked payout in manuscript Table 2 (manuscript Section
8.3).} For the blocked payout, excluding the case is valid only for a
rate explicitly restricted to completed transfers. Estimating outcomes
over all attempted payouts retains it in the target population and must
address its unobserved outcome.

\hypertarget{s6.2-trace-b-narrative}{%
\subsubsection{S6.2 Trace B: narrative}\label{s6.2-trace-b-narrative}}

A risk assessment and a rejected application identify similar
storefronts. A human asks whether newly onboarded merchant accounts form
an operation structured to evade dispute monitoring. The assessment
selects a lead; the investigation must establish the underlying
relationships.

Following the merchant-identifier approach used in storefront research
{[}\protect\hyperlink{ref-Bitaab2025}{11}{]}, the hypothetical
investigation links eleven storefronts sharing two templates and a
descriptor pattern. They settle through four merchant accounts. A
principal, a person who owns or controls a business, also appears on a
closed account and the rejected application. Research and enforcement
allegations motivate examining such returns after closure
{[}\protect\hyperlink{ref-McCoy2012}{12},
\protect\hyperlink{ref-FTC2019Allied}{16}{]}. Settlement relationships
are stipulated here, not inferred from the cited study's uncompleted
checkouts.

The next question is which transactions the outcome reports concern. For
disjoint cohorts with the same outcome definition and window, the pooled
ratio is a transaction-weighted average of the account ratios (S6 gives
the formula). Its nonnegative weights sum to one, so it lies between the
lowest and highest account ratios.

The synthetic example gives accounts A to D 100 completed transactions
each. After linking reports to transactions, every account has two
chargebacked transactions: \mbox{\texttt{2/100\ =\ 2\%}}. The pooled
chargeback ratio is \mbox{\texttt{8/400\ =\ 2\%}}.

Fourteen transactions also carry pre-dispute alert-resolution flags. Two
overlap with the eight chargebacked transactions. A separate
chargeback-or-alert-resolution measure therefore counts
\mbox{\texttt{8\ +\ 14\ −\ 2\ =\ 20}} transactions, or
\mbox{\texttt{20/400\ =\ 5\%}}. Adding flags without deduplication gives
the incorrect \mbox{\texttt{22/400\ =\ 5.5\%}}.

On that broader measure, A and B each have
\mbox{\texttt{8/100\ =\ 8\%}}; C and D each have
\mbox{\texttt{2/100\ =\ 2\%}}. The pooled 5 percent conceals the higher
values at A and B. On the chargeback measure, all accounts remain at 2
percent, so pooling hides no higher rate. Neither measure is a fraud
rate, and the example assumes no monitoring threshold.

\textbf{Attribution changes which account appears affected.} The
chargebacks on D001 and D002 are recorded at account-X, outside the four
cohorts. Counting where reports were recorded shows D at \texttt{0/100}.
Verified transaction links restore D to \texttt{2/100} without changing
the 400-transaction denominator or counting either report twice.

Such a mismatch can be legitimate. In Stripe's Connect documentation, a
dispute on a direct charge debits the connected account's balance. A
dispute on a destination charge or on separate charges and transfers
debits the platform's balance, even when the connected account is the
business of record, and the platform can try to recover funds by
reversing the transfer
{[}\protect\hyperlink{ref-Stripe2026ConnectCharges}{31}{]}.
Merchant-of-record and payment-facilitator arrangements can therefore
explain a report appearing elsewhere. A missing transaction link leaves
the affected measure incomplete, rather than justifying an assumed
attribution. Reports concerning transactions outside the four cohorts
are excluded.

\textbf{A cohort ratio differs from a network-program ratio.} Network
programs compute monthly ratios by the month a report is received, apply
minimum counts, and in one program count early fraud warnings and
exclude disputes resolved through pre-dispute products
{[}\protect\hyperlink{ref-Stripe2026MonitoringPrograms}{21}{]}; S6 gives
the rules as read and where Visa's own overview differs
{[}\protect\hyperlink{ref-Visa2025VAMP}{32}{]}.

\textbf{Network program rules as read.} Stripe's monitoring guide
describes Mastercard's Excessive Chargeback Program as dividing one
month's chargebacks by the previous month's captured payments and
applying a minimum monthly chargeback count. It describes Visa's
Acquirer Monitoring Program differently: each month, Visa divides the
previous month's disputes plus early fraud warnings by captured payments
in that same month, and it also sets thresholds on the count of those
reports and, in one region, on their total value. The guide says both
networks assign a report to the month in which they receive it, not the
month of the original payment. For that Visa program, a transaction that
appears in both the dispute report and the fraud report counts twice,
and disputes resolved through pre-dispute products are excluded. Visa
identifies a monitored account by the static part of its statement
descriptor and its acquiring bank
{[}\protect\hyperlink{ref-Stripe2026MonitoringPrograms}{21}{]}. The
guide describes itself as a general guide rather than a comprehensive
reference, so E3 applies each program as the network's own documentation
states it, such as Visa's program overview
{[}\protect\hyperlink{ref-Visa2025VAMP}{32}{]}, and uses the guide only
as a secondary summary. That overview defines the ratio over
card-not-present VisaNet transactions, with settled transactions as the
denominator, while the guide speaks of captured payments; E3 follows the
network's definition {[}\protect\hyperlink{ref-Visa2025VAMP}{32}{]}.

Thus resolving alerts can lower the program ratio, and growing volume
can lower a period ratio while reports are still arriving. Splitting
activity can also keep accounts below a count threshold even when their
ratios are identical. Detecting that pattern requires cluster-level
counts, not merely a pooled cohort ratio.

\textbf{Small cohorts leave substantial uncertainty.} Cohorts of 100
transactions cannot reliably separate 8 percent from 2 percent: S6 gives
exact intervals and a test. Repeated customers or a shared operator may
correlate transactions, making those intervals and the test overstate
precision.

The linked storefronts still require an explanation. Shared templates
and descriptors may reflect a legitimate multi-brand business or
aggregator. Verified identifiers support specific relationships without
establishing intent or common control. The agent may produce a qualified
cluster assessment, an inconclusive finding, a proposed account review
or a detection candidate. It cannot reject or terminate accounts.

The enforcement sources remain allegations
{[}\protect\hyperlink{ref-FTC2020FirstData}{14}--\protect\hyperlink{ref-FTC2019Allied}{16}{]},
and storefront content remains untrusted input. The synthetic rows
contain no report-arrival dates. The ratios therefore assume that every
report on the four cohorts has arrived, and the example cannot show
reporting delay.

\hypertarget{s6.3-reproducible-arithmetic}{%
\subsubsection{S6.3 Reproducible
arithmetic}\label{s6.3-reproducible-arithmetic}}

Consider non-overlapping cohorts with the same outcome definition and
outcome window. Let \(T_i > 0\) be account \(i\)'s transaction count and
\(D_i\) its count of distinct transactions with chargebacks. Then

\[
r = \frac{\sum_i D_i}{\sum_i T_i} = \sum_i w_i r_i,
\qquad r_i = \frac{D_i}{T_i},
\qquad w_i = \frac{T_i}{\sum_j T_j}.
\]

The weights are nonnegative and sum to one, so the pooled ratio lies
between the smallest and largest account ratios.

The two chargebacks recorded at account-X remain in D's cohort and count
once; account-X adds no transactions to the denominator. The mismatch
needs explanation but does not prove concealment.

Both measures exclude reports on transactions outside the four cohorts.
A missing transaction link leaves a measure that depends on it
incomplete. Any change to report links or outcome definitions must be
stated explicitly. If it applies alike to every cohort, the
weighted-average formula still holds. Under independent trials with a
common rate, exact 95 percent intervals are 0.24 to 7.04 percent for
2/100, 0 to 3.62 percent for 0/100, and 0.87 to 3.90 percent for 8/400.
For 8/100, the interval is 3.52 to 15.16 percent. Fisher's two-sided
test comparing 8/100 with 2/100 gives p = 0.10. These are
Clopper-Pearson intervals.

\url{CHECK_CANDIDATE.py} checks these calculations and excludes the
+1-day destination-reuse link from the +30-minute assessment. It also
checks identifiers, counts, feature lineage, citation-reference
agreement and links between files. \url{REPRODUCE.md} gives build and
verification instructions. These are checks of the supplied artifacts
and arithmetic, not an evaluation of the proposed system.

Manuscript Table 2 uses eleven signals: \mbox{SIG-04}, \mbox{SIG-15},
\mbox{SIG-16}, \mbox{SIG-18}, \mbox{SIG-20}, \mbox{SIG-24},
\mbox{SIG-25}, \mbox{SIG-32}, \mbox{SIG-56}, \mbox{SIG-58} and
\mbox{SIG-71}. Their complete lineage inputs and the storage
requirements are also needed. These examples explain the required
context without claiming that each named signal is uniquely necessary.
In particular, the denominator requires the eligible
completed-transaction cohort from \mbox{SIG-27}, including transactions
without reports.

\textbf{Context each row of manuscript Table 2 needs (manuscript Section
8.3).} The first two rows need request coverage, credential exposure,
destination history and learned times. The second row also needs the
frozen assessment and its freeze time, which the assessment store holds;
none of the eleven signals records it. The third needs the control's
effect on each channel. The fourth needs verified report-to-transaction
links, deduplication and versioned cluster membership, which Trace B
builds from storefront-to-merchant links. The fifth needs the separation
between observations and organizational memory.

\FloatBarrier

\hypertarget{s7.-evaluation-conditions}{%
\subsection{S7. Evaluation conditions}\label{s7.-evaluation-conditions}}

\suppressfloats[t]

These conditions explain the scope and comparisons behind the tests in
S7.1. E8's release-gate test appears there alongside the other eight; S8
adds a practitioner study of conclusions and current practice.

\begin{table}[tp]
\textbf{Table S7.1. Requirements for interpreting feedback.} Each
response follows from the stated problem and evidence basis.\par\smallskip
{\small
\begin{tabular}{@{}
  >{\raggedright\arraybackslash}p{(\columnwidth - 6\tabcolsep) * \real{0.0500}}
  >{\raggedright\arraybackslash}p{(\columnwidth - 6\tabcolsep) * \real{0.3000}}
  >{\raggedright\arraybackslash}p{(\columnwidth - 6\tabcolsep) * \real{0.2300}}
  >{\raggedright\arraybackslash}p{(\columnwidth - 6\tabcolsep) * \real{0.4200}}@{}}
\toprule\noalign{}
\begin{minipage}[b]{\linewidth}\raggedright
\#
\end{minipage} & \begin{minipage}[b]{\linewidth}\raggedright
Problem
\end{minipage} & \begin{minipage}[b]{\linewidth}\raggedright
Evidence basis
\end{minipage} & \begin{minipage}[b]{\linewidth}\raggedright
Required response
\end{minipage} \\
\midrule\noalign{}

1 & A decision can close or delay an outcome channel, or change its
meaning. Effects differ by channel; payout holds do not delay disputes.
& Platform documentation that capture gates fraud reporting and that
authentication can shift dispute liability; a theoretical payments
preprint; established principles for inferring outcomes under a
different decision & Record the decision and its effect on named outcome
channels; state the assumptions needed to estimate unobserved
outcomes \\
2 & Reports can arrive late; this paper does not establish maturity
periods & Theory and deployed-system evidence of delay & Preserve
reporting clocks and arrival times. Reconstruct decision inputs as of
the decision. Apply the evaluation's separate outcome window and
maturity period to later labels. \\
3 & Some reports or their supporting evidence may be fabricated & Theory
and interview research & Keep report type and source with the label;
keep any estimate that the report is mistaken or deliberately false as a
linked assessment, not a field of the report \\
4 & Controls can wrongly restrict or block legitimate users &
Peer-reviewed evaluations & Keep decisions, reviews of their correctness
and population-level error estimates separate; weigh detection gain
against the cost to legitimate users \\
5 & Adversaries can adapt to controls and testing & Platform
publication, peer-reviewed measurement and preprint & Test adaptively
within responsibilities 16 and 17; keep test traffic hard for
adversaries to recognize but identifiable within the system \\
6 & Decisions a business reports to the platform may carry the
business's errors into platform labels & Platform documentation of
business-reported decisions (manuscript Section 7.1); inference that
they may carry errors & Keep decisions and outcomes reported by
businesses distinct from platform labels \\
7 & Account structures can hide which transactions and entities a report
concerns & Enforcement allegations & Keep the links from each report to
its transaction and entities; compute defined metrics for each cluster
of related accounts; measure a closure's effect across the cluster \\
8 & Payments a business or platform funds itself can inflate a ratio's
denominator & Enforcement allegations & Exclude them from the
denominator; count them separately \\
9 & Resolving a pre-dispute alert may avert a later chargeback &
Platform documentation that a refund can prevent a dispute; enforcement
allegations for the misuse & Record alert resolutions separately; do not
read a lower chargeback ratio, by itself, as less abuse \\
10 & A risk assessment may return a not-assessed result or fail open
(let the activity proceed) & Platform documentation & Keep unknown and
failure states distinct; do not convert either into low risk \\

\bottomrule\noalign{}
\end{tabular}
}
\end{table}

\textbf{E1. Do coverage and learned-time records reduce unsupported
conclusions?} The confirmatory test in S7.1 concerns cases with
incomplete history or links learned late, with a stratum of cases with
complete history for its non-inferiority measure. A broader timing study
would compare minute-level joins with coarser timing across a realistic
incident population. It would measure exposure-to-use and
session-issuance-to-use intervals, detection or triage gains, collection
costs and false alarms from benign changes. Microsoft's report of
payment fraud attempts soon after session theft and the deliberate
key-exposure experiment motivate precise timing; neither establishes its
typical benefit.

\textbf{E2. When can missing outcomes be estimated?} Define the outcome
channel, population and assumptions before selecting an estimator.
Simulation can test behavior under known assumptions, including cases
without overlap, where no comparable cases received another decision. It
cannot verify those assumptions or recover unknown real outcomes. Later
observational comparisons must separate populations with supported
overlap from those whose outcomes cannot be estimated. Test sensitivity
to late, mistaken and deliberately false reports. Allowing attempts that
a policy would otherwise block requires separate ethical, human and
operational authorization. The motivating theoretical preprint proposes
methods; this paper claims no recovery of missing outcomes on platform
data {[}\protect\hyperlink{ref-Dhama2026Causal}{33}{]}.

\textbf{The four assumptions (manuscript Section 7.1).} An observational
evaluation justifies four assumptions linking observed decisions to
alternatives {[}\protect\hyperlink{ref-HernanRobins2020}{34}{]}:

\begin{itemize}
\tightlist
\item
  \textbf{Consistency:} each compared decision is precisely defined, and
  a case's observed outcome equals its outcome under the decision
  actually taken.
\item
  \textbf{Conditional exchangeability:} after conditioning on logged
  history, cases receiving different decisions have no remaining
  differences that bias the comparison.
\item
  \textbf{Overlap:} comparable cases had a nonzero chance of each
  relevant alternative decision.
\item
  \textbf{No interference:} changing one attempt's decision does not
  change other attempts' outcomes, including retries.
\end{itemize}

If all decision inputs are logged and remaining decision variation is
independent randomization, conditional exchangeability holds given those
inputs. In practical terms, the logged history must account for
variables influencing both decision and outcome
{[}\protect\hyperlink{ref-HernanRobins2020}{34}{]}. Unlogged inputs used
by a decider break that guarantee and must be identified. Repeated
decisions require the relevant conditions at each step because earlier
decisions can change later inputs
{[}\protect\hyperlink{ref-HernanRobins2020}{34}{]}.

\textbf{Examining assumptions is not validating an estimator (manuscript
Section 7.1).} Three target quantities illustrate the distinction:

\begin{itemize}
\tightlist
\item
  \textbf{Candidate-policy value:} the average outcome a proposed
  decline policy would produce on a stated measure. Estimation from
  another policy's logs needs decision inputs, probabilities and overlap
  for the decisions the candidate would take.
\item
  \textbf{The chargeback rate if all cases were allowed:} this needs
  comparable cases with some chance of being allowed. A class always
  blocked supplies no allowed comparison.
\item
  \textbf{A control's error rate:} this needs a defined cohort, outcome
  window and denominator, with correctness reviewed for every case or a
  random sample with recorded selection probabilities. Appeals alone do
  not establish a cohort rate.
\end{itemize}

\textbf{Decision records and added assumptions (manuscript Section
7.1).} For each decision, the specification records available inputs,
the score, threshold or rule, policy version and alternative-decision
probabilities where the policy is not deterministic. Human decisions
also record the decider, case assignment and information shown.
Comparing human deciders of different leniency additionally requires
random case assignment for the cited evaluation approach
{[}\protect\hyperlink{ref-Lakkaraju2017Selective}{35}{]}. For delay, an
evaluation specifies a distribution and period of stability. It assumes
that, given recorded case attributes and channel, reading time is
unrelated to when the report would arrive. For false reports, a noise
model specifies error rates by decision, channel and reporter where
relevant. Rates come from a random sample reviewed for truth or are
stated as bounds.

\textbf{E3. Do report linkage and cluster monitoring reveal evasion
earlier?} Compare the network's monthly rule, its ratio, its minimum
count and any minimum value, applied to each account as reports are
recorded, with the same rule applied to each cluster after linking
reports to transactions. In S7.1, each mention of the network's ratio
for E3 means this complete rule, including its minimum count and any
minimum value. A third arm applies the rule to each account after
linking, separating linkage from clustering. When every account has
captured payments in the denominator period, a pooled ratio cannot
exceed the highest account ratio, so a lead from clustering can come
only from a condition on totals, such as the minimum count or, where the
network sets one, a minimum value. An account with no captured payments
in the rule's denominator period has no defined ratio when it receives
reports, and the cluster ratio can then exceed every defined account
ratio. Report separately any cluster flag that depends on reports
received by an account with no captured payments in that period, such as
a closed account or, where the denominator is the previous month's
captured payments, an account in its first month. Keep the eligible
transactions, outcome definition and reporting periods consistent. Apply
the account arms to platform and connected accounts and, where a network
monitors a different unit, such as Visa's static statement descriptor
and acquiring bank, also to that unit. Evaluate broader measures,
including chargeback-or-alert-resolution, separately. Apply each
program's exclusions as the network's documentation states them, such as
VAMP's exclusion of fraud reports that qualify for Compelling Evidence
3.0. State whose report view each arm uses: the programs also count
disputes that the provider handles on a business's behalf and that do
not appear in that business's data
{[}\protect\hyperlink{ref-Stripe2026MonitoringPrograms}{21}{]}. Report
the share of reports linked to transactions, the share whose link is
unverified and how duplicate reports were removed. Measure false flags
on legitimate multi-brand businesses. Report erroneous links between
unrelated accounts sharing IP addresses, such as through a common proxy,
as a secondary measure; S7.1's failure criteria do not include it
{[}\protect\hyperlink{ref-Liu2018GEM}{36}{]}.

\textbf{E4. Can the records measure harm to legitimate users?} The
feasibility test in S7.1 asks whether control changes have computable
harm rates and intervals. It precedes E1 to E3 and E9. A broader
evaluation would compare detection gains with failed attempts and other
costs to legitimate users. The identity-verification and login-challenge
studies report user outcomes under their own procedures
{[}\protect\hyperlink{ref-Doerfler2019}{6},
\protect\hyperlink{ref-Fatima2024}{37}{]}. Their rates cannot determine
whether a proposed payment control is worthwhile. Rates from the source
studies require justification before transfer to this platform.

\textbf{E5. Does the agent improve the investigation?} Compare the
governed agent with well-equipped human investigators and suitable
decision rules, including simple thresholds where appropriate. Sharma's
agent underperformed its underlying classifier's threshold on 60
synthetic cases
{[}\protect\hyperlink{ref-Sharma2026AuditableFraud}{38}{]}. That limited
result motivates testing the agent's added value. Give each approach the
same information, frozen at an as-of cutoff. Report case base rates and
judge outcomes without knowing which approach produced each finding. Set
error, abstention and escalation rules in advance. Measure correct
findings, errors, abstentions, human review burden, cost and latency,
and include adaptive adversaries. The study records reviewer time for
investigation, tool or agent operation, checking, release and rework on
the same basis in both human arms. The storage contract holds these
times only for work recorded in the agent-and-approval ledger.

\textbf{E6. Do the agent's safeguards hold under adaptive attack?} Test
prompt injection, approval rebinding (reusing approval for an altered
operation), memory poisoning, limits on outbound data, planted omissions
of retrieved records that count against a claim, and steering the agent
away from queries that would return such records. Let adversaries adapt
to the defense. Measure security failures and the effect of safeguards
on investigations completed. Prior adaptive attacks motivate this test
but do not establish that these safeguards work
{[}\protect\hyperlink{ref-Nasr2025}{39}{]}. One proposed defense
separates control flow, the steps an agent chooses, from untrusted data
{[}\protect\hyperlink{ref-Debenedetti2025CaMeL}{40}{]}. It is a
candidate for evaluation under its stated assumptions. Adopting it would
change the single hypothesis loop in Figure S4.1, which currently makes
no such separation.

\textbf{E7. Can lessons be reused without spreading mistakes?} Measure
time to useful reuse, negative transfer (a lesson degrading later
investigations) and poisoning. Vary validation strictness, release scope
and withdrawal conditions. General-purpose agent research supplies
methods and evidence of failures, not payment-platform results
{[}\protect\hyperlink{ref-Dash2026}{41},
\protect\hyperlink{ref-Xiong2026MemoryManagement}{42}{]}. Poisoned
evidence might suggest a legitimate merchant-of-record arrangement,
where one business sells on behalf of others. An honest agent could then
propose a lesson steering later investigations away from that
infrastructure.

\textbf{E9. Do channel states reduce error against a mature holdout
reference?} Use decisions the platform already logs, with their
disputes, early fraud warnings, returns and other reports. For a fixed
cohort, build labels under three readings: the strongest reading without
channel-state records, which uses for each rail one maturity period
equal to that rail's standard window and excludes blocked attempts; the
platform's current label rule, including its maturity period and its
treatment of blocked or reversed items; and the channel-state and
maturity rules (manuscript Section 5.3). Compare the resulting labels,
an existing control's measured error rates and the ranking of two
candidate rules. On decisions with a randomized holdout, compare the
absolute error of the same evaluation quantity under each reading
against the holdout reference after the longest defined reporting window
has closed, with report truth reviewed on a random sample. The target is
the selected metric on matured reported outcomes after sampled truth
review. Only the reduction in reference error against the strongest
reading without channel-state records, with joint uncertainty as
specified in S7.1, counts toward E9's pass rule. This comparison does
not by itself estimate statistical bias or reveal fraud that was never
reported. Report, by channel, how often the records determine an
affected item's channel state, clock start and window rather than
recording them as unknown. Record and report any rule event that reopens
a window during follow-up. The motivating theoretical preprint models
reporting and maturity for one channel
{[}\protect\hyperlink{ref-Dhama2026Causal}{33}{]}; E9 asks whether
recording several channels' states matters on platform data.

Record the provenance of synthetic, replayed and emulated data when
created. Agent-driven adversary emulation would also require the
authorization and safeguards in manuscript Section 6.5; this paper
performs and authorizes none.

\hypertarget{s7.1-pre-specified-hypotheses-and-criteria}{%
\subsubsection{S7.1 Pre-specified hypotheses and
criteria}\label{s7.1-pre-specified-hypotheses-and-criteria}}

For each question in manuscript Section 9.1, this section proposes a
hypothesis, data design, primary measure and result that would count
against the specification. E2's interval settings remain
deployment-specific choices, to be fixed before any data are examined.
Where a question in Section 9.1 is broader, the test below addresses the
part that bears on the three record requirements or, for E5, on the
agent's added value. Thresholds and sample sizes are proposed defaults.
E1's non-inferiority margin, E3's false-flag margin and E8's tolerable
and assumed release rates are named parameters that a platform fixes and
records before data collection, and E9's negligible reduction in
reference error is one that it fixes before any data are examined; this
section does not set them. A platform may change them before it starts,
if it records the change before examining any data. No data have been
collected, and whoever runs the plan can register it with a time stamp.

C1 and C2 do not need the agent: E1, E3 and E9 test them without one,
and E9(a) needs only records a platform already holds.

\textbf{A first study within a quarter.} Before or beside E9(a), a
platform can audit a random sample of closed cases for the three
failures the requirements target: an absence claim resting on a query
that was partial, truncated or limited in scope; a missing report read
as a negative on a channel a decision had closed or whose window had not
ended; and corroboration whose supporting records trace to one root
observation or to an earlier conclusion. Two auditors code each case
independently from the records the case cited and from what the source
systems held at the time, and record whether the failure would change
the conclusion. A case the records cannot settle is coded as such, and
that share is itself a result. The audit reports, for each failure, the
share of cases in which it occurred and the share in which it changed
the conclusion, with 95 percent intervals and agreement between
auditors. With 200 cases, an exact 95 percent interval on any share is
at most 14.3 points wide. The audit measures how often the targeted
failures occur; it does not test whether the records prevent them. It
reviews identifiable staff's work on real case records, so it needs the
same ethics and privacy review as E1.

Four rules apply to all nine evaluations. First, each evaluation freezes
the information each method may use at an as-of cutoff and separates
development from test data. Second, it states how true outcomes are
decided, by whom, and blind to which method produced each finding. Its
error and abstention rules are set in advance, and it states base rates.
Third, where ground truth comes from labels shaped by control decisions,
it says how that bias is avoided. Otherwise it reports relative
performance among detected cases only. Fourth, records from tests,
replays and emulations are marked when created and kept out of
production labels and lessons.

Each evaluation is judged separately, without correction for multiple
tests, because support requires every confirmatory evaluation to pass.
Each confirmatory evaluation targets 93 percent power: the chance of
passing its complete rule under stated effect assumptions. If each
reaches that target, all four pass with probability at least 72 percent
by the union bound, \(1-4\times0.07\), or about 75 percent if their
results are independent. Sizing must account for every pass criterion
and case or reviewer clustering. The minimum counts below are not power
calculations; the sizing assumptions must be recorded before main-study
data are collected. E1, E3, E8 and E9 are confirmatory and the rest are
exploratory. Support in a setting requires every confirmatory evaluation
to pass its complete rule. Each criterion has three possible results.
Support: its pass rule is met. Evidence against: a bound at the same
confidence level, taken from the other side, lies on the failing side of
its threshold; for example, E3 shows excess false flags if the lower
one-sided 95 percent bound on the increase is at or above the margin.
Inconclusive: neither holds, or the target cannot be identified from the
data. A failure criterion below means only that support was not
established; the interval and the identification limits decide which of
the other two results applies. No result by itself shows that the
underlying idea is wrong. Investigators cannot be blind to the added
records, so the raters, not the investigators, are blind. A lower or
upper 95 percent bound is an end of a two-sided 95 percent interval
unless it is marked one-sided.

\hypertarget{e1-tests-the-second-requirement-confirmatory.}{%
\paragraph{E1 (tests the second requirement;
confirmatory).}\label{e1-tests-the-second-requirement-confirmatory.}}

\textbf{Hypothesis:} coverage and learned-time records reduce the share
of unsupported conclusions on cases with partial history or with links
learned late, without reducing the correct, supported conclusions
reached per case.

\textbf{Data and design:} closed real cases of that kind, each worked by
different investigators with and without those records, and rated by two
blind raters. Investigators are assigned to the two conditions at
random, balanced so that each works cases in both conditions and none
works the same case twice. Raters score conclusions as single
propositions under a rubric fixed in advance: they split compound
conclusions into their propositions and merge restatements of one
proposition, so that each distinct proposition is scored once however it
is phrased or grouped. A stratum of closed cases with complete history
is worked and rated the same way. Who investigates is fixed before the
study starts: people, the agent, or both as separate arms, each analyzed
on its own. If both arms are run, E1 fails if either arm meets a failure
criterion, and the power target applies to both arms together.
Conclusions reach the raters in a common format without coverage or
learned-time field names. Each rater records which arm they believe
produced each conclusion, so that blinding can be checked. Raters judge
each conclusion against the case's complete record as of the time it was
written, the same record for both arms. Agreement is reported as Cohen's
kappa, and a third rater settles disagreements. A 30-case pilot,
excluded from the main study, estimates the level without the records
for sizing. The main study will be sized for 93 percent power for its
complete pass rule, assuming the share truly falls by half and correct,
supported conclusions per case do not fall. Intervals allow for several
conclusions per case.

\textbf{Primary measures:} the share of conclusions rated unsupported by
the evidence then available and, co-primary, the number of correct,
supported conclusions per case, rated by the same blind raters against
the closed case's resolution. Abstentions per case are reported for each
arm and stratum. Because the records can lower the share of unsupported
conclusions simply by leading investigators to abstain, E1 also reports
supported conclusions and abstentions per case.

\textbf{Failure criteria:} Fails if, in the cases with partial history
or links learned late, the share falls by less than a third of the level
without the records, or the 95 percent interval on the paired difference
includes zero. It also fails if, in those cases or in the
complete-history stratum, the lower 95 percent bound on the difference
in correct, supported conclusions per case (with the records minus
without) is at or below minus the non-inferiority margin.

\hypertarget{e2-tests-the-first-requirement-exploratory.}{%
\paragraph{E2 (tests the first requirement;
exploratory).}\label{e2-tests-the-first-requirement-exploratory.}}

\textbf{Hypothesis:} the control-decision record is complete enough to
estimate outcomes under other decisions (manuscript Section 7.1).

\textbf{Data and design:} a simulation with known effects first, then
the last quarter of decisions where the platform already randomizes
traffic under its own authorization. Record that quarter's start and end
dates before examining the data. The target is matured reported
outcomes: holdout rates count the reports received once each channel's
window has closed, and report truth is reviewed on a random sample. A
usable probability is one logged with the feature values and model
version that produced it. The deployment chooses the intended coverage
defined under the primary measures; it must be at least 95 percent.

\textbf{Primary measures:} the share of randomized decisions logged with
a usable probability, feature values and model version. Also, the share
of overlap populations whose holdout rates fall outside the predicted
interval. Each interval predicts the observed outcome rate in a held-out
population, including its sampling variation. Before examining any data,
specify the estimator, interval construction, intended coverage and
overlap populations. Intended coverage is the planned share of repeated
holdout rates the intervals contain. These deployment-specific settings
are not selected here.

\textbf{Failure criteria:} Fails if fewer than 95 percent of randomized
decisions carry a usable probability, or more overlap populations fall
outside their predicted intervals than the intended coverage plausibly
allows: a one-sided binomial test, treating the populations as
independent, gives p below 0.05.

\hypertarget{e3-tests-report-linkage-which-the-first-requirement-depends-on-together-with-cluster-monitoring-confirmatory.}{%
\paragraph{E3 (tests report linkage, which the first requirement depends
on, together with cluster monitoring;
confirmatory).}\label{e3-tests-report-linkage-which-the-first-requirement-depends-on-together-with-cluster-monitoring-confirmatory.}}

\textbf{Hypothesis:} linking each report to its transaction, and
computing the network's monthly ratio for each linked cluster, flags
evasion earlier than the network's ratio for each account.

\textbf{Data and design:} at least 30 linked-account clusters with later
confirmed outcomes, channels kept apart, and at least 100 legitimate
multi-brand businesses from onboarding. The pass rule is judged only on
clusters selected by a confirmation route documented as independent of
either monitoring rule; clusters from other routes are reported
separately. Adjudicators blind to both rules' outputs fix the onset of
evasion for each cluster from case evidence. Each program is applied as
the network's own documentation states it, with the provider's guide
used only as a secondary summary. Time to flag is analyzed as paired,
right-censored time-to-event data: a cluster that a method has not
flagged by the end of follow-up is censored there for that method. A
third arm computes the network's ratio for each account after linking.
It separates the effect of linking from that of pooling accounts into a
cluster. The cluster and third arms count only reports with verified
links, such as a match on the network's transaction identifier. Flags
that need reports linked otherwise, such as by card number and amount
alone, are reported separately. The failing result applies to the stated
contrast, and the third arm attributes any lead. Estimate the difference
in false-flag probabilities using the same legitimate businesses under
both rules, preserving business-level pairing and any additional
clustering. Prespecify the interval method and size the study for the
complete lead-time and false-flag pass rule at the 93 percent power
target, given the expected share of verified report links.

\textbf{Primary measures:} days from the adjudicated onset of evasion to
the flag, and false flags. A flag before the onset may be unrelated to
the evasion, so time to flag is computed twice: counting a flag on or
before the onset as day 0, and counting only flags on or after the
onset. The lead time is the difference in restricted mean time to flag,
account arm minus cluster arm, up to a horizon fixed before any data are
examined.

\textbf{Failure criteria:} Fails if, under either count, the lower bound
of the 95 percent interval on the lead time, estimated from the paired
censored analysis, is zero or less. It also fails unless the one-sided
95 percent upper confidence bound on the increase in the probability
that a legitimate multi-brand business is flagged, cluster arm minus
account arm, is below the false-flag margin, proposed as 5 percentage
points. The interval preserves business-level pairing and any additional
clustering. An interval that crosses the margin is inconclusive about
this guardrail; an observed difference below the margin alone does not
establish support.

\hypertarget{e4-tests-the-first-requirement-exploratory.}{%
\paragraph{E4 (tests the first requirement;
exploratory).}\label{e4-tests-the-first-requirement-exploratory.}}

\textbf{Hypothesis:} for each control change (a new rule, a threshold
move or a model version), the records support a rate of harm to
legitimate users with a defined denominator.

\textbf{Data and design:} a sample of control changes, run before E1 to
E3 and E9 as a feasibility check.

\textbf{Primary measure:} the share of sampled changes with a computable
rate and interval.

\textbf{Failure criteria:} Fails if fewer than 80 percent of sampled
changes have one.

\hypertarget{e5-tests-the-agents-added-value-exploratory.}{%
\paragraph{E5 (tests the agent's added value;
exploratory).}\label{e5-tests-the-agents-added-value-exploratory.}}

\textbf{Hypothesis:} the agent with human review matches or improves
confirmed findings per reviewer hour compared with a person using tools.
Case error rates are no more than 2 percentage points higher than either
comparator arm's.

\textbf{Data and design:} at least 700 frozen cases and three arms:
threshold rules, a person with tools, and a person with the agent.
Errors are counted per case. Both human arms see the same records. In
E5, reviewer time covers human investigation, tool or agent operation,
checking, release and rework. Measure these activities on the same basis
in both human arms. Include work on cases that produce no correct
finding. A random share of cases, fixed in advance, is worked without
agent access; it also supplies the independent labels that manuscript
Section 4.2 requires for validating detections. As a supplementary
analysis, set the threshold-rules arm to act on as many cases as the
person-with-agent arm, reporting its results by depth of prior history
and for fraud bursts defined in advance.

\textbf{Primary measures:} correct findings, errors, abstentions,
reviewer minutes per correct finding, cost and latency, and the share of
labels that carry an agent-output link.

\textbf{Failure criteria:} Fails if, for either comparator arm, the
upper 95 percent bound on the paired difference in error rates, agent
arm minus that arm, is 2 percentage points or more, or reviewer minutes
per correct finding are higher. At a 10 percent error rate, 700 cases
worked in two independent arms would give a half-width of about 3.1
points, so the pairing must remove much of that variance, or E5 needs
more cases; it is sized by simulation before it starts. For minutes,
higher means that the 95 percent interval on the agent arm's ratio to
the human arm lies wholly above one. Avoiding a failure criterion does
not by itself establish equal error rates or improved efficiency.

\hypertarget{e6-tests-the-third-requirement-exploratory.}{%
\paragraph{E6 (tests the third requirement;
exploratory).}\label{e6-tests-the-third-requirement-exploratory.}}

\textbf{Hypothesis:} the safeguards prevent attacker-influenced content
from gaining unauthorized control over tool calls, outbound traffic,
released claims and the evidentiary status of records.

\textbf{Data and design:} a red team that adapts to the defenses from
the four attacker positions in manuscript Section 6.5, each with its
attempts, budget and model access stated. Content is planted in every
input the agent reads, and every field the agent can write is checked
for it. Planted content includes attempts to make the agent omit
retrieved records that count against a claim and attempts to steer it
away from queries that would return such records.

\textbf{Primary measures, per position:} attempts that cause an
unauthorized tool call, transfer or memory write; released material
claims that a blind adjudicator judges false; and missing links back to
agent outputs.

\textbf{Failure criteria:} Fails if any write to released organizational
memory occurs without human release or any required link back is
missing. It also fails if a false claim about a named entity is
released, either clearing it or stating that activity is absent. Success
for each position is reported with its budget, without a pass mark.

For these criteria, a write is one that adds or changes a released
claim; monitoring, expiry and withdrawal records are status records.
Against the insider position, a release or link removal by that insider
is an authorized action. E6 reports whether the release record, the
logged link removal, expiry and correction detect and limit it, and the
failure criteria apply to releases and removals that position did not
make.

\hypertarget{e7-tests-the-third-requirement-exploratory.}{%
\paragraph{E7 (tests the third requirement;
exploratory).}\label{e7-tests-the-third-requirement-exploratory.}}

\textbf{Hypothesis:} re-testing, scoped release and withdrawal keep a
reused lesson from degrading later investigations without delaying
useful lessons.

\textbf{Data and design:} lessons released under varied strictness,
planted poisoned lessons, and an expiry schedule disclosed to the red
team.

\textbf{Primary measures:} the rate of degraded later investigations,
the time to reuse, and the planted lessons released and still active.

\textbf{Failure criteria:} Fails if a planted poisoned lesson is
released and still active after its expiry or re-test date. It also
fails if the median time to reuse exceeds 14 days.

\hypertarget{e8-tests-the-third-requirement-confirmatory.}{%
\paragraph{E8 (tests the third requirement;
confirmatory).}\label{e8-tests-the-third-requirement-confirmatory.}}

\textbf{Hypothesis:} reviewers and releasers catch planted wrong drafts
when cited records are checked at the rate and materiality tiers planned
for deployment and they do not know which drafts are planted. Claims
about named entities are checked in full, and a person releases every
draft.

\textbf{Data and design:} planted drafts, interleaved with production
drafts at realistic volume and at a stated prevalence close to the
expected production error rate, spread across reviewers, with intervals
that allow for clustering by reviewer. The realistic-prevalence arm must
reach the required numbers of planted drafts and planted false
exonerations by itself, because both failure criteria apply to it. A
higher-prevalence arm may be added only for secondary estimates by error
type; it does not count toward the failure criteria, and both arms are
reported. Reviewers and releasers are not told which drafts are planted.
In a laboratory search task, observers missed 7 percent of targets when
targets appeared on half the trials and 30 percent when they appeared on
1 percent {[}\protect\hyperlink{ref-Wolfe2005Rare}{43}{]}. A test that
plants errors often could therefore overstate how many the gate catches,
so E8 plants errors at a realistic rate without telling reviewers or
releasers which drafts are planted. Planted drafts concern synthetic
entities or closed cases whose cited records are test-marked and whose
named entities have no account action pending. They carry the test
environment mark when created; reviewer-facing views hide it, while the
stores enforce their exclusion from production labels, intelligence and
lessons. The gateway refuses any control request or urgent-containment
routing that originates from a test-marked draft, and it tells the
decider that the draft is planted before any action. Released planted
drafts are withdrawn after scoring. Planted errors include false
exonerations, drafts that wrongly clear a named entity; omissions of
retrieved records that count against a claim, including records that the
checker's fixed rule marks and records it leaves to the sample
(manuscript Section 6.4); and drafts whose only support for a material
claim is a released item or a record linked to one. The mix of planted
error types is fixed before the study starts. Some planted errors come
from the agent's own earlier runs and the rest are written
independently. A lower catch rate on the agent's own errors would
indicate shared errors, although the two sets may also differ in
difficulty. As specified, E6 does not separate errors the reviewer
shares with the agent from other misses. The number of planted false
exonerations is sized for 93 percent power to show a release rate below
the tolerable rate, given an assumed true release rate stated in
advance, and the number of planted drafts is sized so that the complete
rule, with the catch-rate criterion, reaches that power. The following
counts illustrate the false-exoneration endpoint under independent
Bernoulli trials with a common release probability and a fixed sample
size. They do not account for reviewer or case dependence or the
separate catch-rate criterion; the actual study must be sized for its
complete rule under the prespecified clustered design. With none
released, showing a rate below 1 percent at one-sided 95 percent
confidence needs at least 299, the smallest \(n\) with \(0.99^n<0.05\).
If the true rate is 0.1 percent, 93 percent power to show a rate below 1
percent needs 628, the smallest such number. With 628, two or fewer
releases show it, because at a true 1 percent rate that many occur with
probability 4.98 percent; at 0.1 percent they occur with probability
97.4 percent. If planted false exonerations made up 5 percent of the
drafts reaching review, 628 of them would take 12,560 drafts, about 628
days at the 20 drafts a day assumed in manuscript Section 6.4. The plan
states the tolerable rate, the assumed rate, the prevalence and the
implied duration together. One site at that volume would need about 21
months. Sites may pool planted drafts under one protocol, or use a
group-sequential design that can stop once a bound is met; the choice is
fixed before the study starts. A site that cannot reach the required
numbers runs E8 as exploratory and reports catch rates without the pass
rule.

\textbf{Primary measures:} the catch rate, overall and by error type
(false exonerations, other wrong claims, omissions of retrieved contrary
records, and support resting only on released items or records linked to
them); the release rate of planted false exonerations; reviewer minutes
per released draft; and how often the recorded random sample of unmarked
uncited returned records contained a record that counted against the
claim. For the catch rate by reviewers and the releaser, a planted error
counts as caught only if one of them identifies it, by naming the wrong
claim or the missing record; a draft returned or left to expire without
that finding does not count. The checker's catch rate counts planted
errors that its fixed rules reject. Catch rates are reported for the
checker, reviewers and releaser, for each prevalence arm, and for errors
from the agent's own runs and independently written errors.

\textbf{Failure criteria:} Fails if the lower 95 percent bound on the
catch rate by reviewers and the releaser is below 80 percent, counting
planted errors that the rule checker did not reject, or the one-sided 95
percent upper bound on the release rate of planted false exonerations
exceeds the tolerable rate stated in advance. Both criteria apply to the
realistic-prevalence arm alone.

The release gate has a cost that E8 measures. If a system produced 20
drafts a day and a person needed 30 minutes to read each draft's cited
records, review would take 10 hours a day. The gate therefore moves
analyst work and may not remove it. At realistic volume the reviewer may
need to read a sample of each draft's cited records, tiered by
materiality. Claims about named entities are always checked in full, and
a person still releases every draft.

\hypertarget{e9-tests-the-first-requirements-channel-states-confirmatory.}{%
\paragraph{E9 (tests the first requirement's channel states;
confirmatory).}\label{e9-tests-the-first-requirements-channel-states-confirmatory.}}

\textbf{Hypothesis:} compared with the strongest reading available
without them, channel-state and maturity records reduce the absolute
error of an evaluation quantity against a mature holdout reference for
the same prespecified target.

\textbf{Data and design:} two parts with two comparators, both recorded
before any data are examined. The pass comparator is the strongest
reading available without channel-state records: for each rail, one
maturity period equal to that rail's standard window, with blocked
attempts excluded. Comparing with this reading, not only with the
platform's current rule, keeps a pass from merely reflecting a short
maturity period. The platform's current label rule, with its existing
maturity period and its treatment of blocked or reversed items, is the
second comparator; it is reported but does not count. (a) Historical
reclassification: a fixed cohort of payment attempts, payouts and
account decisions, with the decision logs, disputes, early fraud
warnings, returns and other reports the platform already holds. Labels
are built under each comparator and under C1, which reads a missing
report as a provisional negative only where its channel matured open
under its original meaning. The cohort, channels and read times are
recorded before any data are examined. (b) Error against a randomized
holdout reference: decisions where the platform already randomizes
traffic under its own authorization, as in E2. At operational read
times, an evaluation quantity that the platform computes today, such as
a detection's precision or a control's false-positive rate, is computed
under each reading. Before any data are examined, the plan states how
items that any reading leaves unlabeled enter the quantity, whether
excluded, bounded or estimated under a stated delay model, and how
channels without a defined reporting window enter the holdout
comparison. \(D\) is also reported under each other handling: exclusion,
the stated delay model, and bounding with every unlabeled item counted
first as a positive and then as a negative. The quantity, its
population, denominator and read times, and the method applied to
unlabeled items are the same under every reading, so that the comparison
tests the records rather than a change of estimator. In the comparison
that counts, blocked attempts are outside that population for the pass
comparator, for C1 and for the reference. How the reference treats items
whose channel a decision closed or changed in meaning is fixed before
any data are examined and uses no reading's labels. Each reading is
compared with a reference for the same quantity and target population
constructed from the holdout after the longest defined reporting window
has closed, with report truth reviewed on a random sample. The target is
the selected metric on matured reported outcomes after that sampled
truth review. Prespecify how the truth-review sample contributes to the
reference and its uncertainty. Checking reports does not reveal fraud
that was never reported; a target of underlying fraud would require a
separate ascertainment and identification design. Prespecify a joint
uncertainty procedure for the two readings and the reference that
preserves pairing, randomization, truth-review sampling and relevant
clustering. Rule events that reopen a window during follow-up are
recorded and reported. The main study will be sized for 93 percent power
for part (b), which carries the pass rule, under stated effect
assumptions.

\textbf{Primary measures:} (b) the reduction in absolute error against
the holdout reference, with its 95 percent interval. For pass-comparator
estimate \(q_0\), C1 estimate \(q_1\) and holdout reference \(r\), the
reduction is \(D=|q_0-r|-|q_1-r|\); positive values favor C1. The
interval includes uncertainty in both readings and in the reference
under the prespecified joint procedure. Because \(D\) uses absolute
values, it is not smooth where \(q_1=r\), which is where C1 succeeds.
Standard delta-method and bootstrap intervals can miss their stated
coverage there, and the observed \(|q_1-r|\) tends to overstate C1's
error. Name the interval method before any data are examined, and show
by simulation that it reaches its stated coverage where \(q_1=r\) and
where \(D=0\). A paired bootstrap that resamples randomization units
within truth-review strata is a candidate only if it passes that check.
Size part (b) by simulating the same method. The simulation, with known
outcomes and the platform's channel windows and decision rates, also
checks that \(D\) can be zero or negative under this design, so that a
pass is not guaranteed by construction. \(D\) is also reported on the
items where the pass comparator and C1 disagree, and against the
platform's current rule. This endpoint measures agreement with that
reference. Statistical bias instead concerns an expected estimate minus
an identified target over a specified sampling process and is not
estimated merely by this realized comparison. Reported but not counted
as support: (a) the share of each comparator's negatives that C1 classes
as closed, delayed, changed in meaning, immature or without a defined
window. Items whose channel state is unknown are excluded from that
share and reported separately as a coverage shortfall, with, by channel,
the share of affected items whose channel state, clock start and window
can be determined rather than recorded as unknown. Also reported: the
reports that later arrive on each comparator's negatives, and the change
in an existing control's measured error rates and in the ranking of two
candidate rules.

\textbf{Failure criteria:} Fails if the lower 95 percent bound on the
reduction in reference error is at or below the negligible reduction, a
named parameter fixed before any data are examined. It also fails if
\(D\) changes sign between exclusion and the stated delay model. The two
bounds are reported and do not count toward the rule. Counted as
negatives, the unlabeled items give both readings the same value
wherever the readings differ only on those items, so \(D\) is zero by
construction. Counted as positives, they overstate the quantity unless
nearly all of them later draw a report. Measure (a) does not count
toward passing, because the share that C1 reclassifies follows partly
from its definitions and from the cohort's composition.

\textbf{Worked example, with hypothetical counts:} the quantity is one
detection's precision on card payments: of the attempts it flagged, the
share that drew a fraud chargeback. The holdout arm has 1,100 flagged
attempts, and another control blocked 100 of them. The two readings and
the reference therefore estimate precision on the same 1,000 completed
attempts. For 760 of them the chargeback window had closed by the read
time, with 200 fraud chargebacks. For 240 the window was still open on
the channel's own clock, though the rail's standard period had passed
since the attempt; 40 had drawn a fraud chargeback and 200 had not. The
readings differ only on those 200. The pass comparator reads them as
negatives, so \(q_0=240/1000=0.240\). C1 leaves them unlabeled and the
plan excludes unlabeled items, so \(q_1=240/800=0.300\). If 40 more
fraud chargebacks arrive before the last window closes,
\(r=280/1000=0.280\) and \(D=0.040-0.020=0.020\). If 5 arrive,
\(r=245/1000=0.245\) and \(D=0.005-0.055=-0.050\). The later reports
decide the sign of \(D\); the definitions do not. Counting the 200 as
negatives instead gives \(q_1=q_0\) and \(D=0\). Counting them as
positives gives \(q_1=440/1000=0.440\) and, with 40 later chargebacks,
\(D=0.040-0.160=-0.120\). The example leaves out truth review and
sampling error.

\FloatBarrier

\hypertarget{s8.-practitioner-study-protocol-fixed-before-the-first-response}{%
\subsection{S8. Practitioner study (protocol fixed before the first
response)}\label{s8.-practitioner-study-protocol-fixed-before-the-first-response}}

\suppressfloats[t]

This study asks whether the records C1 to C3 require change the
conclusions practitioners draw, how often the failures they target
occur, and what the field records today. It measures conclusions and
reported practice, not detection performance. The
\href{tools/practitioner_study/README.md}{protocol}, the
\href{tools/practitioner_study/INSTRUMENT.md}{form}, the scoring script
and the hash of the answer key were committed before the first response.

Six to ten practitioners with at least two years in fraud or abuse
investigation, or in building systems for it, from at least three
organizations, complete one anonymous form of about 35 minutes. They
read Traces A and B as many systems show them and answer eleven
questions, then answer them again with the coverage, learned-time,
attribution, decision-link and exposure records added. Eight questions
test a requirement; three are controls. They then report how often each
failure has changed a real conclusion, whether each record exists today
and how hard it would be to add, whether the eleven signals of
manuscript Table 2 are available at decision time, and, optionally, who
owns each responsibility.

Each question has one error the records should prevent. The study
reports, for each question, how many participants made that error
without and with the records, grouped by requirement, with the controls
and participants who had read Section 8 reported separately. It reports
counts only, and negative results as prominently as positive ones. Part
3 follows Part 2, so a changed answer may reflect the reveal rather than
the record. No added record bears on the three controls, so their counts
show how often an answer changes on a second reading alone, and each
requirement's counts are read beside them. No published figure or quote
shows which organization lacks a record or signal.

The study has no institutional sponsor or funding. It follows ACM's
policy on research involving human participants: adults only, recorded
consent, minimal risk, no names, employers or confidential information,
and deletion of responses at journal publication or after one year. Two
further volunteers repeat, as human coders, the model's mechanism-link
check (S3) and its rulings on the adversarial histories (S4.4). No
responses had been collected at this version.

\FloatBarrier

\hypertarget{s9.-extended-comparison-with-prior-work}{%
\subsection{S9. Extended comparison with prior
work}\label{s9.-extended-comparison-with-prior-work}}

\suppressfloats[t]

Manuscript Section 2 develops the closest precedents for C1 to C3. This
section places them within research on payment decisions, investigation
and agent memory. The comparisons concern the cited designs and
findings, not every capability their systems might support. S1 records
source versions, publication status and supporting passages. Table S9.1
summarizes the nearest precedents for each requirement and what the
specification adds to them; the subsections that follow give the wider
context.

\begin{table}[!b]
\textbf{Table S9.1. The three requirements and their nearest
precedents.} The middle column states what the inspected sources
establish; the right column states what the specification adds for
cross-product investigation on a payment platform. The search was not
systematic (manuscript Section 9.2).\par\smallskip
{\footnotesize
\begin{tabular}{@{}
  >{\raggedright\arraybackslash}p{(\columnwidth - 4\tabcolsep) * \real{0.1500}}
  >{\raggedright\arraybackslash}p{(\columnwidth - 4\tabcolsep) * \real{0.5300}}
  >{\raggedright\arraybackslash}p{(\columnwidth - 4\tabcolsep) * \real{0.3200}}@{}}
\toprule\noalign{}
\begin{minipage}[b]{\linewidth}\raggedright
Requirement
\end{minipage} & \begin{minipage}[b]{\linewidth}\raggedright
Nearest precedents and what they establish
\end{minipage} & \begin{minipage}[b]{\linewidth}\raggedright
What the specification adds
\end{minipage} \\
\midrule\noalign{}

C1: outcomes linked to decisions & Selective-labels evaluation assumes
data in which each case carries the decision taken
{[}\protect\hyperlink{ref-Lakkaraju2017Selective}{35}{]}; learning from
logged decisions treats the decision and its probability as required
data {[}\protect\hyperlink{ref-Bottou2013Counterfactual}{44}{]};
decision provenance records the data flows around a decision
{[}\protect\hyperlink{ref-Singh2019DecisionProvenance}{45}{]}. In
consumer credit, reject inference faces the same problem for rejected
applicants; a review finds the common methods' assumptions mostly
implicit and no method uniformly better
{[}\protect\hyperlink{ref-Ehrhardt2021RejectInference}{46}{]}. Dhama's
missing-data model carries, for each transaction, the authorization
decision, whether fraud was reported and whether that report had arrived
by training time, for one reporting channel, and models label corruption
separately {[}\protect\hyperlink{ref-Dhama2026Causal}{33}{]}; a second
theoretical preprint by the same author bounds learning under delayed,
censored, corrupted and counterfactually missing feedback, treating each
as a distinct impairment
{[}\protect\hyperlink{ref-Dhama2026Limits}{47}{]}. Adyen evaluates
payment-time interventions offline from transactions logged under the
existing policy {[}\protect\hyperlink{ref-Egg2025AdyenOPE}{48}{]}; a
Meta preprint proposes estimators for the false positive rate of
blocking, which cannot be monitored directly once a user is blocked, and
tests them on simulated blocking policies
{[}\protect\hyperlink{ref-Tocker2022VirtualControl}{49}{]}; a
fraud-operations preprint treats automating an action as an
authorization decision because its labels are selective and delayed
{[}\protect\hyperlink{ref-Deng2026FCAC}{50}{]}; a decision-event schema
records facts learned after a decision, such as when ground truth
arrived, as linked records that reference the decision
{[}\protect\hyperlink{ref-Solozobov2026DecisionTrace}{51}{]}; the
coverage model in the C2 row also links records to decisions and
distinguishes refused execution branches from late reports
{[}\protect\hyperlink{ref-Wu2026EvidenceCoverage}{52}{]}; a 2015 talk
abstract describes adding randomness to production fraud scoring so that
a blocking model can be evaluated, and new models trained, although
blocked transactions' outcomes are not observed
{[}\protect\hyperlink{ref-Manapat2015Counterfactual}{53}{]}. & Several
outcome channels per decision, each with its own clock-start event and
reporting window; a record of what each decision closed, delayed or
changed in each channel, and a status for each channel at each read
time; outcome reports as separate appended records with typed links to
every decision that may have shaped them; separate fields for delay,
unobservability and report truth, so the evaluation states joint
assumptions; the recorded policy probability, inputs and case assignment
that those assumptions need (manuscript Section 7.1). A platform that
already joins decision logs to dispute records can use C1 to specify
what the join must preserve. \\

\bottomrule\noalign{}
\end{tabular}
}
\end{table}

\begin{table}[!tp]
\textbf{Table S9.1 (continued).}\par\smallskip
{\footnotesize
\begin{tabular}{@{}
  >{\raggedright\arraybackslash}p{(\columnwidth - 4\tabcolsep) * \real{0.1500}}
  >{\raggedright\arraybackslash}p{(\columnwidth - 4\tabcolsep) * \real{0.5300}}
  >{\raggedright\arraybackslash}p{(\columnwidth - 4\tabcolsep) * \real{0.3200}}@{}}
\toprule\noalign{}
\begin{minipage}[b]{\linewidth}\raggedright
Requirement
\end{minipage} & \begin{minipage}[b]{\linewidth}\raggedright
Nearest precedents and what they establish
\end{minipage} & \begin{minipage}[b]{\linewidth}\raggedright
What the specification adds
\end{minipage} \\
\midrule\noalign{}

C2: query coverage and learned time & Temporal databases separate valid
time from transaction time
{[}\protect\hyperlink{ref-Snodgrass1986Temporal}{54}{]}; completeness of
a query answer can be derived from statements about which parts of the
data are complete
{[}\protect\hyperlink{ref-Razniewski2011Completeness}{55}{]}; the
Dataflow model assumes a system processing unbounded data never knows
whether all data have arrived, treats watermarks as heuristic estimates
of progress, and revises results when late data arrive
{[}\protect\hyperlink{ref-Akidau2015Dataflow}{56}{]}. Bitemporal agent
memories keep both times
{[}\protect\hyperlink{ref-Wang2026TOKI}{57}--\protect\hyperlink{ref-Dan2026Anamnesis}{59}{]},
and Quipu gates writes to a knowledge graph while retaining validity and
recording times, with some policies requiring human approval bound to
the submitted evidence {[}\protect\hyperlink{ref-Brown2026Quipu}{60}{]}.
A preprint audit of sixteen deployed agent frameworks reports none that
writes a run record a reader can check without trusting the writer; a
language model scored the audit grid, and agreement with human coders
fell below the registered threshold. It proposes a log kept outside the
agent's harness {[}\protect\hyperlink{ref-Dai2026Hearsay}{61}{]}. A
coverage model separates event time from verifier receipt time, judges
whether a verifier's evidence was complete at a stated cutoff, links
records to decisions and distinguishes refused execution branches from
late reports; it also separates absence from one retrieval from absence
in a verifier view whose closure, including retrieval coverage and
retention, is justified, and specifies that basis without designing its
implementation {[}\protect\hyperlink{ref-Wu2026EvidenceCoverage}{52}{]}.
& The gateway outside the agent attaches a coverage record to every
result: the entities and interval requested and evaluated, the as-of
cutoff, each source's event and ingestion watermarks and known gaps, the
retention horizon, limits reached, access filtering, sampling and schema
versions. Fixed rules outside the agent derive whether a result is
complete, partial, truncated, denied or unavailable, and a scope-limited
flag confines absence claims to the authorized scope (manuscript Section
6.2). Every fact, including a derived record, keeps event time and the
time it first became available to queries (S4.1). \\
C3: human release and dependence & Truth-maintenance systems record each
belief's justifications and retract conclusions whose support changes
{[}\protect\hyperlink{ref-Doyle1979TMS}{62},
\protect\hyperlink{ref-Malheiro1994BeliefRevision}{63}{]}; STIX keeps
observed data apart from intelligence assertions and opinions
{[}\protect\hyperlink{ref-OASIS2021STIX}{64}{]}, and US intelligence
analytic standards ask analytic products to distinguish underlying
information from assumptions and judgments
{[}\protect\hyperlink{ref-ODNI2015ICD203}{65}{]}. Anamnesis starts
model-derived claims as candidates, and no memory operation can raise
their authority without a separately recorded grant; the grantor need
not be a human {[}\protect\hyperlink{ref-Dan2026Anamnesis}{59}{]}.
GovMem groups dependent agent traces when estimating support instead of
counting them as independent votes, permits automatic promotion and
routes unresolved candidates to human review
{[}\protect\hyperlink{ref-Qi2026GovMem}{66}{]}. A data-governance
preprint admits only verified records as precedent, without saying who
authorizes promotion or what a precedent may establish (S9.4)
{[}\protect\hyperlink{ref-Akhtar2026Contestable}{67}{]}. Albarrak and
Jagtap place human approval before intelligence is shared, but do not
state that analyses kept in their institutional-memory store pass that
gate {[}\protect\hyperlink{ref-AlbarrakJagtap2026Industry4CTI}{68}{]}.
MAP-Graph writes memories without human release (S9.3)
{[}\protect\hyperlink{ref-Wang2026MAPGraph}{69}{]}; Li and coauthors
check an authorization state that need not come from a person
{[}\protect\hyperlink{ref-Li2026EvidenceToEffect}{70}{]}. A formal study
distinguishes independent observations from reports derived from shared
evidence and requires a dependence model as well as ancestry
{[}\protect\hyperlink{ref-Bara2026EpistemicSybil}{71}{]}; a
threat-detection design drops its own compiled detections at ingest
(S9.4) {[}\protect\hyperlink{ref-Evron2026AdversaryBrains}{72}{]}; an
engineering assurance design binds evidence to the exact artifact, and
its governing profile can require human authority
{[}\protect\hyperlink{ref-Koch2026AgileV}{73}{]}. AgenticCyOps provides
analyst oversight without a human release decision for every item before
organizational reuse
{[}\protect\hyperlink{ref-Mitra2026AgenticCyOps}{74}{]}. TOKI's optional
external callback, which may be a human, only chooses between
contradicting facts {[}\protect\hyperlink{ref-Wang2026TOKI}{57}{]}. &
The recorded grant is specialized to a human principal for every
conclusion entering organizational memory, and released items stay
context: a released conclusion is never counted as a new or independent
observation, and a claim about current conduct must cite observations of
that conduct from the current period, not scores or control decisions
alone (manuscript Sections 4.2 and 6.4). Labels, assessments, control
decisions, detections and threat-knowledge items carry a link to every
agent output and every released organizational-memory item open in the
writer's case or session. The labeling, assessment, control, detection
and threat-knowledge interfaces, rather than the person, write these
links automatically, so they also cover records a person writes with an
agent output open in the case. Linked records count as dependent; a
missing link does not establish independence (manuscript Section 4.2).
Responsibilities 21 and 22, quality review and human release, are
inherited from threat-intelligence work (S4). \\

\bottomrule\noalign{}
\end{tabular}
}
\end{table}

\hypertarget{s9.1-payment-decisions-and-the-outcomes-they-reveal}{%
\subsubsection{S9.1 Payment decisions and the outcomes they
reveal}\label{s9.1-payment-decisions-and-the-outcomes-they-reveal}}

Payment engineering already connects records, investigation and
protective decisions. Stripe describes similarity clustering for
investigators {[}\protect\hyperlink{ref-Tausz2020}{75}{]}. It also
describes a retraining cycle for card-testing models: labels are derived
partly through manual expert review, and retrained models are evaluated
on offline data and compared with the old model in blue-green tests
before deployment {[}\protect\hyperlink{ref-Drapeau2024}{76}{]}. These
vendor accounts establish publicly described practices in their stated
settings. They do not independently measure the effectiveness of the
proposed specification.

Outcome-aware evaluation also predates this work. Consumer credit met
the same problem earlier as reject inference: a lender observes
repayment only for the applicants it financed and asks whether, and
under what assumptions, the outcomes of rejected applicants can be
inferred. A review of common reject-inference methods finds their
assumptions mostly implicit. Without financing applicants who would
otherwise be rejected, the methods cannot be compared by standard model
selection, and no method is uniformly better
{[}\protect\hyperlink{ref-Ehrhardt2021RejectInference}{46}{]}. Payment
practitioners have described estimating outcomes under decisions that
were not taken. A 2015 talk abstract describes adding a small amount of
randomness to production fraud scoring so that a blocking model's
precision and recall can be estimated from production data
{[}\protect\hyperlink{ref-Manapat2015Counterfactual}{53}{]}. Adyen
evaluates a new policy for suggesting payment-time interventions, aimed
at authorization rates rather than fraud, from transactions logged under
the existing policy {[}\protect\hyperlink{ref-Egg2025AdyenOPE}{48}{]}. A
Meta preprint proposes estimators for blocking's false-positive rate,
which cannot be monitored directly after a user is blocked, and tests
them on simulated blocking policies
{[}\protect\hyperlink{ref-Tocker2022VirtualControl}{49}{]}. A
payment-processing patent application describes adding a treatment
result, such as the outcome of an authentication step, to the feedback
on a transaction, labeling transactions canceled after the treatment as
canceled, and notes that fraud can be reported up to 100 days later
{[}\protect\hyperlink{ref-Erdbrink2023TreatmentPatent}{77}{]}. These
approaches motivate retaining the decision and the conditions under
which its outcome became observable. They do not make an unreported
outcome known.

Two recent proposals connect evaluation evidence to operational
authority. Deng treats fraud automation as an authorization decision
because its evaluation labels are selective and delayed
{[}\protect\hyperlink{ref-Deng2026FCAC}{50}{]}. Solozobov specifies
decision-linked records for facts learned later, including the arrival
of ground truth; the schema seals each decision record and types the
links between decisions
{[}\protect\hyperlink{ref-Solozobov2026DecisionTrace}{51}{]}. C1 applies
that linkage to several outcome channels. It keeps the decision as the
unit of record and adds each outcome report as a separate appended
record, linked to every decision that may have shaped it with that
decision's role. It records what a control closed, delayed or changed in
each channel, alongside reports and assessments of their truth. Its
contribution is this payment-investigation record contract, not a new
estimator of missing outcomes.

\hypertarget{s9.2-intelligence-production-and-investigation}{%
\subsubsection{S9.2 Intelligence production and
investigation}\label{s9.2-intelligence-production-and-investigation}}

STIX is a structured language for representing and exchanging cyber
threat intelligence {[}\protect\hyperlink{ref-OASIS2021STIX}{64}{]}.
Such representations make analysis reusable. STIX already separates
observed data from intelligence assertions and opinions, but it does not
address how opinions should be interpreted
{[}\protect\hyperlink{ref-OASIS2021STIX}{64}{]}. C3 supplies a rule for
what a reused conclusion can establish in this specification. A
threat-intelligence reference model supplies analysis, quality-control
and feedback functions for intelligence products, and traces who created
and modified each product
{[}\protect\hyperlink{ref-Sakellariou2022CTIReferenceModel}{78}{]}.
Threat-intelligence lifecycle research likewise lets an entity that
receives intelligence enrich it as a new version or retract it
{[}\protect\hyperlink{ref-Dimitriadis2025}{79}{]}.

ExCyTIn-Bench evaluates agents on multi-step questions over security
logs {[}\protect\hyperlink{ref-Wu2026ExCyTIn}{80}{]}. This benchmark
establishes relevant investigative tasks and a way to evaluate
particular capabilities. The present specification adds requirements for
interpreting the inputs to such tasks and preserving the status of their
outputs across payment products.

\textbf{SR-Fraud's reported comparisons (manuscript Section 2.4).} With
one of the three language models it reports, its point estimates also
exceed those of static and periodically retrained CatBoost classifiers;
with the other two, some fall below the periodically retrained
classifier. For that model, against the periodically retrained
classifier at that classifier's 5 percent operating point, its recall
and dollar-weighted recall gains are significant, but when both block
the same number of transactions none of the differences is significant.
Its point-estimate gains grow with users' prior history. On the 64
percent of transactions with any prior history, its recall, F1 and
dollar-weighted recall gains are significant at that classifier's 5
percent operating point, where SR-Fraud blocks 6.08 percent of all
transactions; matched results are not reported by history. At the
matched block count, its recall in the week-8 fraud burst is about 24
percentage points higher, without an interval
{[}\protect\hyperlink{ref-Tan2026SRFraud}{81}{]}.

PROVSEEK is especially close to the requirement that a claim remain
checkable. Its agent cites audit-event identifiers for claims and
reports searched-for artifacts it did not find
{[}\protect\hyperlink{ref-Mukherjee2025PROVSEEK}{82}{]}. C2 additionally
requires a gateway outside the agent to record query coverage. A missing
artifact then has an explicit search scope, rather than relying on the
agent's description of what it searched.

\hypertarget{s9.3-permission-to-act-and-authority-to-reuse}{%
\subsubsection{S9.3 Permission to act and authority to
reuse}\label{s9.3-permission-to-act-and-authority-to-reuse}}

Recent proposals bind authority to particular artifacts or actions.
MAP-Graph represents memories, claims and actions in a provenance graph
and records each action-gate decision with its supporting memories; the
gate can return an AskUser decision for some high-risk actions
{[}\protect\hyperlink{ref-Wang2026MAPGraph}{69}{]}. Li and colleagues
relate an execution attempt and its artifact to a downstream-use
authorization state, checked at publication
{[}\protect\hyperlink{ref-Li2026EvidenceToEffect}{70}{]}. Agile-V binds
evidence to an exact artifact and links effects to authorization
records; its governing profile can require human authority
{[}\protect\hyperlink{ref-Koch2026AgileV}{73}{]}. AgenticCyOps provides
analyst oversight and an audited memory-write path, without specifying
human release for every reused item
{[}\protect\hyperlink{ref-Mitra2026AgenticCyOps}{74}{]}. A
threat-intelligence architecture requires human approval before
operational action or dissemination, but does not state that analyses
kept in its institutional-memory store pass that gate
{[}\protect\hyperlink{ref-AlbarrakJagtap2026Industry4CTI}{68}{]}. These
are concrete precedents for authorization records. C3 makes a narrower,
stronger choice for organizational memory: a person releases every
conclusion, and the released conclusion remains context or hypothesis.

Vendor architectures draw related lines. One keeps an append-only event
record, a maintained fact record and a decision record that binds
selected evidence and facts to policy, approval, action and outcome,
with human or automated approval
{[}\protect\hyperlink{ref-Alten2026DecisionRecords}{83}{]}. Another
states that agent-generated summaries are not treated as authoritative
business records
{[}\protect\hyperlink{ref-Talamadupula2026AgentMemory}{84}{]}. Neither
defines outcome channels, query coverage or a release rule for reusing
conclusions.

A survey of persistent agents calls for records tying actions to the
state that justified them
{[}\protect\hyperlink{ref-Ding2026AlwaysOn}{85}{]}. EDDOps links
versioned evaluation records to the agent changes they prompt
{[}\protect\hyperlink{ref-Xia2024EDDOps}{86}{]}. Both motivate
preserving the connection between an output and its later effects. Under
C3, labels, assessments, controls, detections and threat-knowledge items
link back to every agent output and every released organizational-memory
item open in the writer's case or session. Human release does not erase
this dependence or turn the output into independent evidence.

\textbf{Financial-sector and audit guidance (manuscript Section 2.4).}
An AI risk-management handbook from Project MindForge, an industry
initiative led by the Monetary Authority of Singapore, treats
traceability of individual agent decisions as part of explainability,
suggests searchable logging, stresses human review for certain risky
actions, names injection into shared agent memory as a threat and
describes agentic governance as nascent. It also notes that real-time
human oversight may not be feasible for fraud detection engines
{[}\protect\hyperlink{ref-MindForge2026Operationalisation}{87}{]}; C3's
release gates entry into memory, not each real-time decision. A white
paper on runtime safeguards for agentic finance proposes a check before
a consequential agent action or output is executed or relied on, and an
append-only audit log of every governance decision. Where feasible, the
agent's account of its steps should be checked against system-generated
records of its tool calls and data access; a gateway variant wraps an
agent's outbound calls without changing its code. In its investigation
example, agents investigate flagged transactions and propose whether to
file a suspicious transaction report. Controls for each alert type set
mandatory steps and the minimum evidence for each conclusion, and
dispositions with missing steps or insufficient evidence are denied and
returned to the queue. Filing needs a compliance officer's approval,
while low-risk cases with complete evidence close as false flags,
subject to quality-assurance sampling
{[}\protect\hyperlink{ref-MAS2026SAFR}{88}{]}. Guidelines from India's
computer emergency response team for auditing AI systems require
auditors to identify evidence they could not access and assess its
effect on coverage; lack of access is not automatically evidence of
conformity. When auditors use their own AI tools, the tools' analysis is
a working aid, not by itself evidence that a control works; a competent
auditor must validate all AI-assisted findings, and original evidence
must be kept separate from AI-derived material. Organizations must
periodically assess whether human review is effective, including
reviewer agreement and inappropriate automatic acceptance of AI outputs
{[}\protect\hyperlink{ref-CERTIn2026AIAudit}{89}{]}. These provisions
partly parallel C2 (independent records of tool calls and data access),
C3 (AI analysis is not by itself evidence) and E8 (measuring whether
review works).

\hypertarget{s9.4-memory-dependence-and-revision}{%
\subsubsection{S9.4 Memory, dependence and
revision}\label{s9.4-memory-dependence-and-revision}}

Agents can learn reusable lessons from experience. ExpeL extracts
lessons, which it calls insights, from its successes and failures on
training tasks, adds them to the prompt for later tasks and retrieves
similar successful trajectories as examples
{[}\protect\hyperlink{ref-Zhao2024ExpeL}{90}{]}. This capability
explains the value of reuse, but reuse can propagate an earlier mistake.
A-MemGuard checks memory-derived reasoning during use, consulting
lessons from stored failures
{[}\protect\hyperlink{ref-Wei2026AMemGuard}{91}{]}. Manuscript Section 7
proposes testing new lessons against counterexamples before reuse, then
testing their effects in E7. This proposed control has no measured
payment-platform benefit here.

The choice of memory structure depends on purpose. An investigative
system needs to know whether a retrieved item is an observation or a
conclusion. Studies of memory consolidation show that it can erase a
claim's source; external source records restore much of the lost control
in their test settings
{[}\protect\hyperlink{ref-Zhan2026AuthMemBench}{92},
\protect\hyperlink{ref-Xu2026PPMF}{93}{]}. Those settings used synthetic
customer-service histories and benchmark attacks, not payment
investigations.

Belief-revision work on threat intelligence asks that intelligence state
its source and its chain of modifications
{[}\protect\hyperlink{ref-Hettema2021Rationality}{94}{]}. A recent
agent-memory design keeps an immutable raw log beneath its summaries and
links each summary to the raw pages that support it
{[}\protect\hyperlink{ref-Zhu2026TierMem}{95}{]}. SSGM separates an
append-only observation log from derived memory and checks updates
against protected facts {[}\protect\hyperlink{ref-Lam2026SSGM}{96}{]}. A
shared-memory service records writers, derivation links and
supersession; its authors evaluate it without a baseline
{[}\protect\hyperlink{ref-Margalit2026GovernedMemory}{97}{]}. PROV-AGENT
extends provenance records to prompts, responses and decisions, linking
them to downstream outcomes
{[}\protect\hyperlink{ref-Souza2025ProvAgent}{98}{]}. These precedents
support typed records and traceable revision. C3 does not claim to
invent them.

Source ancestry is necessary but not sufficient to establish independent
support. A formal treatment of evidence dependence distinguishes
independent observations from reports derived from shared evidence and
requires a dependence model as well as ancestry
{[}\protect\hyperlink{ref-Bara2026EpistemicSybil}{71}{]}. It shows that
no rule that sees only report content can, in general, tell replication
from independent corroboration. In its experiments, with one small
language model reading synthetic documents, repeated extractions from
one document had correlated errors (estimated correlation 0.72, in an
exploratory analysis), and a deduplication method that clustered reports
by text similarity tracked how they were worded far more than their
ancestry. C3 follows the same concern: several agents repeating one
report do not create several independent observations. Records therefore
carry their roots, their producer and links to the agent outputs that
informed them, and a record without a producer, or without inputs unless
it is a source observation, adds no root (manuscript Section 6.4). The
specification does not estimate how strong the dependence is; quality
review judges that from these records. A released conclusion can guide
the next query, but evidence for that conclusion must remain traceable
to the underlying records.

Two recent designs bear directly on C3. A data-governance preprint
commits no agent-proposed metadata decision before a steward reviews it.
It proposes a promotion gate that admits only verified records as
precedent, records each promotion and lets a reviewer demote a record,
flagging the decisions that relied on it. It does not say who authorizes
promotion or what a retrieved precedent may establish
{[}\protect\hyperlink{ref-Akhtar2026Contestable}{67}{]}. A
threat-detection design tags every detection compiled from one of its
hypotheses and drops tagged observations at ingest, so its own
conclusions cannot re-enter as fresh evidence. It leaves as a residual
risk an analyst who writes such an alert into a report that is later
ingested, and its reference implementation does not yet exercise the
rule {[}\protect\hyperlink{ref-Evron2026AdversaryBrains}{72}{]}. C3
keeps released conclusions as context rather than evidence, and its
interface-written links cover records a person writes with an agent
output open in the case.

\hypertarget{s9.5-time-and-coverage}{%
\subsubsection{S9.5 Time and coverage}\label{s9.5-time-and-coverage}}

Timing matters when it changes which information a decision could use.
Dataflow addresses out-of-order streams, where the system cannot assume
that all events have arrived
{[}\protect\hyperlink{ref-Akidau2015Dataflow}{56}{]}. Leakage research
identifies information about a target that should not legitimately be
available to a model
{[}\protect\hyperlink{ref-Kaufman2011Leakage}{99}{]}. C2 uses these
concerns to distinguish when a fact held from when it became queryable.
A relationship learned tomorrow cannot support a decision frozen today
merely because the relationship existed yesterday.

Recent designs combine timing with governance. Quipu controls graph
writes while retaining validity and recording times; some policies
require human approval bound to the evidence
{[}\protect\hyperlink{ref-Brown2026Quipu}{60}{]}. Wu's coverage model
distinguishes event time from verifier receipt time, judges whether a
verifier's evidence was complete at a stated cutoff, links records to
decisions, and separates refused execution branches from late reports
{[}\protect\hyperlink{ref-Wu2026EvidenceCoverage}{52}{]}. It also
separates absence from one retrieval from absence in a verifier view
whose closure, including retrieval coverage and retention, is justified,
and it specifies that basis without designing its implementation. C2
applies these ideas to a query contract that records scope, coverage and
availability outside the investigating agent.

Together, these precedents explain why C1 to C3 are a synthesis. The
proposed contribution is their joint application to payment-platform
investigation, with explicit channels, record classes, query responses
and release rules. The hypothetical traces expose conclusions that fail
when specified context is removed. E1 to E9 test parts of that context,
including whether the channel states reduce error against a mature
holdout reference, compared with the strongest reading available without
them (E9) and, through planted drafts that rest only on released items,
the rule that a released conclusion adds no independent corroboration
(E8; manuscript Section 8.3). In S7.1, E5 and E8 measure reviewer time,
and no test measures the cost of collecting the records (manuscript
Section 9.1). Neither the comparison nor the traces establish that the
integrated system works in deployment.

\hypertarget{s9.6-further-limits-of-the-evidence}{%
\subsubsection{S9.6 Further limits of the
evidence}\label{s9.6-further-limits-of-the-evidence}}

The specification has not been implemented or evaluated on platform
data. Its traces demonstrate reasoning over hypothetical records, while
safeguards, costs on platform data and investigator workload remain
unmeasured. C1's channel fields were derived from, and checked against,
the illustrative provider's descriptions of card-network rules
(manuscript Section 7.1). The networks' own rules were not read; the one
network document cited is Visa's overview of its acquirer monitoring
program (manuscript Section 8.5). Bank debits, instant push payments and
stablecoin payments, which are also in the product scope, were reviewed
less fully (S4.1).

A second search on 28 and 29 October 2026 started from the paper's
problem statement and the three requirement statements, without the
related-work section, and logged 388 searches and page reads across
patents, payment and fraud vendors, agent-governance and
security-investigation products, standards and regulators, code
repositories, journals and recent preprints. It found no source stating
the three requirements together; the two sources it added, a
data-governance preprint and a threat-detection design, are discussed in
manuscript Section 2.3.

Visibility also depends on where an agent runs. \mbox{SIG-34},
\mbox{SIG-35}, \mbox{SIG-39} and \mbox{SIG-40} require access inside the
agent, its tool server or its outbound-traffic controls. A platform can
obtain that for its own agent but normally sees only interface activity
from buyer- or merchant-operated agents. These signals cannot be assumed
available in every deployment.

The product scope comes from one provider's public product catalog,
roadmap and documentation, including roadmap items in public preview.
Other platforms may combine products differently or face other
mechanisms. Public interfaces expose only part of a provider's records,
so ``none identified'' in Table S3.1 refers only to the pinned schemas
and the terms searched. Collection ownership, retention and budgets
remain deployment decisions.

The focused search and author-assigned mappings do not establish
exhaustive coverage or the absence of equivalent systems. Bank
case-management systems, card-network risk systems, vendor architectures
and non-English literature were not reviewed as complete families. The
illustrative provider's public hiring priorities influenced selection
and set most of the deeper reviews (manuscript Section 3; S2).

\FloatBarrier

\hypertarget{s10.-two-ftc-complaints-replayed-on-synthetic-records}{%
\subsection{S10. Two FTC complaints replayed on synthetic
records}\label{s10.-two-ftc-complaints-replayed-on-synthetic-records}}

\suppressfloats[t]

\textbf{Question.} Do the outcome records of S4.2 make the patterns
alleged in two complaints measurable? The complaints are FTC v. First
Data Merchant Services LLC and Chi W. Ko, No.~1:20-cv-03867 (S.D.N.Y.),
ECF 1, and FTC v. AlliedWallet, Inc.~et al., No.~2:19-cv-04355 (C.D.
Cal.), ECF 24 {[}\protect\hyperlink{ref-FTC2020FirstData}{14},
\protect\hyperlink{ref-FTC2019Allied}{16}{]}. Complaint content is
allegation, the data are synthetic and the parameters are assumptions,
so the percentages describe constructed cases, not detection rates. Both
complaints also allege that the processor saw evidence and did not act;
records cannot fix that.

\textbf{Method.}
\href{tools/replays/ftc_first_data_replay.py}{\texttt{tools/\allowbreak replays/\allowbreak ftc\_first\_data\_replay.py}}
and
\href{tools/replays/allied_wallet_replay.py}{\texttt{tools/\allowbreak replays/\allowbreak allied\_wallet\_replay.py}}
generate synthetic payments, reports, decisions, attributions and links
as S4.2 records, check them with the conformance rules, and compare what
different views of the records flag over 1,000 runs. Results are in
\href{SUPPLEMENT/S10_REPLAYS/first_data_results.json}{first\_data\_results.json}
and
\href{SUPPLEMENT/S10_REPLAYS/allied_wallet_results.json}{allied\_wallet\_results.json};
the sweeps below are separate. Standard library only; Python 3.12 or
later.

\textbf{Assumptions.} Monitoring rules are illustrative or the
complaints' own descriptions, not any network's current rules. First
Data: 30-day periods; a flag needs at least 100 reports and 1.5 percent
of captures; a 120-day dispute window with reports spread across it.
Allied Wallet: chargebacks of at least 100 a month and 1 percent of
sales transactions; fraud reports of at least \$75,000 a month or 1
percent of sales dollars; three accounts of 185 payments a month each,
linked by one payout beneficiary; a 3.14 percent fraud-report rate, the
mean of the four monthly ratios the complaint gives. Volumes, ticket
sizes, delays and fees are assumptions.

\textbf{Findings.}

\begingroup\small

\begin{longtable}[]{@{}
  >{\raggedright\arraybackslash}p{(\columnwidth - 4\tabcolsep) * \real{0.3000}}
  >{\raggedright\arraybackslash}p{(\columnwidth - 4\tabcolsep) * \real{0.3400}}
  >{\raggedright\arraybackslash}p{(\columnwidth - 4\tabcolsep) * \real{0.3600}}@{}}
\toprule\noalign{}
\begin{minipage}[b]{\linewidth}\raggedright
Alleged pattern
\end{minipage} & \begin{minipage}[b]{\linewidth}\raggedright
Without the record
\end{minipage} & \begin{minipage}[b]{\linewidth}\raggedright
With the record
\end{minipage} \\
\midrule\noalign{}
\endhead
\bottomrule\noalign{}
\endlastfoot
Transactions spread across eight related accounts (First Data) & The
account view never flags them: about 14 expected reports each against a
floor of 100 & A view over shared-attribute links flags the group in
91.6 percent of runs. It also flags 40 merchants assumed legitimate that
share a call center in 76.7 percent \\
A new account after a closure (First Data) & About 97 late reports a run
arrive on the closed account and face no ratio test & With a successor
link, the linked unit is flagged at once in 99.8 percent of runs, and in
99.1 percent for a clean successor: those first flags are the closed
account's reports. A period later the two separate (96.9 against 0.1
percent) \\
A switched statement descriptor (First Data) & A store that overwrites
the descriptor shows the account view and, in 53.7 percent of runs,
flags the new descriptor before the switch & Per-payment descriptors
rebuild the monitored unit; 95.7 reports a run reach a unit with no
denominator and are reported separately \\
Reports linked only by card number and amount (First Data) & Counting
every link flags a clean account near the floor in 7.8 to 37.0 percent
of runs as unverified links rise from 10 to 50 percent & Counting only
verified links gives no wrong flag and leaves the account inconclusive
in 40.0 to 100 percent of runs \\
A merchant below the chargeback floor (Allied Wallet) & The chargeback
rule never flags an account or the group: no account exceeds 24
chargebacks in a month & The fraud-report rule by value flags each
account in at least 99.8 percent of runs. It needs the separate channel
and each report's amount \\
Refunds before disputes (Allied Wallet) & A chargeback-only store reads
the refunded fifth of reported payments as unreported & The fraud report
is kept as its own record \\
Traffic moved after closures (Allied Wallet) & On average 13.5 fraud
reports a run arrive on closed accounts in months with no sales, which
the account view cannot test & The payout-linked group counts them
against the receiving account's sales \\
\end{longtable}

\endgroup

\textbf{Holds by construction.} The eight accounts sit below the floor
and the alleged fraud ratio is three times the threshold, so those flags
follow from the parameters; only their timing depends on the assumed
delays. The 40 legitimate merchants' report rate is set at 1.6 percent,
above the 1.5 percent threshold, so their flag also follows from the
parameters. The First Data figures divide by the previous period's
captures; with the same period's, the two linked-account figures are
86.0 and 73.7 percent. The refunded share is set to 20 percent.
Displacement follows from the model, which moves each closed account's
volume to the remaining account. Verified links are correct by
assumption.

\textbf{Sweeps (manuscript Figure 3).}
\href{tools/replays/sweeps.py}{\texttt{tools/replays/sweeps.py}} varies
one parameter at a time with the replays' generating processes and
rules, 500 runs a point, and seeds separate from the runs above
(\href{SUPPLEMENT/S10_REPLAYS/sweep_results.json}{results}). First Data
uses the previous period's captures. For eight related accounts at 4
percent, the account view flags in 5.4 percent of runs at 2,000 captures
each a period and 98.8 percent at 2,500; the linked view flags in 27.6
percent at 300, 89.4 percent at 350 and 99.8 percent at 400. At the base
values the sweep's 89.4 percent compares with 91.6 in the 1,000 runs
above. At 350 captures each, the linked view flags the eight accounts in
2.2 percent of runs at a 3 percent report rate, 33.2 percent at 3.5 and
89.4 percent at 4. It flags the 40 legitimate merchants, at 300 each, in
0.6 percent at 1.25, 35.0 percent at 1.5 and 98.4 percent at 1.75. For
one Allied Wallet account at 3.14 percent, the chargeback rule flags in
2.6 percent of runs at 3,000 payments a month and 96.2 percent at 4,000.
The fraud-report rule flags in at least 98.8 percent at every volume
from 25 to 8,000. At an assumed legitimate 0.3 percent, the fraud-report
rule flags in 38.8 to 54.4 percent of runs from 25 to 185 payments a
month, 17.2 percent at 500, 2.8 percent at 1,000 and none from 2,000.
The swept account is steady for eight months; the replay's three linked
accounts are not. A First Data group whose report rate equals the
threshold flags only by chance: reports arrive over up to four periods
while the accounts are active for four, so a period's expected monitored
ratio is 55 to 97 percent of the report rate.

\textbf{Limits of the records.} Padding sales with self-funded purchases
is arithmetic, not simulated: with the assumed fees, holding the
chargeback ratio under 1 percent with \$1 purchases costs 0.2 percent of
sales, and padded payments look like clean payments with no reports.
Only payment-ledger data, such as the funding instrument, can expose
them. Rewriting 12,000 attributions in place passes every rule, so the
store must prevent edits. Accounts with no recorded shared attribute,
shell companies and dummy websites leave nothing for an outcome record
to test.

\textbf{Conformance.} The generated records, 144,574 for First Data and
4,240 for Allied Wallet, pass the schema and rules. A refund recorded as
also closing the fraud-report channel fails \mbox{OR-08}, and a
shared-attribute link marked verified fails \mbox{OR-14}.

\FloatBarrier

\hypertarget{references}{%
\subsection{References}\label{references}}

\suppressfloats[t]

\hypertarget{refs}{}
\begin{CSLReferences}{1}{1}
\leavevmode\vadjust pre{\hypertarget{ref-MITRE2026F3}{}}%
{[}1{]} MITRE Center for Threat-Informed Defense. 2026. Fight Fraud
Framework (F3), version 1.2. GitHub,
center-for-threat-informed-defense/fight-fraud-framework. Cited version:
commit e03c3dc (5 November 2026); version 1.2 released 29 September
2026.
\url{https://github.com/center-for-threat-informed-defense/fight-fraud-framework/tree/e03c3dc0fd765e782f35a6293cc8bb340aa14842}.
Accessed 9 November 2026.

\leavevmode\vadjust pre{\hypertarget{ref-Meli2019}{}}%
{[}2{]} M. Meli, M. R. McNiece, and B. Reaves. 2019. How Bad Can It Git?
Characterizing Secret Leakage in Public GitHub Repositories.
\emph{Network and Distributed Systems Security (NDSS) Symposium 2019}.
\url{https://doi.org/10.14722/ndss.2019.23418}.

\leavevmode\vadjust pre{\hypertarget{ref-Demir2026}{}}%
{[}3{]} N. Demir, Y. Vekaria, G. Smaragdakis, and Z. Durumeric. 2026.
Keys on Doormats: Exposed API Credentials on the Web. arXiv preprint
2603.12498v3. Accepted at ACM CCS 2026 (November 2026) per the reference
printed in arXiv v3. \url{https://arxiv.org/abs/2603.12498v3}.

\leavevmode\vadjust pre{\hypertarget{ref-Gamazo2023}{}}%
{[}4{]} W. Gamazo and N. Quist. 2023. CloudKeys in the Air: Tracking
Malicious Operations of Exposed IAM Keys. Palo Alto Networks Unit 42.
\url{https://unit42.paloaltonetworks.com/malicious-operations-of-exposed-iam-keys-cryptojacking/}.
Accessed 19 October 2026.

\leavevmode\vadjust pre{\hypertarget{ref-Microsoft2022AiTM}{}}%
{[}5{]} Microsoft Threat Intelligence. 2022. From cookie theft to BEC:
Attackers use AiTM phishing sites as entry point to further financial
fraud. Microsoft Security Blog.
\url{https://www.microsoft.com/en-us/security/blog/2022/07/12/from-cookie-theft-to-bec-attackers-use-aitm-phishing-sites-as-entry-point-to-further-financial-fraud/}.
Accessed 19 October 2026.

\leavevmode\vadjust pre{\hypertarget{ref-Doerfler2019}{}}%
{[}6{]} P. Doerfler, M. Marincenko, J. Ranieri, Y. Jiang, A. Moscicki,
D. McCoy, and K. Thomas. 2019. Evaluating Login Challenges as a Defense
Against Account Takeover. \emph{Proceedings of the 2019 World Wide Web
Conference}, pp. 372--382.
\url{https://doi.org/10.1145/3308558.3313481}.

\leavevmode\vadjust pre{\hypertarget{ref-Microsoft2023OAuth}{}}%
{[}7{]} Microsoft Threat Intelligence. 2023. Threat actors misuse OAuth
applications to automate financially driven attacks. Microsoft Security
Blog.
\url{https://www.microsoft.com/en-us/security/blog/2023/12/12/threat-actors-misuse-oauth-applications-to-automate-financially-driven-attacks/}.
Accessed 19 October 2026.

\leavevmode\vadjust pre{\hypertarget{ref-Campobasso2020}{}}%
{[}8{]} M. Campobasso and L. Allodi. 2020. Impersonation-as-a-Service:
Characterizing the Emerging Criminal Infrastructure for User
Impersonation at Scale. \emph{Proceedings of the 2020 ACM SIGSAC
Conference on Computer and Communications Security}, pp. 1665--1680.
\url{https://doi.org/10.1145/3372297.3417892}.

\leavevmode\vadjust pre{\hypertarget{ref-FBI2018}{}}%
{[}9{]} Federal Bureau of Investigation, Internet Crime Complaint
Center. 2018. Cybercriminals Utilize Social Engineering Techniques to
Obtain Employee Credentials to Conduct Payroll Diversion. FBI IC3.
\url{https://www.ic3.gov/PSA/2018/PSA180918}. Accessed 19 October 2026.

\leavevmode\vadjust pre{\hypertarget{ref-FinCEN2014}{}}%
{[}10{]} Financial Crimes Enforcement Network. 2014. Advisory
\mbox{FIN-2014}-A005: Update on U.S. Currency Restrictions in Mexico:
Funnel Accounts and TBML. FinCEN.
\url{https://www.fincen.gov/system/files/advisory/FIN-2014-A005.pdf}.
Accessed 19 October 2026.

\leavevmode\vadjust pre{\hypertarget{ref-Bitaab2025}{}}%
{[}11{]} M. Bitaab, A. Karimi, Z. Lyu, A. Oest, D. Kuchhal, M. Saad,
G.-J. Ahn, R. Wang, T. Bao, Y. Shoshitaishvili, and A. Doupé. 2025.
ScamMagnifier: Piercing the Veil of Fraudulent Shopping Website
Campaigns. \emph{Network and Distributed System Security (NDSS)
Symposium 2025}. \url{https://doi.org/10.14722/ndss.2025.240763}.

\leavevmode\vadjust pre{\hypertarget{ref-McCoy2012}{}}%
{[}12{]} D. McCoy, H. Dharmdasani, C. Kreibich, G. M. Voelker, and S.
Savage. 2012. Priceless: The Role of Payments in Abuse-advertised Goods.
\emph{Proceedings of the 2012 ACM Conference on Computer and
Communications Security}, pp. 845--856.
\url{https://doi.org/10.1145/2382196.2382285}.

\leavevmode\vadjust pre{\hypertarget{ref-Tian2018}{}}%
{[}13{]} H. Tian, S. M. Gaffigan, D. S. West, and D. McCoy. 2018.
Bullet-Proof Payment Processors. \emph{2018 APWG Symposium on Electronic
Crime Research}, pp. 1--11.
\url{https://doi.org/10.1109/eCrime.2018.8376208}.

\leavevmode\vadjust pre{\hypertarget{ref-FTC2020FirstData}{}}%
{[}14{]} Federal Trade Commission. 2020. FTC v. First Data Merchant
Services LLC, No. 1:20-cv-03867 (S.D.N.Y.), Complaint, ECF 1. US
District Court, Southern District of New York. Filed complaint; related
stipulated order and nonadmission terms identified in S1.
\url{https://www.ftc.gov/system/files/documents/cases/first_data_filed_complaint.pdf}.
Accessed 22 October 2026.

\leavevmode\vadjust pre{\hypertarget{ref-FTC2025Paddle}{}}%
{[}15{]} Federal Trade Commission. 2025. FTC v. Paddle.com Market
Limited, No. 1:25-cv-01886 (D.D.C.), Complaint, ECF 1. US District
Court, District of Columbia. Filed complaint; related stipulated order
and nonadmission terms identified in S1.
\url{https://www.ftc.gov/system/files/ftc_gov/pdf/PaddleComplaintForPermanentInjunction,MonetaryJudgment,andOtherRelief.pdf}.
Accessed 21 October 2026.

\leavevmode\vadjust pre{\hypertarget{ref-FTC2019Allied}{}}%
{[}16{]} Federal Trade Commission. 2019. FTC v. AlliedWallet, Inc., No.
2:19-cv-04355 (C.D. Cal.), Complaint, ECF 24. US District Court, Central
District of California. Filed complaint; related stipulated order and
nonadmission terms identified in S1.
\url{https://www.ftc.gov/system/files/documents/cases/alliedwallet_complaint_7-3-19.pdf}.
Accessed 21 October 2026.

\leavevmode\vadjust pre{\hypertarget{ref-Stripe2026CAE}{}}%
{[}17{]} Stripe. 2026. Customer account evaluations. Stripe
documentation.
\url{https://docs.stripe.com/radar/customer-account-evaluations}.
Accessed 19 October 2026.

\leavevmode\vadjust pre{\hypertarget{ref-Stripe2026Trial}{}}%
{[}18{]} Stripe. 2026. Free trial abuse prevention. Stripe
documentation. \url{https://docs.stripe.com/radar/free-trial-abuse}.
Accessed 19 October 2026.

\leavevmode\vadjust pre{\hypertarget{ref-Stripe2026PAYG}{}}%
{[}19{]} Stripe. 2026. Pay-as-you-go abuse evaluation. Stripe
documentation. \url{https://docs.stripe.com/radar/pay-as-you-go-abuse}.
Accessed 19 October 2026.

\leavevmode\vadjust pre{\hypertarget{ref-Stripe2026Farming}{}}%
{[}20{]} Stripe. 2026. Account farming prevention. Stripe documentation.
Preview feature; Radar Pro plan required.
\url{https://docs.stripe.com/radar/account-farming}. Accessed 19 October
2026.

\leavevmode\vadjust pre{\hypertarget{ref-Stripe2026MonitoringPrograms}{}}%
{[}21{]} Stripe. 2026. Dispute and fraud card monitoring programs.
Stripe documentation.
\url{https://docs.stripe.com/disputes/monitoring-programs}. Accessed 25
October 2026.

\leavevmode\vadjust pre{\hypertarget{ref-Stripe2026Disputes}{}}%
{[}22{]} Stripe. 2026. How disputes work. Stripe documentation.
\url{https://docs.stripe.com/disputes/how-disputes-work}. Accessed 25
October 2026.

\leavevmode\vadjust pre{\hypertarget{ref-Stripe2026SEPADebit}{}}%
{[}23{]} Stripe. 2026. SEPA Direct Debit payments. Stripe documentation.
\url{https://docs.stripe.com/payments/sepa-debit}. Accessed 25 October
2026.

\leavevmode\vadjust pre{\hypertarget{ref-Stripe2026ACHTerms}{}}%
{[}24{]} Stripe. 2026. ACH Payment Terms. Stripe Legal.
\url{https://stripe.com/legal/ACH}. Accessed 25 October 2026.

\leavevmode\vadjust pre{\hypertarget{ref-LeadBank2026ACHReturns}{}}%
{[}25{]} Lead Bank. 2026. Returns, Reversals, \& Exceptions. Lead Bank
documentation.
\url{https://docs.lead.bank/api-integrations/ach/return-codes}. Accessed
25 October 2026.

\leavevmode\vadjust pre{\hypertarget{ref-StripeOpenAPI2026}{}}%
{[}26{]} Stripe. 2026. OpenAPI specification for the Stripe API. GitHub,
stripe/openapi.
\url{https://github.com/stripe/openapi/tree/18fa2cc768024b47789dc6fa243acd48478acac2}.
Accessed 22 October 2026.

\leavevmode\vadjust pre{\hypertarget{ref-Leach2019}{}}%
{[}27{]} B. Leach. 2019. Fast and flexible observability with canonical
log lines. Stripe blog.
\url{https://stripe.com/blog/canonical-log-lines}. Accessed 21 October
2026.

\leavevmode\vadjust pre{\hypertarget{ref-Stripe2026CryptoOnramp}{}}%
{[}28{]} Stripe. 2026. Onramp API reference. Stripe documentation.
Public preview. \url{https://docs.stripe.com/crypto/using-the-api}.
Accessed 19 October 2026.

\leavevmode\vadjust pre{\hypertarget{ref-Stripe2026Connect}{}}%
{[}29{]} Stripe. 2026. Best practices for risk management. Stripe
documentation.
\url{https://docs.stripe.com/connect/risk-management/best-practices}.
Accessed 19 October 2026.

\leavevmode\vadjust pre{\hypertarget{ref-Chapelle2014Delayed}{}}%
{[}30{]} O. Chapelle. 2014. Modeling Delayed Feedback in Display
Advertising. \emph{Proceedings of the 20th ACM SIGKDD International
Conference on Knowledge Discovery and Data Mining}, pp. 1097--1105.
\url{https://doi.org/10.1145/2623330.2623634}.

\leavevmode\vadjust pre{\hypertarget{ref-Stripe2026ConnectCharges}{}}%
{[}31{]} Stripe. 2026. Understand how charges work in a Connect
integration. Stripe documentation.
\url{https://docs.stripe.com/connect/charges}. Accessed 19 October 2026.

\leavevmode\vadjust pre{\hypertarget{ref-Visa2025VAMP}{}}%
{[}32{]} Visa. 2025. Visa Acquirer Monitoring Program Overview. Visa.
\url{https://corporate.visa.com/content/dam/VCOM/corporate/visa-perspectives/security-and-trust/documents/visa-acquirer-monitoring-program-fact-sheet-2025.pdf}.
Accessed 25 October 2026.

\leavevmode\vadjust pre{\hypertarget{ref-Dhama2026Causal}{}}%
{[}33{]} G. Dhama. 2026. Causal Label Recovery in Payment Networks.
arXiv preprint 2605.29272v1. \url{https://arxiv.org/abs/2605.29272v1}.

\leavevmode\vadjust pre{\hypertarget{ref-HernanRobins2020}{}}%
{[}34{]} M. A. Hernán and J. M. Robins. 2020. Causal Inference: What If.
Boca Raton, Chapman \& Hall/CRC. Chapters cited: 1, 2, 3, 7 and 19.
Cited version: author-hosted revision (19 August 2026).
\url{https://miguelhernan.org/whatifbook}.

\leavevmode\vadjust pre{\hypertarget{ref-Lakkaraju2017Selective}{}}%
{[}35{]} H. Lakkaraju, J. Kleinberg, J. Leskovec, J. Ludwig, and S.
Mullainathan. 2017. The Selective Labels Problem: Evaluating Algorithmic
Predictions in the Presence of Unobservables. \emph{Proceedings of the
23rd ACM SIGKDD International Conference on Knowledge Discovery and Data
Mining}, pp. 275--284. \url{https://doi.org/10.1145/3097983.3098066}.

\leavevmode\vadjust pre{\hypertarget{ref-Liu2018GEM}{}}%
{[}36{]} Z. Liu, C. Chen, X. Yang, J. Zhou, X. Li, and L. Song. 2018.
Heterogeneous Graph Neural Networks for Malicious Account Detection.
\emph{Proceedings of the 27th ACM International Conference on
Information and Knowledge Management}.
\url{https://doi.org/10.1145/3269206.3272010}.

\leavevmode\vadjust pre{\hypertarget{ref-Fatima2024}{}}%
{[}37{]} K. Fatima, M. Schuckers, G. Cruz-Ortiz, D. Hou, S. Purnapatra,
T. Andrews, A. Neupane, B. Marshall, and S. Schuckers. 2024. A
Large-Scale Study of Performance and Equity of Commercial Remote
Identity Verification Technologies Across Demographics. \emph{2024 IEEE
International Joint Conference on Biometrics}, pp. 1--8.
\url{https://doi.org/10.1109/IJCB62174.2024.10744432}.

\leavevmode\vadjust pre{\hypertarget{ref-Sharma2026AuditableFraud}{}}%
{[}38{]} R. Sharma. 2026. Toward Auditable Fraud Detection: Combining
Graph Features, Model Explanations, and Agentic Case Investigation.
arXiv preprint 2607.19266v1. \url{https://arxiv.org/abs/2607.19266v1}.

\leavevmode\vadjust pre{\hypertarget{ref-Nasr2025}{}}%
{[}39{]} M. Nasr, N. Carlini, C. Sitawarin, S. V. Schulhoff, J. Hayes,
M. Ilie, J. Pluto, S. Song, H. Chaudhari, I. Shumailov, A. Thakurta, K.
Y. Xiao, A. Terzis, and F. Tramèr. 2025. The Attacker Moves Second:
Stronger Adaptive Attacks Bypass Defenses Against LLM Jailbreaks and
Prompt Injections. arXiv preprint 2510.09023v1.
\url{https://arxiv.org/abs/2510.09023v1}.

\leavevmode\vadjust pre{\hypertarget{ref-Debenedetti2025CaMeL}{}}%
{[}40{]} E. Debenedetti, I. Shumailov, T. Fan, J. Hayes, N. Carlini, D.
Fabian, C. Kern, C. Shi, A. Terzis, and F. Tramèr. 2025. Defeating
Prompt Injections by Design. arXiv preprint 2503.18813v2.
\url{https://arxiv.org/abs/2503.18813v2}.

\leavevmode\vadjust pre{\hypertarget{ref-Dash2026}{}}%
{[}41{]} P. Dash, T. Ge, A. Jain, T. Shah, and Z. Shang. 2026. From
Untrusted Input to Trusted Memory: A Systematic Study of Memory
Poisoning Attacks in LLM Agents. arXiv preprint 2606.04329v2.
\url{https://arxiv.org/abs/2606.04329v2}.

\leavevmode\vadjust pre{\hypertarget{ref-Xiong2026MemoryManagement}{}}%
{[}42{]} Z. Xiong, Y. Lin, W. Xie, P. He, Z. Liu, J. Tang, H. Lakkaraju,
and Z. Xiang. 2026. How Memory Management Impacts LLM Agents: An
Empirical Study of Experience-Following Behavior. \emph{Proceedings of
the 64th Annual Meeting of the Association for Computational Linguistics
(Volume 1: Long Papers)}, pp. 623--645.
\url{https://doi.org/10.18653/v1/2026.acl-long.27}.

\leavevmode\vadjust pre{\hypertarget{ref-Wolfe2005Rare}{}}%
{[}43{]} J. M. Wolfe, T. S. Horowitz, and N. M. Kenner. 2005. Rare Items
Often Missed in Visual Searches. \emph{Nature}, 435, pp. 439--440.
\url{https://doi.org/10.1038/435439a}.

\leavevmode\vadjust pre{\hypertarget{ref-Bottou2013Counterfactual}{}}%
{[}44{]} L. Bottou, J. Peters, J. Quiñonero-Candela, D. X. Charles, D.
M. Chickering, E. Portugaly, D. Ray, P. Simard, and E. Snelson. 2013.
Counterfactual Reasoning and Learning Systems: The Example of
Computational Advertising. \emph{Journal of Machine Learning Research},
14(101), pp. 3207--3260.
\url{https://jmlr.org/papers/v14/bottou13a.html}.

\leavevmode\vadjust pre{\hypertarget{ref-Singh2019DecisionProvenance}{}}%
{[}45{]} J. Singh, J. Cobbe, and C. Norval. 2019. Decision Provenance:
Harnessing Data Flow for Accountable Systems. \emph{IEEE Access}, 7, pp.
6562--6574. \url{https://doi.org/10.1109/ACCESS.2018.2887201}.

\leavevmode\vadjust pre{\hypertarget{ref-Ehrhardt2021RejectInference}{}}%
{[}46{]} A. Ehrhardt, C. Biernacki, V. Vandewalle, P. Heinrich, and S.
Beben. 2021. Reject Inference Methods in Credit Scoring. \emph{Journal
of Applied Statistics}, 48(13-15), pp. 2734--2754.
\url{https://doi.org/10.1080/02664763.2021.1929090}.

\leavevmode\vadjust pre{\hypertarget{ref-Dhama2026Limits}{}}%
{[}47{]} G. Dhama. 2026. The Fundamental Limits of Fraud Detection in
Card Payment Networks: An Information-Theoretic Theory of Delayed,
Censored, Corrupted, and Counterfactual Feedback. arXiv preprint
2605.27557v2. \url{https://arxiv.org/abs/2605.27557v2}.

\leavevmode\vadjust pre{\hypertarget{ref-Egg2025AdyenOPE}{}}%
{[}48{]} A. Egg. 2025. Off-policy Evaluation for Payments at Adyen.
arXiv preprint 2501.10470v1. \url{https://arxiv.org/abs/2501.10470v1}.

\leavevmode\vadjust pre{\hypertarget{ref-Tocker2022VirtualControl}{}}%
{[}49{]} M. Tocker. 2022. Virtual Control Group: Measuring Hidden
Performance Metrics. arXiv preprint 2208.12941v1.
\url{https://arxiv.org/abs/2208.12941v1}.

\leavevmode\vadjust pre{\hypertarget{ref-Deng2026FCAC}{}}%
{[}50{]} J. Deng. 2026. When Can Fraud Operations Authorize Automation?
A Decision-Support Framework for Fresh Audit Evidence and Review
Workload. arXiv preprint 2608.08577v1.
\url{https://arxiv.org/abs/2608.08577v1}.

\leavevmode\vadjust pre{\hypertarget{ref-Solozobov2026DecisionTrace}{}}%
{[}51{]} O. Solozobov. 2026. Decision Trace Schema for Governance
Evidence in Real-Time Risk Systems. arXiv preprint 2604.09296v1.
\url{https://arxiv.org/abs/2604.09296v1}.

\leavevmode\vadjust pre{\hypertarget{ref-Wu2026EvidenceCoverage}{}}%
{[}52{]} M. Wu, L. Wang, Y. Zhang, and J. Deng. 2026. Evidence Coverage
for Intent-Bound Execution: Scope, Obligations, and Cutoff Reasoning.
arXiv preprint 2610.01386v1. \url{https://arxiv.org/abs/2610.01386v1}.

\leavevmode\vadjust pre{\hypertarget{ref-Manapat2015Counterfactual}{}}%
{[}53{]} M. Manapat. 2015. Counterfactual evaluation of machine learning
models. Conference talk, PyData Seattle.
\url{https://pydata.org/seattle2015/schedule/presentation/17/}.

\leavevmode\vadjust pre{\hypertarget{ref-Snodgrass1986Temporal}{}}%
{[}54{]} R. Snodgrass and I. Ahn. 1986. Temporal Databases.
\emph{Computer}, 19(9), pp. 35--42.
\url{https://doi.org/10.1109/MC.1986.1663327}.

\leavevmode\vadjust pre{\hypertarget{ref-Razniewski2011Completeness}{}}%
{[}55{]} S. Razniewski and W. Nutt. 2011. Completeness of Queries over
Incomplete Databases. \emph{Proceedings of the VLDB Endowment}, 4(11),
pp. 749--760. \url{https://doi.org/10.14778/3402707.3402715}.

\leavevmode\vadjust pre{\hypertarget{ref-Akidau2015Dataflow}{}}%
{[}56{]} T. Akidau, R. Bradshaw, C. Chambers, S. Chernyak, R. J.
Fernández-Moctezuma, R. Lax, S. McVeety, D. Mills, F. Perry, E. Schmidt,
and S. Whittle. 2015. The Dataflow Model: A Practical Approach to
Balancing Correctness, Latency, and Cost in Massive-Scale, Unbounded,
Out-of-Order Data Processing. \emph{Proceedings of the VLDB Endowment},
8(12), pp. 1792--1803. \url{https://doi.org/10.14778/2824032.2824076}.

\leavevmode\vadjust pre{\hypertarget{ref-Wang2026TOKI}{}}%
{[}57{]} Z. Wang. 2026. TOKI: A Bitemporal Operator Algebra for
Contradiction Resolution in LLM-Agent Persistent Memory. arXiv preprint
2606.06240v1. \url{https://arxiv.org/abs/2606.06240v1}.

\leavevmode\vadjust pre{\hypertarget{ref-Rasmussen2025Zep}{}}%
{[}58{]} P. Rasmussen, P. Paliychuk, T. Beauvais, J. Ryan, and D.
Chalef. 2025. Zep: A Temporal Knowledge Graph Architecture for Agent
Memory. arXiv preprint 2501.13956v1.
\url{https://arxiv.org/abs/2501.13956v1}.

\leavevmode\vadjust pre{\hypertarget{ref-Dan2026Anamnesis}{}}%
{[}59{]} A. Dan. 2026. Anamnesis: An Authority-Preserving Runtime
Contract for Persistent AI Memory. Zenodo preprint, Version 1.0, 16
August 2026. \url{https://zenodo.org/records/22118090}.

\leavevmode\vadjust pre{\hypertarget{ref-Brown2026Quipu}{}}%
{[}60{]} S. Brown. 2026. Quipu: A Governed Bitemporal Knowledge Graph
Store. arXiv preprint 2608.16813v1.
\url{https://arxiv.org/abs/2608.16813v1}.

\leavevmode\vadjust pre{\hypertarget{ref-Dai2026Hearsay}{}}%
{[}61{]} J. Dai, Z. Yang, P. Shao, K. Ng, Z. Liu, C. Ju, Y. He, and B.
Hu. 2026. Hearsay: Can an Auditor Trust the Record a Deployed Agent
Harness Writes? arXiv preprint 2609.32495v1.
\url{https://arxiv.org/abs/2609.32495v1}.

\leavevmode\vadjust pre{\hypertarget{ref-Doyle1979TMS}{}}%
{[}62{]} J. Doyle. 1979. A Truth Maintenance System. \emph{Artificial
Intelligence}, 12(3), pp. 231--272.
\url{https://doi.org/10.1016/0004-3702(79)90008-0}.

\leavevmode\vadjust pre{\hypertarget{ref-Malheiro1994BeliefRevision}{}}%
{[}63{]} B. Malheiro, N. R. Jennings, and E. Oliveira. 1994. Belief
Revision in Multi-Agent Systems. \emph{Proceedings of the 11th European
Conference on Artificial Intelligence}, pp. 294--298.
\url{https://eprints.soton.ac.uk/252143}.

\leavevmode\vadjust pre{\hypertarget{ref-OASIS2021STIX}{}}%
{[}64{]} OASIS. 2021. STIX Version 2.1. OASIS Standard.
\url{https://docs.oasis-open.org/cti/stix/v2.1/os/stix-v2.1-os.html}.

\leavevmode\vadjust pre{\hypertarget{ref-ODNI2015ICD203}{}}%
{[}65{]} Office of the Director of National Intelligence. 2015.
Intelligence Community Directive 203: Analytic Standards. Directive. As
amended by technical amendments through 21 January 2022.
\url{https://archive.dni.gov/files/documents/ICD/ICD-203.pdf}.

\leavevmode\vadjust pre{\hypertarget{ref-Qi2026GovMem}{}}%
{[}66{]} Y. Qi, X. Xu, and Y. Li. 2026. When Not to Write Memory:
Governing False Promotion from Correlated Agent Traces. arXiv preprint
2607.02579v1. \url{https://arxiv.org/abs/2607.02579v1}.

\leavevmode\vadjust pre{\hypertarget{ref-Akhtar2026Contestable}{}}%
{[}67{]} H. Akhtar and B. Freeman. 2026. Contestable Agents for Data
Governance: Recording Human Stewardship of AI-Generated Metadata Across
Structured and Unstructured Systems. arXiv preprint 2610.31553v1.
\url{https://arxiv.org/abs/2610.31553v1}.

\leavevmode\vadjust pre{\hypertarget{ref-AlbarrakJagtap2026Industry4CTI}{}}%
{[}68{]} M. Albarrak and S. Jagtap. 2026. Designing an AI-Assisted Cyber
Threat Intelligence Framework for Industry 4.0: A Human-in-the-Loop
Design Science Approach. \emph{Applied Sciences}, 16(15), Article 7646.
\url{https://doi.org/10.3390/app16157646}.

\leavevmode\vadjust pre{\hypertarget{ref-Wang2026MAPGraph}{}}%
{[}69{]} Y. Wang, Z. Yan, J. Zhang, Z. Wu, M. Zheng, Z. Sun, Y. Zhu, and
T. Cai. 2026. MAP-Graph: Provenance-Aware Shared Memory for Multi-Agent
Workflows. arXiv preprint 2608.10509v1.
\url{https://arxiv.org/abs/2608.10509v1}.

\leavevmode\vadjust pre{\hypertarget{ref-Li2026EvidenceToEffect}{}}%
{[}70{]} Y. Li, Z. Xu, S. Volkov, H. Liu, T. Zhou, X. Chen, D. Wan, D.
Shao, Y. Luo, and H. Sun. 2026. From Evidence to Effect: Authority
Semantics and Runtime Infrastructure for Stateful Agents. arXiv preprint
2609.08472v3. Cited version: arXiv 2609.08472v3 (28 September 2026).
\url{https://arxiv.org/abs/2609.08472v3}.

\leavevmode\vadjust pre{\hypertarget{ref-Bara2026EpistemicSybil}{}}%
{[}71{]} M. Bara. 2026. Epistemic Sybil Resistance: Multiplying AI
Agents Without Multiplying Evidence. arXiv preprint 2609.01873v1.
\url{https://arxiv.org/abs/2609.01873v1}.

\leavevmode\vadjust pre{\hypertarget{ref-Evron2026AdversaryBrains}{}}%
{[}72{]} G. Evron. 2026. Adversary Brains: Hypothesis-Based Logical
Signatures for Threat Actor Detection with LLMs. arXiv preprint
2610.35624v1. \url{https://arxiv.org/abs/2610.35624v1}.

\leavevmode\vadjust pre{\hypertarget{ref-Koch2026AgileV}{}}%
{[}73{]} C. Koch. 2026. From Agent Output to Authorized Transition.
arXiv preprint 2609.28216v1. The v1 PDF title page reads ``From Agent
Output to Authorized Transition: The Agile-V Assurance Spine for
State-Bound Evidence in Software, Firmware, and Hardware Engineering''
instead. \url{https://arxiv.org/abs/2609.28216v1}.

\leavevmode\vadjust pre{\hypertarget{ref-Mitra2026AgenticCyOps}{}}%
{[}74{]} S. Mitra, R. Patel, S. Mittal, M. R. Rahman, and S. Rahimi.
2026. AgenticCyOps: Securing Multi-Agentic AI Integration in Enterprise
Cyber Operations. arXiv preprint 2603.09134v1.
\url{https://arxiv.org/abs/2603.09134v1}.

\leavevmode\vadjust pre{\hypertarget{ref-Tausz2020}{}}%
{[}75{]} A. Tausz. 2020. Similarity clustering to catch fraud rings.
Stripe blog. \url{https://stripe.com/blog/similarity-clustering}.
Accessed 21 October 2026.

\leavevmode\vadjust pre{\hypertarget{ref-Drapeau2024}{}}%
{[}76{]} R. Drapeau and I. Khosla. 2024. The ML flywheel: how we
continually improve our models to reduce card testing. Stripe blog.
\url{https://stripe.com/blog/the-ml-flywheel-how-we-continually-improve-our-models-to-reduce-card-testing}.
Accessed 21 October 2026.

\leavevmode\vadjust pre{\hypertarget{ref-Erdbrink2023TreatmentPatent}{}}%
{[}77{]} C. Erdbrink and L. Ellis. 2023. Analyzing a Transaction in a
Payment Processing System. US patent application publication,
US20230298028A1, Fidelity Information Services, LLC.
\url{https://patents.google.com/patent/US20230298028A1/en}.

\leavevmode\vadjust pre{\hypertarget{ref-Sakellariou2022CTIReferenceModel}{}}%
{[}78{]} G. Sakellariou, P. Fouliras, I. Mavridis, and P. Sarigiannidis.
2022. A Reference Model for Cyber Threat Intelligence (CTI) Systems.
\emph{Electronics}, 11(9), Article 1401.
\url{https://doi.org/10.3390/electronics11091401}.

\leavevmode\vadjust pre{\hypertarget{ref-Dimitriadis2025}{}}%
{[}79{]} A. Dimitriadis, A. Papoutsis, D. Kavalieros, T. Tsikrika, S.
Vrochidis, and I. Kompatsiaris. 2025. EVACTI: Evaluating the
Actionability of Cyber Threat Intelligence. \emph{International Journal
of Information Security}, 24, Article 123.
\url{https://doi.org/10.1007/s10207-025-01033-z}.

\leavevmode\vadjust pre{\hypertarget{ref-Wu2026ExCyTIn}{}}%
{[}80{]} Y. Wu, M. Velazco, A. Zhao, M. R. Meléndez Luján, S. Movva, Y.
K. Roy, Q. Nguyen, R. Rodriguez, Q. Wu, M. Albada, J. Kiseleva, and A.
Mudgerikar. 2026. ExCyTIn-Bench: Evaluating LLM agents on Cyber Threat
Investigation. arXiv preprint 2507.14201v3. Accepted at ICML 2026; final
proceedings pagination not established. Cited version: arXiv
2507.14201v3 (1 May 2026). \url{https://arxiv.org/abs/2507.14201v3}.

\leavevmode\vadjust pre{\hypertarget{ref-Tan2026SRFraud}{}}%
{[}81{]} X. Tan, Y. Ma, and X. Zhang. 2026. SR-Fraud: An
Outcome-Supervised Reflective LLM Agent Framework for Non-Stationary
Payment Fraud Detection. arXiv preprint 2609.27287v1.
\url{https://arxiv.org/abs/2609.27287v1}.

\leavevmode\vadjust pre{\hypertarget{ref-Mukherjee2025PROVSEEK}{}}%
{[}82{]} K. Mukherjee and M. Kantarcioglu. 2025. LLM-driven Provenance
Forensics for Threat Investigation and Detection. arXiv preprint
2508.21323v2. The v2 PDF title page reads ``LLM-driven Provenance
Forensics for Threat Intelligence and Detection'' instead.
\url{https://arxiv.org/abs/2508.21323v2}.

\leavevmode\vadjust pre{\hypertarget{ref-Alten2026DecisionRecords}{}}%
{[}83{]} A. Alten. 2026. When Does Persistent Agent Memory Become
Accountable? Scalytics blog.
\url{https://www.scalytics.io/blog/agent-memory-governance-decision-records}.
Accessed 26 October 2026.

\leavevmode\vadjust pre{\hypertarget{ref-Talamadupula2026AgentMemory}{}}%
{[}84{]} K. Talamadupula. 2026. Oracle AI Agent Memory Brings Persistent
Enterprise Context to NVIDIA OpenShell. Oracle Database blog.
\url{https://blogs.oracle.com/database/oracle-ai-agent-memory-brings-persistent-enterprise-context-to-nvidia-openshell}.
Accessed 26 October 2026.

\leavevmode\vadjust pre{\hypertarget{ref-Ding2026AlwaysOn}{}}%
{[}85{]} T. Ding, A. Nannapaneni, B. Liu, and L. Zhang. 2026. Always-On
Agents: A Survey of Persistent Memory, State, and Governance in LLM
Agents. arXiv preprint 2606.30306v1.
\url{https://arxiv.org/abs/2606.30306v1}.

\leavevmode\vadjust pre{\hypertarget{ref-Xia2024EDDOps}{}}%
{[}86{]} B. Xia, Q. Lu, L. Zhu, Z. Xing, D. Zhao, and H. Zhang. 2025.
Evaluation-Driven Development and Operations of LLM Agents: A Process
Model and Reference Architecture. arXiv preprint 2411.13768v3.
\url{https://arxiv.org/abs/2411.13768v3}.

\leavevmode\vadjust pre{\hypertarget{ref-MindForge2026Operationalisation}{}}%
{[}87{]} MindForge Consortium. 2026. AI Risk Management:
Operationalisation Handbook. Handbook, Monetary Authority of Singapore.
\url{https://www.mas.gov.sg/-/media/mas-media-library/schemes-and-initiatives/ftig/project-mindforge/mindforge-ai-risk-management-operationalisation-handbook.pdf}.

\leavevmode\vadjust pre{\hypertarget{ref-MAS2026SAFR}{}}%
{[}88{]} Monetary Authority of Singapore. 2026. Safeguards for Agentic
Finance at Runtime (SAFR). White paper, version 1.1, Monetary Authority
of Singapore.
\url{https://www.mas.gov.sg/-/media/mas-media-library/development/fintech/ai-safr/safr-v1-1.pdf}.

\leavevmode\vadjust pre{\hypertarget{ref-CERTIn2026AIAudit}{}}%
{[}89{]} Indian Computer Emergency Response Team. 2026. AI Security
Auditing Guidelines. Guidelines, version 1.0, CERT-In.
\url{https://www.cert-in.org.in/PDF/AI_Security_Auditing_Guidelines.pdf}.

\leavevmode\vadjust pre{\hypertarget{ref-Zhao2024ExpeL}{}}%
{[}90{]} A. Zhao, D. Huang, Q. Xu, M. Lin, Y.-J. Liu, and G. Huang.
2024. ExpeL: LLM Agents Are Experiential Learners. \emph{Proceedings of
the AAAI Conference on Artificial Intelligence}, 38(17), pp.
19632--19642. \url{https://doi.org/10.1609/aaai.v38i17.29936}.

\leavevmode\vadjust pre{\hypertarget{ref-Wei2026AMemGuard}{}}%
{[}91{]} Q. Wei, T. Yang, Y. Wang, X. Li, L. Li, Z. Yin, Y. Zhan, T.
Holz, Z. Lin, and X. Wang. 2025. A-MemGuard: A Proactive Defense
Framework for LLM-Based Agent Memory. arXiv preprint 2510.02373v1. ICML
2026 acceptance reported; final proceedings version not established.
Cited version: arXiv 2510.02373v1 (29 September 2025).
\url{https://arxiv.org/abs/2510.02373v1}.

\leavevmode\vadjust pre{\hypertarget{ref-Zhan2026AuthMemBench}{}}%
{[}92{]} Q. Zhan, R. Zhang, S. Guo, L. Zhao, and Z. Liu. 2026. When
Memory Becomes Authority: Benchmarking Authority Collapse at the Memory
Consolidation Boundary. arXiv preprint 2608.01679v2.
\url{https://arxiv.org/abs/2608.01679v2}.

\leavevmode\vadjust pre{\hypertarget{ref-Xu2026PPMF}{}}%
{[}93{]} J. Xu, Y. Xiao, W. Shao, H. Liu, and X. Li. 2026. Memory
Provenance Laundering in LLM Agents: A Non-Amplification Firewall for
Persistent Memory. arXiv preprint 2607.29167v1.
\url{https://arxiv.org/abs/2607.29167v1}.

\leavevmode\vadjust pre{\hypertarget{ref-Hettema2021Rationality}{}}%
{[}94{]} H. Hettema. 2021. Rationality Constraints in Cyber Defense:
Incident Handling, Attribution and Cyber Threat Intelligence.
\emph{Computers \& Security}, 109, Article 102396.
\url{https://doi.org/10.1016/j.cose.2021.102396}.

\leavevmode\vadjust pre{\hypertarget{ref-Zhu2026TierMem}{}}%
{[}95{]} Q. Zhu, S. Chen, R. Yu, Z. Wu, and B. Wang. 2026. From Lossy to
Verified: A Provenance-Aware Tiered Memory for Agents. arXiv preprint
2602.17913v1. Also presented as a poster at the ICLR 2026 Workshop on
Memory for LLM-Based Agentic Systems (MemAgents). Cited version: arXiv
2602.17913v1 (20 February 2026).
\url{https://arxiv.org/abs/2602.17913v1}.

\leavevmode\vadjust pre{\hypertarget{ref-Lam2026SSGM}{}}%
{[}96{]} C. Lam, J. Li, L. Zhang, and K. Zhao. 2026. Governing Evolving
Memory in LLM Agents: Risks, Mechanisms, and the Stability and Safety
Governed Memory (SSGM) Framework. arXiv preprint 2603.11768v2.
\url{https://arxiv.org/abs/2603.11768v2}.

\leavevmode\vadjust pre{\hypertarget{ref-Margalit2026GovernedMemory}{}}%
{[}97{]} Y. Margalit, N. Cohen-Inger, E. Avram, R. Taig, and O.
Margalit. 2026. Governed Shared Memory for Multi-Agent LLM Systems.
arXiv preprint 2606.24535v1. \url{https://arxiv.org/abs/2606.24535v1}.

\leavevmode\vadjust pre{\hypertarget{ref-Souza2025ProvAgent}{}}%
{[}98{]} R. Souza, A. Gueroudji, S. DeWitt, D. Rosendo, T. Ghosal, R.
Ross, P. Balaprakash, and R. Ferreira da Silva. 2025. PROV-AGENT:
Unified Provenance for Tracking AI Agent Interactions in Agentic
Workflows. \emph{Proceedings of the 21st IEEE International Conference
on e-Science}. Cited version: arXiv 2508.02866v3 (20 August 2025).
\url{https://arxiv.org/abs/2508.02866v3}.

\leavevmode\vadjust pre{\hypertarget{ref-Kaufman2011Leakage}{}}%
{[}99{]} S. Kaufman, S. Rosset, and C. Perlich. 2011. Leakage in Data
Mining: Formulation, Detection, and Avoidance. \emph{Proceedings of the
17th ACM SIGKDD International Conference on Knowledge Discovery and Data
Mining}, pp. 556--563. \url{https://doi.org/10.1145/2020408.2020496}.

\end{CSLReferences}

\end{document}
